\documentclass[10pt, letterpaper]{article}

% Inhaltsverzeichnis für Pakettypen (nur für Übersicht im Header, wird nicht im Dokument angezeigt)
% 1. Seitenlayout und Ränder
% 2. Sprache und Zeichensatz
% 3. Mathematik und Theorem-Umgebungen
% 4. Eigene Makros
% 5. Diagramme und Grafiken
% 6. Tabellen und Aufzählungen
% 7. Inhaltsverzeichnis
% 8. Abschnittsüberschriften
% 9. Abstrakt-Umgebung
% 10. Todos/Notizen
% 11. Rahmen/Box-Umgebungen
% 12. Python-Integration
% 13. Literaturverwaltung
% 14. Hyperlinks
% 15. Absatzeinstellungen
% 16. Umgebungen
% 17  Graphik
% 00. Titel und Autor

% --- 1. Seitenlayout und Ränder ---
\usepackage[margin=3cm]{geometry}

% --- 2. Sprache und Zeichensatz ---
\usepackage[english]{babel}
\usepackage[T1]{fontenc}
\usepackage[utf8]{inputenc}

% --- 3. Mathematik und Theorem-Umgebungen ---
\usepackage{amsmath, amssymb, amsthm}
\usepackage{mathrsfs}
\DeclareMathOperator{\WF}{WF}

% --- 4. Eigene Makros ---
\usepackage{xcolor}
\newcommand{\SKP}{\langle\cdot,\cdot\rangle}
\newcommand{\R}{\mathbb{R}}
\newcommand{\N}{\mathbb{N}}
\newcommand{\Q}{\mathbb{Q}}
\newcommand{\Z}{\mathbb{Z}}
\newcommand{\C}{\mathbb{C}}
\newcommand{\entwurf}[1]{\textcolor{red}{#1}}

% --- 5. Diagramme und Grafiken ---
\usepackage{graphicx}
\usepackage{tikz}
\usetikzlibrary{decorations.pathreplacing, arrows.meta, positioning}
\usepackage{tikz-cd}

% --- 6. Tabellen und Aufzählungen ---
\usepackage{enumitem}
\setlist[itemize]{left=0.5cm}

\newenvironment{romanenum}[1][]
  {%
    \ifx&#1&
    \else
      \textbf{#1}\quad
    \fi
    \begin{enumerate}[label=\roman*)]
  }
  {%
    \end{enumerate}%
  }

% --- 7. Inhaltsverzeichnis ---
\usepackage{tocloft}
\renewcommand{\cftsecfont}{\footnotesize}
\renewcommand{\cftsubsecfont}{\footnotesize}
\renewcommand{\cftsubsubsecfont}{\footnotesize}
\renewcommand{\cftsecpagefont}{\footnotesize}
\renewcommand{\cftsubsecpagefont}{\footnotesize}
\renewcommand{\cftsubsubsecpagefont}{\footnotesize}
\usepackage{etoc}

% --- 8. Abschnittsüberschriften ---
\usepackage{titlesec}
\titleformat{\section}{\normalfont\large\bfseries}{\thesection}{1em}{}
\titleformat{\subsection}{\normalfont\normalsize\bfseries}{\thesubsection}{0.5em}{}
\titleformat{\subsubsection}{\normalfont\normalsize\bfseries}{\thesubsubsection}{0.5em}{}
\setcounter{secnumdepth}{4}

% --- 9. Abstrakt-Umgebung ---
\usepackage{changepage}
\renewenvironment{abstract}
  {
    \begin{adjustwidth}{1.5cm}{1.5cm}
    \small
    \textsc{Abstract. –}%
  }
  {
    \end{adjustwidth}
  }

% --- 10. Todos/Notizen ---
\usepackage{todonotes}

% --- 11. Rahmen/Box-Umgebungen ---
\usepackage{mdframed}
\usepackage{tcolorbox}
\colorlet{shadecolor}{gray!25}

\newenvironment{customTheorem}
  {\vspace{10pt}%
   \begin{mdframed}[
     backgroundcolor=gray!20,
     linewidth=0pt,
     innertopmargin=10pt,
     innerbottommargin=10pt,
     skipabove=\dimexpr\topsep+\ht\strutbox\relax,
     skipbelow=\topsep,
   ]}
  {\end{mdframed}
   \vspace{10pt}%
  }

% --- 12. Python-Integration ---
% (Deaktiviert in dieser Version, aktiviere bei Bedarf)
% \usepackage{pythontex}
% \usepackage[makestderr]{pythontex}

% --- 13. Literaturverwaltung ---
\usepackage{csquotes}
\usepackage[backend=biber, style=alphabetic, citestyle=alphabetic]{biblatex}
\addbibresource{bibliography.bib}

% --- 14. Hyperlinks ---
\usepackage{hyperref}
\hypersetup{
  colorlinks   = true,
  urlcolor     = blue,
  linkcolor    = blue,
  citecolor    = blue,
  frenchlinks  = true
}

% --- 15. Absatzeinstellungen ---
\usepackage[parfill]{parskip}
\sloppy

% --- 16. Umgebungen ---
\usepackage{thmtools}

\newcommand{\CustomHeading}[3]{%
  \par\medskip\noindent%
  \textbf{#1 #2} \textnormal{(#3)}.\enskip%
}


\newenvironment{DEF}[2]{\begin{unitbox}\CustomHeading{Definition}{#1}{#2}}{\end{unitbox}}
\newenvironment{PROP}[2]{\begin{unitbox}\CustomHeading{Proposition}{#1}{#2}}{\end{unitbox}}
\newenvironment{THEO}[2]{\begin{unitbox}\CustomHeading{Theorem}{#1}{#2}}{\end{unitbox}}
\newenvironment{LEM}[2]{\begin{unitbox}\CustomHeading{Lemma}{#1}{#2}}{\end{unitbox}}
\newenvironment{KORO}[2]{\begin{unitbox}\CustomHeading{Corollar}{#1}{#2}}{\end{unitbox}}
\newenvironment{REM}[2]{\begin{unitbox}\CustomHeading{Remark}{#1}{#2}}{\end{unitbox}}
\newenvironment{EXA}[2]{\begin{unitbox}\CustomHeading{Example}{#1}{#2}}{\end{unitbox}}
\newenvironment{STUD}[2]{\begin{unitbox}\CustomHeading{Study}{#1}{#2}}{\end{unitbox}}
\newenvironment{CONC}[2]{\begin{unitbox}\CustomHeading{Concept}{#1}{#2}}{\end{unitbox}}
\newenvironment{OTH}[2]{\begin{unitbox}\CustomHeading{Other}{#1}{#2}}{\end{unitbox}}
\newenvironment{EXE}[2]{\begin{unitbox}\CustomHeading{Exercise}{#1}{#2}}{\end{unitbox}}
\newenvironment{MOT}[2]{\begin{unitbox}\CustomHeading{Motivation}{#1}{#2}}{\end{unitbox}}
\newenvironment{PROOF}[2]{\begin{unitbox}\CustomHeading{Proof}{#1}{#2}}{\end{unitbox}}


% --- Unit Umgebung für Source-Inhalte ---
\usepackage{mdframed}
\newmdenv[
  linewidth=1pt,
  topline=false,
  bottomline=false,
  rightline=false,
  leftmargin=0cm,
  rightmargin=0cm,
  skipabove=10pt,
  skipbelow=10pt,
  innertopmargin=0.5\baselineskip,
  innerbottommargin=0.5\baselineskip,
  backgroundcolor=gray!10,
  linecolor=gray
]{unitbox}

\newenvironment{unit}[1]
  {\begin{unitbox}\textbf{Unit #1}\par\smallskip}
  {\end{unitbox}}

% --- 17. Graphik ---
\usepackage{graphicx}
\graphicspath{ {./images/} }

% --- 00. Titel und Autor ---
\title{Eichfeldtheorie 1}
\author{Tim Jaschik}
\date{\today}

\begin{document}

\maketitle
\rule{\textwidth}{0.5pt}
\begin{abstract}
Kurze Beschreibung …
\end{abstract}
\rule{\textwidth}{0.5pt}
\vspace{0.5cm}

\tableofcontents

\pagebreak


\section{1 Faserbündel}
Soweit nichts anderes gesagt ist, sind in dieser Vorlesung Mannigfaltigkeiten und Abbildungen stets als differenzierbar vorausgesetzt.

\subsection*{1.1 Definition}

\begin{DEF}{BÜN-K25a-02-01}{Lokale triviale Faserung mit typischen Fasern auf Mfk}
Seien $E, M$ und $F$ differenzierbare Mannigfaltikeiten und $\pi: E \rightarrow M$ eine differenzierbare Abbildung. Dann heißt ( $E, \pi, M$ ) eine lokal triviale Faserung mit typischer Faser $F$, wenn es zu jedem $x \in M$ eine offene Umgebung $U$ gibt und einen Diffeomorphismus $\varphi: \pi^{-1}(U):=E \mid U \rightarrow$ $U \times F$, sodass
$$
\begin{array}{lll}
E \mid U & \xrightarrow{\varphi} & U \times F \\
\pi \searrow & & \swarrow p r_{1} \\
& U &
\end{array}
$$
kommutiert. Man spricht auch von der lokal trivialen Faserung $E \rightarrow M$ oder $E$.
\end{DEF}



\begin{DEF}{BÜN-K25a-02-02}{Vektorraumbündel}
Ist $F=\mathbb{R}^{k}$ und ist $\pi^{-1}(x)$ ein $k$-dimensionaler Vektorraum und $p r_{2} \circ$ $\left.\varphi\right|_{\pi^{-1}(x)}: \pi^{-1}(x) \rightarrow \mathbb{R}^{k}$ ein Isomorphismus, so heißt $E$ ein Vektorraumbündel der Dimension $k$.
\end{DEF}

\subsection*{1.2 Beispiele}


\begin{EXA}{BÜN-K25a-02-03}{Projektion von Kreuzprodukt ist eine lokal triviale Faserung}
$p r_{1}: U \times F \rightarrow U$ ist eine lokal triviale Faserung.
\end{EXA}


\begin{EXA}{BÜN-K25a-02-04}{Tangentialbündel mit differenzierbarer Struktur ist Vektorraumbündel}
$T M:=\bigcup_{x \in M} T_{x} M \rightarrow M$ mit der üblichen differenzierbaren Struktur ist ein $\operatorname{dim} M$-dimensionales Vektorraumbündel. Denn ist ( $U, h$ ) eine Karte für $M$ und $\left(\partial_{1}^{(h)}, \ldots, \partial_{n}^{(h)}\right)$ Koordinatenbasis auf $U$, so ist
$$
\bigcup_{x \in U} T_{x} U \rightarrow U \times \mathbb{R}^{n}, \sim \sum_{i=1}^{n} a_{i}(x) \partial_{i}^{(h)}(x) \mapsto\left(x, a_{1}(x), \ldots, a_{n}(x)\right) .
$$
eine Bündelkarte.
\end{EXA}



\begin{EXA}{BÜN-K25a-02-05}{Vektorraumbündel zu $S^1$}
Sei $U:=[0,1] / 0 \sim 1 \cong S^{1}$. $E:=[0,1] \times \mathbb{R} /(0, t) \sim(1,-t) \neq S^{1} \times \mathbb{R}$. Dann ist $\pi: E \rightarrow U,[(x, t)] \mapsto[x]$ ein Vektorraumbündel:

Ist $x \neq[0]$, so wähle $U=(0,1) \subset M$. Dann ist
$$
\pi^{-1}(U)=\{(x, t) \mid x \in(0,1), t \in \mathbb{R}\} \underset{\varphi}{\cong} U \times \mathbb{R}
$$

Ist $x=[0]$, so wähle $U=M \backslash\left\{\frac{1}{2}\right\}$ und
$$
\varphi: \pi^{-1}(U) \rightarrow U \times \mathbb{R},\left\{[(x, t)] \left\lvert\, x \neq \frac{1}{2}\right., t \in \mathbb{R}\right\} \mapsto \begin{cases}([x], t), & 0 \leq x<\frac{1}{2} \\ ([x],-t) & \frac{1}{2}<x \leq 1\end{cases}
$$
\end{EXA}


\begin{EXA}{BÜN-K25a-02-06}{Lokale triviale Faserung über $S^1$}
$S^{1} \rightarrow S^{1}, z \mapsto z^{2}$ ist eine lokal triviale Faserung mit $F=\mathbb{Z}_{2}$. (Übungsaufgabe: Was ist $\varphi$ ?)
\end{EXA}

\subsection*{1.3 Definition}


\begin{DEF}{BÜN-K25a-02-07}{Lokale triviale Faserung als Tripel von Totalraum, Basisraum, Bündelprojektion mit typischen Fasern}
Sei $(E, \pi, M)$ eine lokal triviale Faserung wie in 1.1. Dann heißt $E$ Totalraum, $M$ Basis, $\pi$ Bündelprojektion und $F$ typische Faser.
\end{DEF}

\begin{DEF}{BÜN-K25a-02-08}{Reale Fasern in lokal trivialen Faserungen}
Für jedes $x \in M$ heißt $E_{x}=\pi^{-1}(x)$ reale Faser an der Stelle $x$.
\end{DEF}

\begin{DEF}{BÜN-K25a-02-09}{Bündelkarten für offene Teilmengen der Basis}
Für $U \subset M$ offen heißt $\varphi: E \mid U \rightarrow U \times F$ Bündelkarte
\end{DEF}


\begin{DEF}{BÜN-K25a-02-10}{Bündelatlas für lokale triviale Faserungen}
$$
\left\{\left(U_{\lambda}, \varphi_{\lambda}\right) \mid\left(U_{\lambda}, \varphi_{\lambda}\right) \text { Bündelkarte }, \bigcup_{\lambda \in \Lambda} U_{\lambda}=M\right\}
$$
heißt Bündelatlas.
\end{DEF}


\begin{DEF}{BÜN-K25a-02-11}{Faserkarte am Punkt x im Basisraum}
Die Abbildung $\varphi_{x}: E_{x} \rightarrow F, \varphi_{x}:=p r_{2} \circ \varphi \mid E_{x}$ heißt Faserkarte.
\end{DEF}

\begin{DEF}{BÜN-K25a-02-12}{Bündelkartenwechsel zwischen Bündelkarten}
Sind $(U, \varphi)$ und $(V, \psi)$ Bündelkarten, so heißt die Abbildung
$$
\omega: U \cap V \rightarrow \operatorname{Diffeo}(F), x \mapsto \psi_{x} \circ \varphi_{x}^{-1}
$$
der Bündelkartenwechsel zwischen $\varphi$ und $\psi$.
\end{DEF}


\begin{DEF}{BÜN-K25a-02-13}{G-Faserbündel mit Liegruppen als Strukturgruppen}
Ist $G$ eine Liegruppe und $G \times F \rightarrow F$ eine $G$-Aktion, und gibt es zu jedem Bündelkartenwechsel $\omega$ eine differenzierbare Abbildung
$$
g: U \cap V \rightarrow G \text { mit } \omega(x)(f)=g(x) f
$$
so heißt ( $E, \pi, M$ ) ein $G$-Faserbündel mit Strukturgruppe $G$.
\end{DEF}



\begin{DEF}{BÜN-K25a-02-14}{Prinzipalbüdel / Hauptfaserbündel}
Ist ( $E, \pi, M$ ) ein $G$-Faserbündel mit typischer Faser $G$ und der durch die Linksmultiplikation mit $G$ gegebenen $G$-Aktion, so heißt ( $E, \pi, M$ ) ein Prinzipalbündel oder Hauptfaserbündel.
\end{DEF}



\subsection*{1.4 Bemerkung}


\begin{REM}{BÜN-K25a-02-15}{Beziehung zwischen Vektorraumbündeln und GL-Faserbündeln}
Ist $E \xrightarrow{\pi} M$ ein $k$-dimensionales Vektorraumbündel, so ist der Bündelkartenwechsel zwischen zwei Bündelkarten stets eine differenzierbare Abbildung
$$
\omega: U \cap V \rightarrow G L(k, \mathbb{R}) \text { bzw. } G L(k, \mathbb{C}) .
$$
d.h., $E$ ist ein $G L(k, \mathbb{R})$ - bzw. $G L(k, \mathbb{C})$-Faserbündel. Umgekehrt ist jedes $G L(k, \mathbb{R})$-Faserbündel mit typischer Faser $\mathbb{R}^{k}$ ein Vektorraumbündel.
\end{REM}

\subsection*{1.5 Definition}

\begin{DEF}{BÜN-K25a-02-16}{(Differenzierbare) (Lokale) Schnitte in lokal trivialen Faserungen}
Ist $E \xrightarrow{\pi} M$ eine lokal triviale Faserung, $U \subseteq M$, so heißt eine differenzierbare Abbildung $\sigma: U \rightarrow E$ (differenzierbarer) lokaler Schnitt, falls $\pi \circ \sigma=\mathrm{id}_{U}$. $\sigma$ heißt Schnitt, falls zusätzlich $U=M$ gilt. Den Raum der differenzierbaren Schnitte bezeichnet man mit $\Gamma E$.
\end{DEF}



\subsection*{1.6 Beispiele}

\begin{DEF}{BÜN-K25a-02-17}{Raum der differenzierbaren lokalen Schnitte}
$\Gamma(M \times F)=\left\{\sigma: M \rightarrow M \times F \mid \sigma(x)=(x, f(x)), f \in C^{\infty}(M, F)\right\} \cong$ $C^{\infty}(M, F)$.
\end{DEF}

\begin{EXA}{BÜN-K25a-02-18}{Raum der diff. lokalen Schnitte in Kreuzprodukten}
$\Gamma T M=\{$ differenzierbare Vektorfelder auf $M\}$.
\end{EXA}

\begin{EXA}{BÜN-K25a-02-19}{Raum der diff. Lokalen Schnitte im Tangentialbündel}
Jedes Vektorraumbündel hat einen Schnitt $\sigma: M \rightarrow E, x \mapsto O_{x} \in E_{x}$.
\end{EXA}

$T M$ hat im Allgemeinen keinen Schnitt, der nirgends verschwindet, z.B. hat jedes Vektorfeld auf $M=S^{2}$ eine Nullstelle. Aber in $T T^{2}$ existieren zwei an jeder Stelle linear unabhängige Schnitte.\\

\begin{EXA}{BÜN-K25a-02-21}{Im Tangentialbündel existiert kein diff. Schnitt, der nirgends verschwindet}
Für $S^{1} \rightarrow S^{1}, z \mapsto z^{2}$ gibt es keinen Schnitt.
\end{EXA}

\subsection*{1.7 Bemerkungen}

\begin{REM}{BÜN-K25a-02-23}{Raum der diff Schnitte in Vektorraumbündeln ist der Vektorraum von glatten Abbildungen auf M}
Ist $E \rightarrow M$ ein Vektorraumbündel, so ist $\Gamma E$ ein $C^{\infty}(M)$ Vektorraum.
\end{REM}



\begin{REM}{BÜN-K25a-02-24}{Für Bündelkarten in Vektorraumbündeln existieren k lokale Schnitte, die an jeder Stelle eine Basis der realen Faser bilden}
Ist $E \rightarrow M$ ein $k$-dimensionales Vektorraumbündel und $\varphi: E \mid U \rightarrow$ $U \times \mathbb{R}^{k}$ eine Bündelkarte, so gibt es $k$ lokale Schnitte auf $U$, die an jeder Stelle $x \in U$ eine Basis von $E_{x}$ bilden, nämlich $\sigma_{j}\left(x^{\prime}\right)=\varphi^{-1}\left(x^{\prime}, e_{j}\right)$ für $x^{\prime} \in U$.
\end{REM}


\begin{REM}{BÜN-K25a-02-25}{k lokale Schnitte, die bei Punkt eine Basis der Faser bilden, induzieren eine Bündelkarte}
Umgekehrt: Sind auf $U \in M k$ Schnitte gegeben, die an der Stelle eine Basis bilden, so ist auf $U$ eine Bündelkarte durch
$$
E \mid U \rightarrow U \times \mathbb{R}^{k}, e=\Sigma a_{i}(x) \sigma_{i}(x) \mapsto\left(x, a_{i}(x), \ldots, a_{k}(x)\right)
$$
definiert.
\end{REM}


\begin{REM}{BÜN-K25a-02-26}{Bündelkarten in G-Prinzipalbündeln induzieren lokale Schnitte}
Sind $P \rightarrow M$ ein $G$-Prinzipalbündel und $\varphi: P \mid U \rightarrow U \times G$ eine Bündelkarte, so ist $\sigma: U \rightarrow P, x \mapsto \varphi^{-1}(x, 1)$ ein lokaler Schnitt.
\end{REM}

\subsection*{1.8 Definition}


\begin{DEF}{BÜN-K25a-02-27}{Präbündel mit Strukturgruppe G zu Liegruppe G, Mfk, (disj) Vereinigung von punktweise Mfk und Projektion}
Seien $G$ eine Liegruppe, $M$ eine Mannigfaltigkeit, $F$ eine $G$-Mannigfaltigkeit und sei für jedes $x \in M$ eine Mannigfaltigkeit $E_{x} \cong F$ gegeben. Sei $E:=$ $\bigcup_{x \in M} E_{x}$ und $\pi: E \rightarrow M$ die kanonische Projektion. Dann heißt $(E, \pi, M)$ ein Präbündel mit Strukturgruppe $G$, falls es um jedes $x \in M$ eine offene Umgebung $U$ gibt und eine bijektive Abbildung
$$
\varphi: E \mid U=\pi^{-1}(U) \rightarrow U \times F
$$
sodass
$$
\begin{array}{rlll}
\varphi: E \mid U & \longrightarrow & U \times F & \text { kommutiert } \\
\pi \searrow & & \swarrow p r & \\
& U & &
\end{array}
$$
für je zwei solcher Abbildungen $\varphi, \psi$ eine differenzierbare Abbildung $g_{\varphi, \psi}$ : $U \cap V \rightarrow G$ existiert mit
$$
\varphi \circ \psi^{-1}(x, v)=\left(x, g_{\varphi, \psi}(x) v\right)
$$
\end{DEF}

Die Definitionen aus 1.3 werden entsprechend übertragen, z.B. heißt $\varphi$ dann Präbündelkarte.

\subsection*{1.9 Satz}
\begin{PROP}{BÜN-K25a-02-28}{Für Präbündel $(E,\pi,M)$ existiert auf E genau einem Topologie und differenzierbare Struktur, sodass $(E,\pi,M)$ ein Faserbündel mit Strukturgruppe G wird und Präbündelkarten Bündelkarten werden}
Ist $(E, \pi, M)$ ein Präbündel, dann existiert auf $E$ genau eine Topologie und differenzierbare Struktur, sodass ( $E, \pi, M$ ) ein Faserbündel mit Strukturgruppe $G$ wird und die Präbündelkarten Bündelkarten werden.
\end{PROP}

Beweisidee: Man definiert $\Omega \subseteq E$ offen: $\Leftrightarrow$ für jede Präbündelkarte $\varphi$ : $E \mid U \rightarrow U \times F$ ist $\varphi(\Omega \cap E \mid U) \subseteq U \times F$ offen.\\
Ist $(U, \varphi)$ Präbündelkarte von $E$ und o.B.d.A. $(U, h)$ Mannigfaltigkeitskarte für $M$, so definiert man

$$
\Phi: E \mid U \xrightarrow{\varphi} U \times F \xrightarrow{h} U^{\prime} \times F .
$$

Zusammen mit Karten für $F$ erhält man dann eine differenzierbare Struktur auf $E$.

\subsection*{1.10 Beispiele}
\begin{EXA}{BÜN-K25a-02-29}{Bündelstruktur von Tangentialbündel als Ergebnis der Konstruktion von Präbündeln}
Die Bündelstruktur $T M$ wurde wie in Satz 1.9 definiert.
\end{EXA}


\begin{EXA}{BÜN-K25a-02-30}{Präbündel zum GL-Prinzipalbündel}
Sei $E_{x}:=\left\{\left(v_{1}, \ldots, v_{n}\right) \mid\left(v_{1}, \ldots, v_{n}\right)\right.$ Basis von $\left.T_{x} M\right\}, P_{G L}:=\bigcup_{x \in M} E_{x}$. Dann ist $P_{G L}$ ein Präbündel und wird in kanonischer Weise ein $G L(n, \mathbb{R})$ Prinzipalbündel (Beweis in den Übungen).
\end{EXA}


\begin{EXA}{BÜN-K25a-02-31}{Präbündel zum $O(n)$-Prinzipalbündel für Riemannische Mfk}
Analog, falls ( $M, g$ ) Riemannsch ist und $E_{x}:=\left\{\left(v_{1}, \ldots, v_{n}\right) \mid\left(v_{1}, \ldots, v_{n}\right)\right.$ ist Orthonormalbasis von $\left.T_{x} M\right\}$, so ist $P_{O(n)}:=\bigcup_{x \in M} E_{x}$ in kanonischer Weise ein $O(n)$-Prinzipalbündel.
\end{EXA}


\subsection*{1.11 Korollar}

\begin{KORO}{BÜN-K25a-02-32}{Direkte Summe von Vektorraumbündeln ergeben Prävektorraumbündel}
Sind $E \rightarrow M$ und $F \rightarrow M$ zwei Vektorraumbündel der Dimension $k$ und $\ell$, so ist
$$
E \oplus F:=\bigcup_{x \in M} E_{x} \oplus F_{x}
$$
ein $k+\ell$-dimensionales Prävektorraumbündel, also in kanonischer Weise ein Vektorraumbündel.
\end{KORO}

\section*{Beweis:}
Seien ( $U, \varphi$ ) und ( $U, \psi$ ) Bündelkarten für $E$ und $F$. Dann ist durch

$$
\begin{aligned}
\phi: \bigcup_{x \in U} E_{x} \oplus F_{x} & \rightarrow U \times\left(\mathbb{R}^{k} \oplus \mathbb{R}^{\ell}\right) \\
(e, f) & \mapsto\left(\pi_{E}(e)=\pi_{F}(f)=: x, \varphi_{x}(e), \psi_{x}(f)\right)
\end{aligned}
$$
eine Präbündelkarte gegeben.

Sind $(U, \tilde{\varphi})$ und $(U, \tilde{\varphi})$ weitere Bündelkarten mit Bündelkartenwechsel $\omega: U \rightarrow$ $G L(k), \eta: U \rightarrow G L(\ell)$, so ist der Bündelkartenwechsel zwischen $\varphi$ und $\tilde{\varphi}$ durch
$$
U \rightarrow G L(k+\ell), x \mapsto\left(\begin{array}{cc}
\omega(x) & 0 \\
0 & \eta(x)
\end{array}\right)
$$
gegeben, also differenzierbar.


\subsection*{1.12 Beispiele}

\begin{EXA}{BÜN-K25a-02-33}{Hom-Raum für Homomorphismen zwischen Vektorraumbündeln sind Vektorraumbündel}
Sind $E, F$ Vektorraumbündel, so auch

$\operatorname{Hom}(E, F)=\bigcup_{x \in M} \operatorname{Hom}\left(E_{x}, F_{x}\right)$.
\end{EXA}

\begin{EXA}{BÜN-K25a-02-34}{Mult}
$\operatorname{Mult}^{k}(E, F)$.
\end{EXA}

\begin{EXA}{BÜN-K25a-02-35}{Sym}
$\operatorname{Sym}^{k}(E)$.
\end{EXA}

\begin{EXA}{BÜN-K25a-02-36}{Alt}
$\operatorname{Alt}^{k}(E)$ usw.
\end{EXA}

Die Präbündelkartenwechsel können als Übungsaufgabe konstruiert werden.

\subsection*{1.13 Definition}


\begin{DEF}{BÜN-K25a-02-37}{Bündelmetrik auf Totalraum ist ein Schnitt in $Sym^2(E)$, sodass g pw. positiv definit }
Eine Bündelmetrik auf $E$ ist ein Schnitt $g$ in $\operatorname{Sym}^{2}(E)$, sodass $g(x)$ positiv definit ist für jedes $x \in M$.
\end{DEF}

\subsection*{1.14 Bemerkungen}

\begin{EXA}{BÜN-K25a-02-38}{Riemannische Metrik als Bündelmetrik im Tangentialbündel}
Eine Riemannsche Metrik auf $M$ ist eine Bündelmetrik auf $T M$.
\end{EXA}


\begin{EXA}{BÜN-K25a-02-39}{$\Gamma (Alt^k(TM))$}
$\Omega^{k} M=\Gamma \operatorname{Alt}^{k}(T M)$.
\end{EXA}



\subsection*{1.15 Definition}

\begin{DEF}{BÜN-K25a-02-40}{Vektorraumbündel vom endlichen Typ}
Ein Vektorraumbündel $E \rightarrow M$ heißt von endlichem Typ, wenn es ein Vektorraumbündel $F \rightarrow M$ gibt und ein $N \in \mathbb{N}$, sodass $E \oplus F=M \times \mathbb{R}^{N}$ ist.
\end{DEF}


\subsection*{1.16 Beispiel}

\begin{EXA}{BÜN-K25a-02-41}{Tangentialbündel von $S^n$ ist von endlichem Typ}
$T S^{n}$ ist von endlichem Typ, denn
$$
T S^{n} \oplus\left(S^{n} \times \mathbb{R}\right)=S^{n} \times \mathbb{R}^{n-1},(v,(x, \lambda)) \mapsto v+\lambda x .
$$
\end{EXA}

\subsection*{1.17 Definition}


\begin{DEF}{BÜN-K25a-02-43}{Bündelisomorphismus}
Seien $\pi_{i}: E_{i} \rightarrow B_{i}, i=1,2$ Faserbündel mit Strukturgruppe $G$ und typischer Faser $F, f_{0}: B_{1} \rightarrow B_{2}$ eine differenzierbare Abbildung, dann heißt eine differenzierbare Abbildung $f: E_{1} \rightarrow E_{2}$ eine Bündelabbildung über $f_{0}$, falls $f_{0} \circ \pi_{1}=\pi_{2} \circ f$ gilt und $f$ bezüglich Bündelkarten durch differenzierbare Abbildungen $\gamma: U \rightarrow G$, genauer durch
$$
U \times F \rightarrow V \times F,(x, v) \mapsto\left(f_{0}(x), \gamma(x) v\right)
$$
gegeben ist.

Ist zusätzlich $f_{0}=\mathrm{id}$, so heißt $f$ ein Bündelisomorphismus.
\end{DEF}

\begin{DEF}{BÜN-K25a-02-44}{Trivialisierung von Totalraum}
Ein Bündelisomorphismus $f: E \rightarrow B \times F$ heißt eine Trivialisierung von $E$.
\end{DEF}


\begin{DEF}{BÜN-K25a-02-45}{Vektorraumbündelabbildung über diff. Abbildungen zwischen Vektorraumbündeln}
Seien jetzt $E_{i} \rightarrow B_{i}$ zwei Vektorraumbündel $f_{0}: B_{1} \rightarrow B_{2}$ eine differenzierbare Abbildung und $f: E_{1} \rightarrow E_{2}$ eine differenzierbare Abbildung über $f_{0}$, so heißt $f$ eine Vektorraumbündelabbildung über $f$, falls $f_{x}:=f \mid E_{1, x}$ linear ist für jedes $x \in M$.
\end{DEF}


\begin{DEF}{BÜN-K25a-02-46}{Vektorraumbündelisomorphismus}
Ist $f_{0}=\mathrm{id}$, so heißt $f$ ein Vektorraumbündelhomomorphismus, es ist dann $f \in \Gamma \operatorname{Hom}\left(E_{1}, E_{2}\right)$.
\end{DEF}

\subsection*{1.18 Beispiel}


\begin{EXA}{BÜN-K25a-02-47}{Differential von glatten Abbildungen zw. Tangentialbündel von Mfk ist eine Vektorraumbündelabbildung über glatte Abbildung $f$}
Ist $f: M \rightarrow N$ eine differenzierbare Abbildung, so ist $d f: T M \rightarrow T N$ eine Vektorraumbündelabbildung über $f$.
\end{EXA}

\subsection*{1.19 Definition und Satz}
\begin{DEF}{BÜN-K25a-02-48}{Induzierte Bündel durch Abbildungen}
Ist $f: N \rightarrow M$ eine differenzierbare Abbildung, $E \rightarrow M$ ein Faserbündel mit Strukturgruppe $G$, so ist
$$
p r_{1}: f^{*} E=\left\{(x, e) \mid x \in N, e \in E_{f(x)}\right\} \rightarrow N
$$
ein Faserbündel mit Strukturgruppe $G$ über $N$, das von $f$ induzierte Bündel. Schnitte in $f^{*} E$ heißen Schnitte in $E$ längs $f$.
\end{DEF}


\begin{PROP}{BÜN-K25a-02-49}{Schnitte in induzierten Bündeln längs $f$}
Ist $f: N \rightarrow M$ eine differenzierbare Abbildung, $E \rightarrow M$ ein Faserbündel mit Strukturgruppe $G$, so ist
$$
p r_{1}: f^{*} E=\left\{(x, e) \mid x \in N, e \in E_{f(x)}\right\} \rightarrow N
$$
ein Faserbündel mit Strukturgruppe $G$ über $N$, das von $f$ induzierte Bündel. Schnitte in $f^{*} E$ heißen Schnitte in $E$ längs $f$.
\end{PROP}


\section*{Beweis:}
Sei $x \in N$ und $(U, \varphi)$ Bündelkarte um $f(x)$ für $E$. Dann ist $\left(f^{-1}(U), \tilde{\varphi}\right)$ mit

$$
\begin{aligned}
\tilde{\varphi}:\left\{(x, e) \mid x \in f^{-1}(U), e \in E_{f(x)}\right\}=p r_{1}^{-1}\left(f^{-1}(U)\right) & \longrightarrow f^{-1}(U) \times F \\
(x, e) & \mapsto\left(x, \varphi_{f(x)}(e)\right)
\end{aligned}
$$

eine Präbündelkarte gegeben. Die Kartenwechsel sind dieselben wie die für $E \xrightarrow{\pi} M$, also differenzierbar.


\subsection*{1.20 Beispiel}
\begin{EXA}{BÜN-K25a-02-50}{Menge der Vektorfelder längs Kurven}
Ist $\gamma: I \rightarrow M$ differenzierbar, so ist $\Gamma\left(\gamma^{*} T M\right)$ die Menge der Vektorfelder längs $\gamma$.
\end{EXA}


\subsection*{1.21 Beispiel}
\begin{EXA}{BÜN-K25a-02-51}{Vektorraumbündel bzgl Grassmann-Mfk}
Sei $G_{k, N}=O(N) /(O(k) \times O(N-k))$ die Grassmann-Mannigfaltigkeit und
$$
\gamma_{k, N}=\left\{(E, v) \in G_{k, N} \times \mathbb{R}^{N} \mid v \in E\right\}
$$
dann ist $\gamma_{k, N} \rightarrow G_{k, N}$ ein $k$-dimensionales Vektorraumbündel. $\gamma_{k, N}$ ist offenbar von endlichem Typ und ein Vektorraumbündel $E \rightarrow M$ ist genau dann von endlichem Typ, wenn es ein $N$ gibt und eine Abbildung $f: M \rightarrow G_{k, N}$, sodass $E=f^{*} \gamma_{k, N}$ ist.
\end{EXA}

\subsection*{1.22 Bemerkung}


\begin{REM}{BÜN-K25a-02-52}{Bündelabbildungen bzgl induzierte Bündel}
Die Abbildung $\hat{f}: f^{*} E \rightarrow E,(x, e) \mapsto e$ ist eine Bündelabbildung über $f$, sodass $\hat{f} \mid\left(f^{*} E\right)_{x}:\left(f^{*} E\right)_{x} \rightarrow E_{f(x)}$ ein Isomorphismus ist für jedes $x$. Ist $g$ : $\widetilde{E} \rightarrow E$ eine weitere Bündelabbildung über $f$, so ist
$$
\tilde{g}: \widetilde{E} \rightarrow f^{*} E, \tilde{e} \mapsto(\pi(\tilde{e}), g(\tilde{e}))
$$
ein Bündelisomorphismus mit $\hat{f} \circ \tilde{g}=g$.
\end{REM}

\subsection*{1.23 Satz}

\begin{PROP}{BÜN-K25a-02-67}{Pullback-Faserbündel längs homotopen Abbildungen sind isomorph}
Sei $E \rightarrow B$ ein Faserbündel, seien $f_{0}, f_{1}: X \rightarrow B$ homotope Abbildungen, dann ist $f_{0}^{*} E \cong f_{1}^{*} E$.
\end{PROP}

\subsection*{1.24 Korollar}
\begin{KORO}{BÜN-K25a-02-68}{Faserbündel über kontrahierbarer Basis sind trivial}
Ist $B$ zusammenziehbar, also die Identität homotop zur konstanten Abbildung $c: B \rightarrow B, x \mapsto p$, so ist $E \rightarrow B$ trivial, denn $E=\mathrm{id}^{*} E \cong c^{*} E=B \times E_{p}$.
\end{KORO}

\section*{Beweis des Satzes:}

\begin{PROOF}{BÜN-K25a-02-69}{P: Faserbündel über kontrahierbarer Basis sind trivial}
Seien $\xi, \xi^{\prime}$ Faserbündel mit Strukturgruppe $G$ und typischer Faser F. Sei
$$
\operatorname{Iso}_{G}\left(\xi_{x}, \xi_{x}^{\prime}\right):=\left\{f: \xi_{x} \rightarrow \xi_{x}^{\prime} \mid \psi_{x} \circ f \circ \varphi_{x}^{-1}(v)=g \cdot v \text { für ein } g \in G\right\}
$$
Dann ist $\bigcup_{x \in B} \operatorname{Iso}_{G}\left(\xi_{x}, \xi_{x}^{\prime}\right)$ in kanonischer Weise ein Faserbündel mit typischer Faser $G$ und Strukturgruppe $G \times G$, die Schnitte in $\operatorname{Iso}_{G}\left(\xi_{x}, \xi_{x}^{\prime}\right)$ sind gerade die Bündelisomorphismen. 

Benutze nun die Homotopiehochhebungseigenschaft: Ist $(X, A)$ ein $CW$-Paar, (z.B. $(M, \partial M)$ ), $Y \rightarrow[0,1] \times X$ eine lokal triviale Faserung, $\sigma_{0}$ ein Schnitt in $Y \mid((0 \times X) \cup([0,1] \times A))$, dann ist $\sigma$ zu einem Schnitt in $Y$ fortsetzbar, vgl. z.B. Husemoller Fiber Bundles oder Hatcher Algebraic Geometry.

Wende dies an auf das Bündel $\operatorname{Iso}_{G}\left([0,1] \times f_{0}^{*} E, h^{*} E\right)$, wobei $h:[0,1] \times X \rightarrow B$ die Homotopie zwiswchen $f_{0}$ und $f_{1}$ ist. Dann ist $\sigma(0, x)=\operatorname{id}_{E_{f_{0}(x)}}$ fortsetzbar zu einem Schnitt $\sigma$, und $\sigma$ und $\sigma(1, \cdot)$ ist der gesuchte Bündelisomorphismus.
\end{PROOF}

\subsection*{1.25 Definition}

\begin{DEF}{BÜN-K25a-02-54}{Induzierte Bündel bei Einbettungen von UnterMfk}
Ist $i: M_{0} \rightarrow M$ die Einbettung einer Untermannigfaltigkeit $M_{0}$ in $M$, so schreibt man statt $i^{*} E$ auch $\left.E\right|_{M_{0}}$.
\end{DEF}

\subsection*{1.26 Definition}

\begin{DEF}{BÜN-K25a-02-55}{Untervektorraumbündel}
Sei $(E, \pi, M)$ ein $k$-dimensionales Vektorraumbündel. Eine Teilmenge $E_{0} \subseteq E$ heißt ein $m$-dimensionales Untervektorraumbündel, falls es um jedes $x \in M$ eine Bündelkarte $(U, \varphi)$ gibt, sodass $\varphi\left(\pi^{-1}(U) \cap E_{0}\right)=U \times \mathbb{R}^{m} \times 0$ ist.
\end{DEF}

\subsection*{1.27 Bemerkungen}


\begin{REM}{BÜN-K25a-02-56}{Untervektorraumbündel sind Vektorraumbündel}
Ein $m$-dimensionales Untervektorraumbündel ist offenbar ein $m$-dimensionales Vektorraumbündel.
\end{REM}


\begin{REM}{BÜN-K25a-02-57}{Quotienten-Räume bzgl Untervektorraumbündel sind Vektorraumbündel}
Ist $E_{0} \subset E$ ein Untervektorraumbündel, so ist $E / E_{0}$ ein Vektorraumbündel.
\end{REM}

\begin{REM}{BÜN-K25a-02-58}{Untervektorraumbündel bzgl Bündelmetrik}
Ist $M$ mit einer Bündelmetrik $g$ versehen, so ist
$$
E_{0}^{\perp}:=\bigcup_{x \in M} E_{0, x}^{\perp}=\left\{e \in E_{x}: g(\tilde{e}, e)=0 \text { für alle } \tilde{e} \in E_{0, x}\right\}
$$
ein Untervektorraumbündel und es gilt: $E_{0}^{\perp} \cong E / E_{0}$.
\end{REM}




\begin{REM}{BÜN-K25a-02-59}{Tangentialbündel von UnterMfk sind Untervektorraumbündel}
Ist $M_{0}$ eine Untermannigfaltigkeit, dann ist $\left.T M_{0} \subset T M\right|_{M_{0}}:=\bigcup_{x \in M_{0}} T_{x} M$ ein $\operatorname{dim} M_{0}$-dimensionales Untervektorraumbündel von $\left.T M\right|_{M_{0}}$.
\end{REM}

\begin{REM}{BÜN-K25a-02-60}{Normalenbündel von UnterMfk}
$N M_{0}:=\left.T M\right|_{M_{0}} / T M_{0}$ heißt das Normalenbündel von $M_{0}$. Ist $M$ eine Riemannsche Mannigfaltigkeit, so ist $N M_{0} \cong T M_{0}^{\perp}$.
\end{REM}

\subsection*{1.28 Satz}

\begin{PROP}{BÜN-K25a-02-61}{Rang-Satz für Vektorraumhomomorphismen: Konstanter Rang impliziert ker und im sind Untervektorraumbündel}
Ist $f: E \rightarrow F$ ein Vektorraumhomomorphismus von einem Vektorraumbündel der Dimension $k$ und $m$ und ist $\operatorname{rg} f_{x}=$ const.\\
Dann ist $\operatorname{ker} f:=\bigcup_{x \in M} \operatorname{ker} f_{x} \subset E$ ein Untervektorraumbündel von $E$ und Bild $f:=\bigcup_{x \in M}$ Bild $f_{x} \subset F$ ein Untervektorraumbündel von $F$.
\end{PROP}

\section*{Beweis:}
Ohne Einschränkung sei $E=X \times \mathbb{R}^{k}, F=X \times \mathbb{R}^{m}$. Sei $r g f_{x}=r$.

\begin{enumerate}
  \item Fall: $k \geq m$. Wir zeigen: $\operatorname{ker} f$ ist ein Untervektorraumbündel. OBdA $k=m$, sonst ersetze $F$ durch $F \oplus\left(X \times \mathbb{R}^{k-m}\right)$ und $f$ durch $(f, 0)$. Sei $x \in X$. Nach Wahl geeigneter Karten ist
\end{enumerate}

$$
f_{x}=\left(\begin{array}{cccccc}
1 & & & & & \\
& \ddots & & & & \\
& & 1 & & & \\
& & & 0 & & \\
& & & & \ddots & \\
& & & & & 0
\end{array}\right)
$$

Setze

$$
P=\left(\begin{array}{cccccc}
0 & & & & & \\
& \ddots & & & & \\
& & 0 & & & \\
& & & 1 & & \\
& & & & \ddots & \\
& & & & & 1
\end{array}\right)
$$

Also ist $f_{x}+P \in G L(k, \mathbb{R})$. Definiere für $x^{\prime} \in X:(f+P)_{x^{\prime}}=f_{x^{\prime}}+P$. Dann existiert eine Umgebung $U$ von $x$ so, dass $(f+P)_{x^{\prime}} \in G L(k, \mathbb{R})$ für alle $x^{\prime} \in U$. Für $x^{\prime} \in U$ ist $(f+P)_{x^{\prime}}\left(\operatorname{ker} f_{x^{\prime}}\right)=0 \times \mathbb{R}^{k-r}$.\\
" $\subseteq$ " folgt, da $f_{x^{\prime}}\left(\operatorname{ker} f_{x^{\prime}}\right)=0$ und die Gleichheit folgt dann aus Dimensionsgründen.\\
2. Fall: $k \leq m$. Wir zeigen: Bild $f$ ist ein Untervektorraumbündel.

OBdA $k=m$, sonst ersetze $E$ durch $E \oplus\left(X \times \mathbb{R}^{m-k}\right)$ und $f$ durch $f \circ p r_{k}$. Seien $f_{x}$ und $P$ wie im ersten Fall. Dann ist $\left(f_{x^{\prime}}+P\right)\left(\mathbb{R}^{r} \times 0\right) \subseteq \operatorname{Bild} f_{x^{\prime}}$. Da $\left(f_{x^{\prime}}+P\right): \mathbb{R}^{m} \rightarrow \mathbb{R}^{m}$ ein Isomorphismus ist, folgt die Gleichheit aus Dimensionsgründen. Folglich ist $\left(f_{x^{\prime}}+P\right)^{-1}$ die gesuchte Bündelkarte.\\
3. Fall: $k \leq m$. Wir zeigen: $\operatorname{ker} f \subset E$ ist ein Untervektorraumbündel.

Wende den 1. Fall auf $f: E \rightarrow \operatorname{Bild} f$ an.\\
4. Fall: $k \geq m$ : Wir zeigen: Bild $f \subset F$ ist ein Untervektorraumbündel. Wende den 2. Fall auf $\left.E\right|_{\operatorname{kerf}}$ an.



\subsection*{1.29 Korollar}
\begin{KORO}{BÜN-K25a-02-62}{Charakterisierung von Vektorraumbündeln von endlichem Typ}
Ein Vektorraumbündel $E \rightarrow M$ ist genau dann von endlichem Typ, wenn eine der beiden äquivalenten Bedingungen erfüllt ist:
\begin{enumerate}
  \item Es existiert ein surjektiver Vektorraumbündelhomomorphismus $M \times \mathbb{R}^{N} \rightarrow$ $E$.
  \item Es existieren $N$ Schnitte $s_{1}, \ldots, s_{n} \in \Gamma E$, sodass $\left(s_{1}(x), \ldots, s_{n}(x)\right)$ für jedes $x \in M$ die Faser $E_{x}$ erzeugt.
\end{enumerate}
\end{KORO}


\subsection*{1.30 Definition}
\begin{DEF}{BÜN-K25a-02-63}{Reduktionen von Faserbündeln mit Strukturgruppe bzgl abgeschlossener Untergruppe}
Ist $E \rightarrow M$ ein Faserbündel mit Strukturgruppe $G$ und $G_{0} \subseteq G$ eine abgeschlossene Untergruppe. Dann sagt man: $E$ besitzt eine Reduktion auf $G_{0}$ oder eine $G_{0}$-Bündelstruktur, falls es einen Bündelatlas gibt, dessen Bündelkartenwechsel Werte in $G_{0}$ annehmen.
\end{DEF}

\subsection*{1.31 Beispiel}
\begin{EXA}{BÜN-K25a-02-64}{Charakterisierung von orientierten Mfk}
$M$ ist genau dann orientierbar, wenn $T M$ eine $G L^{+}(n, \mathbb{R})$-Bündelstruktur hat. (Übungsaufgabe)
\end{EXA}

Die folgenden beiden Sätze geben oft eine einfache Möglichkeit zu entscheiden, wann eine Bündelstruktur vorliegt.

\subsection*{1.32 Satz (Ehresmannscher Faserungssatz)}

\begin{PROP}{BÜN-K25a-02-65}{Ehresmannscher Faserungssatz: Totalräume mit eigentlich regulären Abbildungen in zusammenhängenden Basisraum implizieren eine lokale triviale Faserung}
Ist $X$ zusammenhängend, $p: E \rightarrow X$ eine eigentliche reguläre Abbildung, so ist $E$ eine lokal triviale Faserung.
\end{PROP}

\subsection*{1.33 Satz (von Hermann)}

\begin{PROP}{BÜN-K25a-02-66}{Satz von Hermann für vollständige zusammenhängende Riemannische Mannigfaltigkeiten}
Ist ( $M, g$ ) vollständig und $\tilde{M}$ zusammenhängend, dann ist jede Riemannsche Submersion $\pi: M \rightarrow \tilde{M}$ ein Faserbündel.
\end{PROP}



\subsection*{1.34 Übungsaufgaben}
\begin{enumerate}
  \item Auf $S^{1} \subset \mathbb{C}$ betrachte man die durch $z \sim \bar{z}$ und $1 \sim-1$ definierte Äquivalenzrelation. Zeigen Sie:\\
a) $S^{1} / \sim$ ist homöomorph zu $S^{1}$.\\
b) Die kanonische Projektion $S^{1} \rightarrow S^{1} / \sim$ ist keine lokal triviale Faserung.
  \item Es sei $E:=\left\{(x, v) \in \mathbb{R P}^{n} \times \mathbb{R}^{n+1} \mid v \in x\right\}$. Geben Sie einen Bündelatlas für die lokal triviale Faserung $E \rightarrow \mathbb{R} \mathbb{P}^{n}$ an.
  \item Es sei $M$ eine $n$-dimensionale differenzierbare Mannigfaltigkeit mit differenzierbarer Struktur D. Beschreiben Sie den kanonischen Prä-Bündelatlas für $\mathrm{Alt}^{K} T M$, d.h. für die Familie $\left\{\mathrm{Alt}^{K} T_{p} M\right\}_{p \in M}$.
  \item $\pi_{1}: E_{1} \rightarrow M$ und $\pi_{2}: E_{2} \rightarrow M$ seien lokal triviale Faserungen mit typischen Fasern $F_{1}$ und $F_{2}$. Dann heißt $E_{1} \times_{M} E_{2}:=\left\{\left(e_{1}, e_{2}\right) \in E_{1} \times\right.$ $\left.E_{2} \mid \pi_{1}\left(e_{1}\right)=\pi_{2}\left(e_{2}\right)\right\}$ mit der kanonischen Abildung $\pi: E_{1} \times_{M} E_{2} \rightarrow M$ das gefaserte Produkt oder das Faserprodukt oder das Produkt über $M$ der beiden Faserungen. Konstruieren Sie aus Bündelatlanten $\mathcal{A}_{1}$ und $\mathcal{A}_{2}$ für die Faktoren einen Bündelatlas $\mathcal{A}$ für das gefaserte Produkt.
  \item a) Es sei ( $M,\langle$,$\rangle ) eine pseudo-Riemannsche n$-dimensionale differenzierbare Mannigfaltigkeit, der Index des Skalarproduktes sei $n-k$. Zeigen Sie, dass diejenigen Bündelkarten des Tangentialbündels, deren Faserkarten Isometrien sind, zusammen eine $O(k, n-k)$-Bündelstruktur für $T M \rightarrow M$ bilden.\\
b) Zeigen Sie: Eine differenzierbare Mannigfaltigkeit besitzt genau dann eine Metrik vom Index $n-k$, wenn TM eine $O(k, n-k)$-Bündelstruktur hat.
  \item Eine differenzierbare Mannigfaltigkeit heißt parallelisierbar, wenn ihr Tangentialbündel trivial ist. Zeigen Sie, dass jede Lie-Gruppe parallelisierbar ist, $S^{2 k}$ für $k>1$ aber nicht. Beweisen Sie ferner, dass ( $S^{n} \times$ $\mathbb{R}) \oplus T S^{n}$ für alle $n$ trivial ist.
  \item Es sei $E \rightarrow B$ ein $n$-dimensionales Vektorraumbündel. Bestimmen Sie die Übergangsfunktionen für $\mathrm{Alt}^{n} E$ aus denen für $E$.
  \item Für Vektorraumbündel über einer Mannigfaltigkeit zeigen Sie: Ist
\end{enumerate}

$$
0 \rightarrow E^{\prime} \rightarrow E \rightarrow E^{\prime \prime} \rightarrow 0
$$

eine exakte Sequenz von Bündelhomomorphismen, so ist $E \cong E^{\prime} \oplus E^{\prime \prime}$.\\
9) Für Bündelhomomorphismen $f: E \rightarrow E$ folgere man aus dem Ranglemma:\\
a) Ist $f \circ f=f$, so sind Kern $f$ und Bild $f$ Teilbündel von $E$.\\
b) Ist $f \circ f=\mathrm{Id}$, so ist $\operatorname{Fix}(f):=\{e \in E \mid f(e)=e\}$ ein Teilbündel von $E$.



\pagebreak

\section{Prinzipalbündel}
\subsection*{2.1 Definition}
Eine $G$-Aktion $G \times M \rightarrow M$ heißt frei, falls aus $g x=x$ für ein $x \in M$ folgt $g=1$.\\
Eine $G$-Aktion $G \times M \rightarrow M$ heißt effektiv, falls gilt: wirkt $g$ als Identität, so ist $g=1$.\\
Eine $G$-Aktion $G \times M \rightarrow M$ heißt transitiv, falls gilt: Zu jedem Paar $(x, y) \in$ $M \times M$ existiert ein $g \in G$ mit $g x=y$.

\subsection*{2.2 Notiz}
Ist $\phi: G \times M \rightarrow M$ eine freie transitive $G$-Aktion, so ist $G \cong M$.

\section*{Beweis:}
Sei $x \in M$. Die Abbildung $o_{x}: G \rightarrow M, g \mapsto g x$ ist bijektiv und differenzierbar und $\left.d o_{x}\right|_{1}$ ist injektiv, also ist $\left.d o_{x}\right|_{g}$ bijektiv für jedes $g \in G$, also ist $o_{x}$ ein Diffeomorphismus (vgl. Lee 7.15: differenzierbare Abbildung von konstantem Rang!).

\subsection*{2.3 Satz}
Auf dem Totalraum eines $G$-Prinzipalbündels gibt es eine $G$-Rechtsaktion, die auf den Fasern frei und transitiv ist. Die Faserkarten $\varphi_{x}$ sind bezüglich dieser Aktion $G$-rechtsäquivariant, d.h., $\varphi_{x}(p g)=\varphi_{x}(p) g$.

\section*{Beweis:}
Sei $p \in P_{x}$. Setze $p g=\varphi_{x}^{-1}\left(\varphi_{x}(p) g\right)$. Dies ist wohldefiniert, denn ist $\psi$ eine weitere Karte, so ist $\psi_{x}(p)=\omega(x) \varphi_{x}(p)$ für ein $\omega(x) \in G$, also $\psi^{-1}\left(x, \psi_{x}(p) g\right)=$ $\psi^{-1}\left(x, \omega(x) \varphi_{x}(p) g\right)=\varphi^{-1}\left(x, \varphi_{x}(p) g\right)$. Offenbar ist $\varphi_{x}$ rechtsäquivariant und die Operation auf der realen Faser frei und transitiv, weil die $G$-Aktion auf $G$ dies ist.

Für Prinzipalbündel sind Schnitte Bündelkarten, genauer:

\subsection*{2.4 Lemma}
Ein Prinzipalbündel ist genau dann trivial, wenn es einen Schnitt besitzt.

\section*{Beweis:}
Ist $\sigma$ ein Schnitt, so setze $P \rightarrow M \times G, \sigma(x) g \mapsto(x, g)$. Umgekehrt: Setze $\sigma(x)=\varphi^{-1}(x, 1)$.

\subsection*{2.5 Bemerkung}
Ist $P \rightarrow M$ ein $G$-Prinzipalbündel, und ist $\varphi$ die durch $\sigma$ defnierte Bündelkarte und $s(x)=\sigma(x) g(x)$, dann wird $s$ bezüglich $\varphi$ durch $g$ beschrieben.\\
Ist $\tilde{\sigma}(x)=\sigma(x) a(x)$ ein weiterer Schnitt und $\tilde{\varphi}$ die durch $\tilde{\sigma}$ gegebene Bündelkarte, so ist $s(x)=\tilde{\sigma}(x) a^{-1}(x) g(x)$, wird also bezüglich $\tilde{\sigma}$ durch $a^{-1} g$ beschrieben.

\subsection*{2.6 Bemerkung}
Die rechtsäquivarianten Abbildungen $f: G \rightarrow G$ sind genau die Linksmultiplikationen mit $f(1) \in G$, denn $f(g)=f(1 g)=f(1) g$.

\subsection*{2.7 Lemma}
Sei $G$ eine Liegruppe, $P \rightarrow M$ ein Faserbündel mit Strukturgruppe $G$. Dann ist äquivalent:\\
a) $P \rightarrow M$ ist ein $G$-Prinzipalbündel.\\
b) Auf dem Totalraum von $P$ ist eine $G$-Rechtsaktion gegeben, die auf den Fasern frei und transitiv operiert.

\section*{Beweis:}
a) $\Rightarrow$ b) $\checkmark$\\
b) $\Rightarrow$ a) Nach 2.2 ist die typische Faser $G$, wie in 2.5 sind rechtsinvariante Faserkarten gegeben, und rechtsäquivariante Bündelkartenwechsel sind nach 2.6. durch Linksmultiplikation mit $g \in G$ gegeben.

\subsection*{2.8 Satz}
Sei $G$ eine Liegruppe, $P \rightarrow M$ eine differenzierbare Abbildung. Dann ist äquivalent:\\
a) $P \rightarrow M$ ist ein $G$-Prinzipalbündel.\\
b) $P$ ist eine Rechts- $G$-Mannigfaltigkeit der Dimension $\operatorname{dim} G+\operatorname{dim} M$. Die $G$-Wirkung ist fasertreu und auf den Fasern frei und transitiv. Ferner existiert eine Überdeckung $\left(U_{\lambda}\right)_{\lambda \in \Lambda}$ von $M$ und lokale Schnitte $s_{\lambda}: U_{\lambda} \rightarrow$ $P$.

\section*{Beweis:}
a) $\Rightarrow$ b) $\checkmark$\\
b) $\Rightarrow$ a) Da lokale Schnitte existieren, ist $\pi: P \rightarrow M$ eine surjektive Submersion, also ist $\pi^{-1}(x) \subseteq P$ eine $\operatorname{dim} G$-dimensionale Untermannigfaltigkeit von $P$,\\
also ist nach $2.2 P_{x} \cong G$. Defniere $\psi: U \times G \rightarrow P \mid U,(x, g) \mapsto s(x) g$. Dies definiert eine Bündelkarte $\varphi=\psi^{-1}$. Die Abbildung $\psi$ ist bijektiv, differenzierbar und rechts- $G$-äquivariant. Das Differential

$$
d \psi_{(x, g)}(X, v)=d R_{g}(X)+d L_{s(x) g}(v)
$$

ist injektiv, denn ist $d \psi_{(x, g)}(X, v)=0$, so ist

$$
X=d \pi_{s(x) g} \circ d \psi_{(x, g)}(X, v)=0
$$

da $(\pi \circ \psi)(x, g)=x$, also ist auch $d L_{s(x)}(v)=0$, also $v=0$ (da $G$ frei operiert).

\subsection*{2.9 Beispiel}
Die Hopffaserung $S^{3} \rightarrow \mathbb{C} P^{1},\left(z_{1}, z_{2}\right) \mapsto\left[z_{1}: z_{2}\right]$ ist ein $S^{1}$-Prinzipalfaserbündel.\\
Insbesondere für kompakte Gruppen ist folgender Satz auch nützlich:

\subsection*{2.10 Satz}
Operiert $G$ frei und eigentlich auf $M$, so ist $M / G$ eine differenzierbare Mannigfaltigkeit und $M \rightarrow M / G$ ein $G$-Prinzipalbündel.

\section*{Beweis:}
Lee, 9.16 (Quotient manifold theorem).

\subsection*{2.11 Definition}
a) Sei $\pi_{1}: P \rightarrow M$ ein $G$-Prinzipalbündel, $\pi_{2}: Q \rightarrow N$ ein $H$-Prinzipalbündel, $\alpha: G \rightarrow H$ ein Homomorphismus von Liegruppen, $f_{0}: M \rightarrow N$ differenzierbar. Dann heißt $f: P \rightarrow Q$ ein $\alpha$-Prinzipalbündelhomomorphismus über $f_{0}$, falls $\pi_{2} \circ f=f_{0} \circ \pi_{1}$ gilt und $f(p g)=f(p) \alpha(g)$ ist.\\
c) Ist $f_{0}=\mathrm{id}$, so heißt ( $P, f$ ) eine $\alpha$-Version von Q. Ist zusätzlich $G=H$ und $\alpha=\mathrm{id}$, so spricht man von einem Prinzipalbündelisomorphismus.\\
d) Zwei $\alpha$-Versionen heißen äquivalent, falls es einen $G$-Prinzipalbündelisomorphismus $g$ über id gibt, sodass $\tilde{f}=g \circ f$ ist.\\
Eine Äquivalenzklasse von $\alpha$-Versionen heißt $\alpha$-Struktur. Ist $G \subseteq H$ eine Untergruppe, $\alpha: G \rightarrow H$ die Inklusion, so heißt eine $\alpha$-Struktur auch eine Reduktion von $Q$ auf $G$.

\subsection*{2.12 Beispiel}
Sei $(M, g)$ eine Riemannsche Mannigfaltigkeit. Dann beschreibt $P_{O(n)}(M) \xrightarrow{i}$ $P_{G L}(M)$ eine $O(n)$-Reduktion von $P_{G L}(M)$.

\subsection*{2.13 Lemma}
a) Ist $Q$ ein $H$-Prinzipalbündel, $G \subseteq H$ eine abgeschlossene Untergruppe und $P \subseteq Q$ eine Teilmenge, sodass gilt:

\begin{enumerate}
  \item Die $G$-Rechtsaktion ist aus $P_{x}:=Q_{x} \cap P$ frei und transitiv.
  \item Es existieren eine offene Überdeckung $\left(U_{\lambda}\right)_{\lambda \in \Lambda}$ von $M$ und lokale Schnitte $\sigma_{\lambda}: U_{\lambda} \rightarrow Q$ mit $\sigma\left(U_{\lambda}\right) \subseteq P$.
\end{enumerate}

Dann ist ( $P, i$ ) eine $G$-Reduktion von $M$.\\
Umgekehrt: Besitzt $Q$ eine $G$-Reduktion $(P, f)$, so erfüllt $f(P) \subseteq Q$ die Bedingungen 1. und 2.\\
b) Ein $H$-Prinzipalbündel $Q$ besitzt genau dann eine Reduktion auf eine abgeschlossene Untergruppe $G \subseteq H$, falls es einen Bündelatlas von Q gibt, dessen Übergangsfunktionen Bilder in $G$ annehmen.

\section*{Beweis von $\mathbf{b}$ ):}
$" \Leftarrow$ " Ist $A=\left\{U_{\lambda}, \varphi_{\lambda} \mid \lambda \in \Gamma\right\}$ ein Bündelatlas wie gefordert, dann setze

$$
P=\bigcup_{\lambda \in \Lambda} \varphi_{\lambda}^{-1}\left(U_{\lambda} \times G\right)
$$

Dies ist eine Teilmenge von Q wie in 1 .\\
" $\Rightarrow$ " Ist $f: P \rightarrow Q$ eine Reduktion auf $G$, so existieren nach a) lokale Schnitte $\sigma_{\lambda}: U_{\lambda} \rightarrow Q$ mit $\sigma_{\lambda}\left(U_{\lambda}\right) \subseteq f(P)$. Diese definieren Bündelkarten für $P$. Ist $\sigma_{\mu}$ ein weiterer solcher Schnitt, so gilt für $x \in U_{\lambda} \cap U_{\mu}$ :

$$
\sigma_{\mu}(x)=\sigma_{\lambda}(x) g_{\lambda \mu}(x)
$$

für $g_{\lambda \mu}: U_{\lambda} \cap U_{\mu} \rightarrow G$. Damit ist der entsprechende Bündelkartenwechsel durch $g_{\lambda \mu}^{-1}: U_{\lambda} \cap U_{\mu} \rightarrow G, x \mapsto\left(g_{\lambda \mu}(x)\right)^{-1}$ gegeben.

\subsection*{2.14 Satz}
Seien $\pi: Q \rightarrow M$ ein $H$-Prinzipalbündel und $G \subseteq H$ eine abgeschlossene Untergruppe. Dann hat man eine kanonische Bijektion zwischen der Menge aller Reduktionen auf $G$ und der Menge aller Schnitte in $Q / G \rightarrow M$.

\section*{Beweis:}
" $\Rightarrow$ "Da $H \rightarrow H / G$ ein $G$-Prinzipalbündel ist (Beweis später), ist $\pi_{1}: Q / G \rightarrow$ $M$ ein Faserbündel mit typischer Faser $H / G$ und $Q \rightarrow Q / G$ ein $G$-Prinzipalbündel. Sei $\sigma$ ein Schnitt in $Q / G$. Dann ist $\sigma^{*} Q$ ein $G$-Prinzipalbündel über $M$, und die gesuchte Reduktion ist durch ( $\sigma^{*} Q, \hat{\sigma}$ ) gegeben. Dabei ist $\hat{\sigma}: \sigma^{*} Q \rightarrow Q$ als Bündelhomomorphismus über $M$ aufzufassen, denn $\hat{\sigma}(x, q)=q$ für $q \in$ $\pi_{1}^{-1}(\sigma(x)) \subseteq \pi^{-1}(x)$.\\
„ $\Leftarrow$ "Sei $f: P \rightarrow Q$ ein Repräsentant einer $G$-Reduktion, also $f(p g)=f(p) g$ für $g \in G$. Dann ist durch $\sigma(x)=[f(p)]_{G}$ für ein $p \in P_{x}$ ein Schnitt in $Q / G$ wohldefiniert. Ist $\tilde{f}: \tilde{P} \rightarrow Q$ ein weiterer Repräsentant, also $\tilde{f}=f \circ F$ für einen Bündelisomorphismus $F: \tilde{P} \rightarrow P$, so existiert für $\tilde{p} \in \tilde{P}_{x}$ und $p \in P_{x}$ ein $g \in G$ mit

$$
\tilde{f}(\tilde{p})=f(F(\tilde{p}))=f(p g)=f(p) \alpha(g)
$$

für ein $g \in G$ also

$$
[\tilde{f}(p)]=[f(p)]
$$

\subsection*{2.15 Korollar}
Ist $G \subseteq H$ eine abgeschlossene Untergruppe einer Liegruppe mit $H / G \cong \mathbb{R}^{n}$ für ein $n \in \mathbb{N}$, so besitzt jedes $H$-Prinzipalbündel eine $G$-Reduktion

\subsection*{2.16 Beispiel}
a) $G L(n, \mathbb{R}) / O(n) \cong \mathbb{R}^{\frac{1}{2} n(n+1)}$.\\
b) $O(k, n-k) / O(k) \times O(n-k) \cong \mathbb{R}^{m}$ für $1 \leq k \leq n-1$ und $m=k(n-k)$.

\subsection*{2.17 Übungsaufgabe}
Verallgemeinern Sie die Hopf-Faserung $S^{3} \rightarrow \mathbb{C P}^{1}$ zu $S^{2 n+1} \rightarrow \mathbb{C} \mathbb{P}^{n}$, (Bündelatlas angeben).


\pagebreak

\section{An- und Abmontieren der Faser}

\subsection*{3.1 Bemerkung und Definition}
Sind $G$ eine Liegruppe, $X$ eine Rechts- $G$-Mannigfaltigkeit und $F$ eine Links-$G$-Mannigfaltigkeit, dann ist auf $X \times F$ eine $G$-Aktion durch $G \times(X \times F) \rightarrow$ $X \times F,(g,(x, v)) \mapsto\left(x g, g^{-1} v\right)$ gegeben. Wir schreiben $(X \times F) / \sim=: X \times{ }_{G} F$. Auf $X \times_{G} F$ existieren zwei kanonische Projektionen:

$$
\begin{gathered}
X \times_{G} F \rightarrow G \backslash F \\
\text { und } X \times_{G} F \rightarrow X / G .
\end{gathered}
$$

Schreiben wir $X \times_{G} F$ und reden von der Projektion, so meinen wir $X \times_{G} F \rightarrow$ $X / G$.

\subsection*{3.2 Beispiel}
a) $G \times_{G} F \rightarrow F,[g, v] \mapsto g v$ ist ein Homöomorphismus und $G \times_{a} F$ trägt eine eindeutig bestimmte differenzierbare Struktur, sodass dies ein Diffeomorphismus ist.\\
b) Ist $X$ eine freie transitive Rechts- $G$-Mannigfaltigkeit, dann ist für jedes $x_{0} \in X$

$$
X \times_{G} F \rightarrow F,\left[x_{0} g, v\right] \mapsto g v
$$

ebenfalls ein Diffeomorphismus.\\
c) Ist $X$ eine freie transitive Rechts- $G$-Mannigfaltigkeit, so ist $U \times X$ kanonisch ebenfalls eine Rechts- $G$-Mannigfaltigkeit und $(U \times X) \times{ }_{G} F \rightarrow$ $U \times F$ ein Diffeomorphismus.

\subsection*{3.3 Satz und Definition}
Sei $P \rightarrow M$ ein $G$-Prinzipalbündel und $F$ eine $G$-Mannigfaltigkeit. Dann ist $P \times_{G} F$ in kanonischer Weise ein Faserbündel über $M$ mit Strukturgruppe $G$ und typischer Faser $F$. Ein Atlas von $P$ induziert einen Atlas von $P \times_{G} F$, dessen Kartenwechsel durch die selben Übergangsfunktionen gegeben sind. Wir nennen den Funktor $P \rightsquigarrow P \times_{G} F$ das Anmontieren oder Assoziieren der Faser $F$ an $P$. Ist die $G$-Aktion auf $F$ mit $\alpha$ bezeichnet, $\alpha: G \times F \rightarrow F$, dann schreibt man auch $P \times_{G} F \equiv P \times_{\alpha} F$.

\section*{Beweis:}
Ist $(U, \varphi)$ eine Bündelkarte für $P$, so definiere eine Präbündelkarte durch:


\begin{gather*}
\tilde{\varphi}: \bigcup_{x \in U} P_{x} \times_{G} F=: P \times_{G} F \mid U \xrightarrow{\varphi}(U \times G) \times_{G} F \rightarrow U \times F  \tag{3.1}\\
{[p, v] \mapsto[\varphi(p), v]=\left[\left(\pi(p), \varphi_{x}(p)\right), v\right] \mapsto\left(\pi(p), \varphi_{x}(p) v\right) .}
\end{gather*}


\subsection*{3.4 Beispiel}
a) $P_{G L} \times_{G L} \mathbb{R}^{n} \stackrel{\cong}{\rightrightarrows} T M,\left[\left(v_{1}, \ldots, v_{n}\right),\left(a_{1}, \ldots, a_{n}\right)\right] \mapsto \sum_{i=1}^{n} a_{i} v_{i}$\\
b) Ist $M$ Riemannsch, so ist ebenso ein Isomorphismus der $O(n)$-Reduktion definiert: $P_{O(n)} \times{ }_{O(n)} \mathbb{R}^{n} \stackrel{\cong}{\rightrightarrows} T M$

\subsection*{3.5 Bemerkung und Notation}
a) Ist $f: P \rightarrow Q$ ein $G$-Prinzipalbündelhomomorphismus, so ist durch $f_{*}$ : $P \times_{G} F \rightarrow Q \times_{G} F,[p, v] \mapsto[f(p), v]$ eine $G$-Bündelabbildung gegeben.\\
b) Ist $\alpha: G \rightarrow H$ ein Homomorphismus, $f: P \rightarrow Q$ eine $\alpha$-Version und $F$ eine $H$-Mannigfaltigkeit, so ist $f_{*}: P \times_{\alpha} F \rightarrow Q \times_{H} F,[p, v] \mapsto[f(p), v]$ ein Isomorphismus von lokal trivialen Faserungen. und bezüglich geeigneten Bündelkarten durch id gegeben. Fassen wir auch $P \times{ }_{\alpha} F$ als Faserbündel mit Strukturgruppe $H$ auf, so ist $f_{*}$ ein $H$-Faserbündelisomorphismus.\\
c) Ist insbesondere ist $f: P_{0} \rightarrow P$ eine $G_{0}$-Reduktion von $P$, so ist $P_{0} \times G_{0}$ $F \rightarrow P \times_{G} F$ ein $G$-Bündelisomorphismus.

\section*{Beweis:}
Ist $\sigma$ ein lokaler Schnitt von $P$, so ist $f \circ \sigma$ ein lokaler Schnitt in $Q$. Für die durch $\sigma$ und $f \circ \sigma$ induzierten Karten $\varphi$ und $f_{*} \varphi$ wie in (3.1) gilt:

$$
\widetilde{f_{*} \varphi} \circ f \circ \tilde{\varphi}^{-1}: \varphi\left(P \times_{G} F \mid U\right) \rightarrow f_{*} \varphi\left(Q \times_{G} F \mid U\right),(x, v) \mapsto(x, v) .
$$

\subsection*{3.6 Erinnerung}
Sei $E \rightarrow M$ ein Faserbündel mit Strukturgruppe $G$, so sagt man, $E$ besitzt eine Reduktion auf $G_{0}$, wenn $E$ einen Bündelatlas besitzt, dessen Bündelabbildungen durch Abbildungen nach $G_{0}$ gegeben sind.

\subsection*{3.7 Bemerkung}
Ist $E=P \times_{G} F$ und $P_{0} \rightarrow P$ eine Reduktion von $P$, so besitzt $E$ eine Reduktion auf $G_{0}$ (nämlich $P_{0} \times{ }_{G_{0}} F$ ).\\
Betrachten wir jetzt den Spezialfall, dass wir als typische Faser wieder eine Liegruppe anmontieren.

\subsection*{3.8 Lemma}
Ist $\alpha: G \rightarrow H$ ein Liegruppenhomomorphismus, so ist $Q=P \times_{\alpha} H$ in kanonischer Weise ein $H$-Prinzipalbündel, und $P$ ist eine $\alpha$-Version von $Q$.\\
Umgekehrt: Ist $P \rightarrow Q$ eine $\alpha$-Version, so ist $Q \cong P \times_{\alpha} H$.

\section*{Beweis:}
$Q$ ist ein $H$-Prinzipalbündel nach Lemma 2.8, denn $H$ operiert auf den Fasern von $P \times_{\alpha} H$ frei und transitiv von rechts. Ein $\alpha$-Prinzipalbündelhomomorphismus $P \rightarrow Q$ ist durch $f: P \rightarrow Q, p \mapsto[p, 1]$ gegeben, denn $f(p g)=[p g, 1]=[p, \alpha(g)]=[p, 1] \alpha(g)$.\\
Ist $P \xrightarrow{f} Q$ eine $\alpha$-Version, so ist durch $\tilde{f}: P \times{ }_{\alpha} H \rightarrow Q,[p, h] \mapsto f(p) \underset{\sim}{h}$ der gesuchte Bündelisomorphismus wohldefiniert, denn $\tilde{f}([p, h] \tilde{h})=\tilde{f}([p, h \tilde{h}])=$ $f(p) h \tilde{h}=\tilde{f}([p, h]) \tilde{h}$.

\subsection*{3.9 Bemerkung}
a) Ist $\sigma$ ein lokaler Schnitt in $P$, so definiert $\sigma$ eine Bündelkarte für $P \times_{G}$ $F \mid U \cong U \times F$. Also wird ein Schnitt $s$ in $P \times{ }_{G} F \mid U$ nach Wahl von $\sigma$ durch eine Abbildung $U \rightarrow F$ beschrieben.\\
Genauer: Ist $s(x)=[\sigma(x), v(x)]$, so wird $s$ bezüglich $\sigma$ durch $v$ beschrieben. Ist $\tilde{\sigma}$ ein anderer Schnitt von $P, \tilde{\sigma}=\sigma g$, so wird $s$ bezüglich $\tilde{\sigma}$ durch $g^{-1} v$ beschrieben:

$$
s(x)=[\sigma(x), v(x)]=\left[\tilde{\sigma} g^{-1}, v(x)\right]
$$

b) Ein Schnitt $s$ in $P \times{ }_{G} F$ definiert eine Abbildung $\bar{s}: P \rightarrow F$ mit $\bar{s}(p g)=$ $g^{-1} s(p)$ durch

$$
s(x)=[p(x), \bar{s}(p(x))]=\left[p(x) g(x), g^{-1}(x) s(p(x))\right]
$$

Umgekehrt: Ist $f: P \rightarrow F$ eine Funktion mit $f(p g)=g^{-1} f(p)$, so definiert $f$ einen Schnitt in $P \times_{G} F$ durch $x \mapsto[p(x), f(p(x))]$ mit $p(x) \in$ $P_{x}$ beliebig.

Jetzt betrachten wir die umgekehrte Konstruktion, die Faserbündeln Prinzipalbündel zuordnet.

\subsection*{3.10 Definition und Satz}
Sei $F$ eine effektive $G$-Mannigfaltigkeit, $E \rightarrow M$ ein Faserbündel mit typischer Faser $F$ und Strukturgruppe $G$. Sei


\begin{align*}
\operatorname{Iso}_{G}\left(F, E_{x}\right):= & \left\{f_{x}: F \rightarrow E_{x} \mid \varphi_{x} \circ f_{x}=g_{\varphi, f}\right. \\
& \text { für ein } \left.g_{\varphi, f} \in G \text { für (eine und dann) jede Bündelkarte } \varphi\right\} \tag{*}
\end{align*}


Dann ist $\operatorname{Iso}_{G}(F, E)=\bigcup_{x \in M} \operatorname{Iso}_{G}\left(F, E_{x}\right)$ in kanonischer Weise ein $G$-Prinzipalbündel, das $E$ zugrundeliegende Prinzipalbündel. Die Bündelkartenwechsel sind genau die selben, wie die von $E$.

\section*{Beweis:}
Sei $(U, \varphi)$ eine Bündelkarte von $E$. Setze

$$
\bigcup_{x \in U} \operatorname{Iso}_{G}\left(F, E_{x}\right) \rightarrow U \times G, f_{x} \mapsto\left(x, g_{\varphi, f}\right)
$$

mit $g_{\varphi, f}$ wie in $(*)$. Ist $\psi_{x}=\omega(x) \cdot \varphi_{x}$, so ist $g_{\psi, f}=\omega(x) g_{\varphi, f}$.

\subsection*{3.11 Satz}
Ist $P$ ein $G$-Prinzipalbündel, $F$ eine effektive $G$-Mannigfaltigkeit, so ist

$$
\operatorname{Iso}_{G}\left(F, P \times_{G} F\right) \cong P .
$$

Ist $E$ ein Faserbündel mit typischer Faser $F$ und Strukturgruppe $G$, wobei $G$ effektiv auf $F$ operiert, so ist

$$
\operatorname{Iso}_{G}(F, E) \times{ }_{G} F \cong E .
$$

Bezüglich geeigneter Bündelkarten sind die Abbildungen durch id gegeben.

\section*{Beweis:}
Definiere

$$
\alpha: P \rightarrow \operatorname{Iso}_{G}\left(F, P \times_{G} F\right), p \mapsto(v \mapsto[p, v])
$$

und

$$
\beta: \operatorname{Iso}_{G}(F, E) \times_{G} F \rightarrow E,[f, v] \mapsto f(v) .
$$

\subsection*{3.12 Beispiel}
Ist $M$ eine $n$-dimensionale Mannigfaltigkeit, so ist $T M=P_{G L}(M) \times{ }_{G L(n, \mathbb{R})} \mathbb{R}^{n}$ und $P_{G L}(M) \cong \operatorname{Iso}_{G L}\left(\mathbb{R}^{n}, T M\right)$

$$
\left(v_{1}, \ldots, v_{n}\right) \mapsto\left(\left(a_{1}, \ldots, a_{n}\right) \mapsto \sum_{i=1}^{n} a_{i} v_{i}\right)
$$

\subsection*{3.13 Bemerkung}
Besitzt $E \rightarrow M$ eine $G_{0^{-}}$-Reduktion, so ist $\operatorname{Iso}_{G_{0}}(F, E) \subseteq \operatorname{Iso}_{G}(F, E)$ eine $G_{0^{-}}$ Reduktion von $\operatorname{Iso}_{G}(F, E)$.

\subsection*{3.14 Beispiel}
Ist auf $M$ eine Riemannsche Metrik gegeben, so ist

$$
\operatorname{Iso}_{O(n)}\left(\mathbb{R}^{n}, T M\right)=P_{O(n)}(M)
$$

\subsection*{3.15 Lemma}
Ist $f: E \rightarrow \tilde{E}$ ein $G$-Bündelisomorphismus, so ist

$$
f_{*}: \operatorname{Iso}_{G}(F, E) \rightarrow \operatorname{Iso}_{G}(F, \tilde{E}), \alpha \mapsto f \circ \alpha
$$

ein $G$-Prinzipalbündelisomorphismus.

\subsection*{3.16 Sprechweise}
Der Vorgang $E \rightsquigarrow \operatorname{Iso}_{G}(F, E)$ heißt Abmontieren der Faser.\\
Anmontieren von Fasern ist auch für nicht effektive Aktionen wichtig:

\subsection*{3.17 Beispiel und Definition}
Ist $P$ ein $G$-Prinzipalbündel, so heißt $\operatorname{Aut}(P):=P \times_{k o n j} G,[p, g]=\left[p \tilde{g}, \tilde{g}^{-1} g \tilde{g}\right]$ das Bündel der Eichtransformationen. Dies ist ein Bündel mit typischer Faser $G$ und Strukturgruppe $G$, aber kein Prinzipalbündel!\\
Ist $f \in \Gamma \operatorname{Aut}(P)$, so ist $f$ ein Bündelautomorphismus, denn für $[p, g] \in$ $\operatorname{Aut}_{x}(P)$ und $\tilde{p} \in P_{x}$ definiert

$$
[p, g](\tilde{p})=:[\tilde{p} \tilde{g}, g](\tilde{p})=\left[\tilde{p}, \tilde{g} g \tilde{g}^{-1}\right](\tilde{p})=\tilde{p} \tilde{g} g \tilde{g}^{-1}
$$

einen Bündelautomorphismus, denn

$$
[p, g](\tilde{p} a)=[\tilde{p} \tilde{g}, g](\tilde{p} a)=\left[\tilde{p}, \tilde{g} g \tilde{g}^{-1}\right](\tilde{p} a)=\left[\tilde{p} a, a^{-1} \tilde{g} g \tilde{g}^{-1} a\right](\tilde{p} a)=\tilde{p} \tilde{g} g \tilde{g}^{-1} a .
$$

Umgekehrt: Ist $f: P \rightarrow P$ ein Bündelautomorphismus, so ist ein Schnitt $s$ in Aut $(P)$ auf folgende Weise gegeben: Es ist für jedes $p \in P$ ist $f(p)=p g$ für ein eindeutig definiertes $g$. Setze $s(x)=[p, g]$. Dies ist wohldefiniert, denn für $\tilde{p}=p a$ ist

$$
f(\tilde{p})=f(p a)=p g a=\tilde{p} a^{-1} g a \text { und }\left[\tilde{p}, a^{-1} g a\right]=[p, g] .
$$

\subsection*{3.18 Übungsaufgaben}
\begin{enumerate}
  \item Es sei $P \rightarrow B$ ein $G$-Prinzipalfaserbündel und $F$ ein $G$-Raum. Was ist
\end{enumerate}

$$
\operatorname{Iso}_{G}\left(F, P \times_{G} F\right) \rightarrow B
$$

für ein Bündel, wenn die Aktion $G \rightarrow \text{Homöo}(F)$ nicht effektiv ist, sondern einen nichttrivialen Kern $G_{0}$ hat?\\
2) Auf $\mathbb{R}^{n}$ operiere $\mathbb{Z}_{2}$ durch die Involution $x \mapsto-x$. Dann ist $E_{n}:=$ $S^{1} \times \mathbb{Z}_{2} \mathbb{R}^{n}$ für $n \geq 1$ als Faserbündel mit Strukturgruppe $\mathbb{Z}_{2}$ natürlich nicht trivial (weshalb nämlich?). Wir betrachten $E_{n}$ jetzt aber als Vektorraumbündel über $S^{1}$. Zeigen Sie: $E_{2}$ ist trivial, $E_{1}$ aber nicht. Verallgemeinerung für $E_{n}$ mit $n \geq 3$ ?\\
3) Es sei $\alpha: H \rightarrow G$ ein Liegruppenepimorphismus, $K$ sein Kern und $P$ ein $H$-Prinzipalfaserbündel. Zeigen Sie: Das $G$-Prinzipalfaserbündel $P \times{ }_{\alpha} G$ ist genau dann trivial, wenn sich die Strukturgruppe von $P$ auf $K$ reduzieren läßt.\\
4) Bestimmen Sie eine Untergruppe $H \subset G L(n, \mathbb{R})$, sodass gilt: Ein $n$ dimensionales Vektorraumbündel $E$ besitzt genau dann ein eindimensionales Untervektorraumbündel, wenn seine Strukturgruppe auf $H$ reduzierbar ist.\\
5) Zeigen Sie: Eine differenzierbare $n$-dimensionale Mannigfaltigkeit besitzt genau dann eine Metrik vom Index $k$, wenn das Tangentialbündel TM als Summe $T M=\xi \oplus \eta$ eines $k$-dimensionalen Vektorraumbündels $\xi$ und eines ( $n-k$ )-dimensionalen Vektorraumbündels $\eta$ geschrieben werden kann.



\pagebreak

\section{Die Parallelverschiebung}
\subsection*{4.1 Definition}
Sei $E \xrightarrow{\pi} M$ eine lokal triviale Faserung. Eine Zuordnung $\tau$, die jedem Weg $\gamma:[a, b] \rightarrow M$ einen Diffeomorphismus $\tau_{\gamma}: E_{\gamma(a)} \rightarrow E_{\gamma(b)}$ zuordnet, heißt Paralleltransport oder Parallelverschiebung, falls gilt:

\begin{enumerate}
  \item Die Abbildung $\tau$ hängt differenzierbar von $\gamma \mathrm{ab}$, d.h. sind $U \subseteq \mathbb{R}^{n}$ offen, $a, b: U \rightarrow \mathbb{R}$ differenzierbar mit $a(u) \leq b(u)$ und ist $\gamma: \bigcup_{u \in U}\{u\} \times$ $[a(u), b(u)] \rightarrow M$ differenzierbar, so ist
\end{enumerate}

$$
\begin{gathered}
\gamma_{\mathrm{Anf}}^{*} E \rightarrow \gamma_{\mathrm{End}}^{*} E, \\
(u, e) \mapsto\left(u, \tau_{\gamma_{u}}(e)\right)
\end{gathered}
$$

differenzierbar. Dabei sind die Anfangs- und Endkurven $\gamma_{\text {Anf,End }}: U \rightarrow$ $M$ durch $\gamma_{\text {Anf }}(u)=\gamma(u, a(u))$ und $\gamma_{\text {End }}(u)=\gamma(u, b(u))$ gegeben und $\gamma_{u}(t)=\gamma(u, t)$.\\
2. Für $a<c<b$ ist $\tau_{\left.\gamma\right|_{[c, b]}} \circ \tau_{\left.\gamma\right|_{[a, c]}}=\tau_{\gamma}$ (Unterteilbarkeit).\\
3. Ist $\varphi:[c, d] \rightarrow[a, b]$ differenzierbar, $\varphi(c)=a, \varphi(d)=b$, so ist $\tau_{\gamma}=\tau_{\gamma \circ \varphi}$. (Invarianz unter Umparametrisierungen)\\
4. Für $e \in E_{x}$ hängt $\left.\frac{d}{d t}\right|_{t=0} \tau_{\left.\gamma\right|_{[a, t]}}(e)$ nur von $\dot{\gamma}(0)$ ab. (Erstes Ordnungsaxiom)

\subsection*{4.2 Notiz}
Ist $\gamma:[a, b] \rightarrow M$ konstant, so ist $\tau_{\gamma}=\operatorname{id}_{E_{\gamma(a)}}$ nach 2. und 3.

\subsection*{4.3 Definition}
Ist $\gamma:[a, b] \rightarrow M$ eine Kurve, $a \leq s_{i} \leq b$, so schreiben wir:

$$
\gamma_{\left[s_{1}, s_{2}\right]}(t)=\gamma\left((1-t) s_{1}+t s_{2}\right), t \in[0,1]
$$

\subsection*{4.4 Notiz}
Für $s>a$ ist $\tau_{\gamma_{[a, s]}}=\tau_{\left.\gamma\right|_{[a, s]}}$ und $\tau_{\gamma_{[a, a]}}=\operatorname{id}_{E_{\gamma(a)}}$.

\subsection*{4.5 Notation}
Ist $\gamma:[a, b] \rightarrow M$ differenzierbar, $e \in E_{\gamma(a)}$, so heißt $\gamma_{e}(t)=\tau_{\gamma_{[a, t]}}(e)$ die (zu e) hochgehobene Kurve.

\subsection*{4.6 Bemerkung}
Ist $f: M \rightarrow N$ eine differenzierbare Abbildung, $E \rightarrow N$ eine lokal triviale Faserung mit Paralleltransport $\tau$, so ist ein Paralleltransport in $f^{*} E \rightarrow M$ durch

$$
\left(f^{*} \tau_{\gamma}\right)(x, e)=\left(\gamma(b), \tau_{f \circ \gamma}(e)\right)
$$

für $\gamma:[a, b] \rightarrow M$ mit $\gamma(a)=x$ und $e \in E_{f(x)}$ wohldefiniert.

\subsection*{4.7 Definition und Satz}
Ist $\tau$ eine Parallelverschiebung auf $E$ und $\gamma: I \rightarrow M$ eine Kurve in $M$, so ist für jedes $e \in E_{p}$

$$
\dot{\tau}_{e}: T_{p} M \rightarrow T_{e} E, v \mapsto \dot{\gamma}_{e}(0)
$$

wobei $\dot{\gamma}(0)=v$ wohldefiniert (1. Ordnungsaxiom). Der dadurch definierte Vektorraumbündelhomomorphismus

$$
\pi^{*} T M \rightarrow T E,(e, v) \mapsto \dot{\tau}_{e}(v)
$$

heißt infinitesimaler Parallelismus. Es gilt $(d \pi)_{e} \circ \dot{\tau}_{e}=\operatorname{id}_{T_{\pi(e)} M}$, insbesondere ist $\dot{\tau}_{e}$ injektiv.

\section*{Beweis:}
Zu zeigen: $\dot{\tau}_{e}$ ist linear. Wir konstruieren durch radialen Paralleltransport einen Schnitt $s$ in $E$ und zeigen, dass $d s=\tau_{e}$ gilt, genauer:\\
Sei $p \in M,(U, h)$ eine Karte um $p$ mit $h(p)=0$ und $\bar{U}_{1}(0)=h(U)$.\\
Sei $\gamma: U \times[0,1] \rightarrow M ; \gamma(x, t):=h^{-1}(\operatorname{th}(x))$. Dann ist

$$
\begin{gathered}
\gamma_{\mathrm{Anf}}(x)=\gamma(x, 0)=p, \text { also } \gamma_{\mathrm{Anf}}^{*} E=U \times E_{p} \\
\gamma_{\mathrm{End}}(x)=\gamma(x, 1)=x \text { also } \gamma_{\mathrm{End}}^{*} E=E \mid U
\end{gathered}
$$

Die Abbildung $\tau$ : $U \times E_{p} \rightarrow E \mid U,(x, e) \mapsto \tau_{\gamma_{x}}(e)$ ist differenzierbar. Sei $s_{e} \in$ $\Gamma(E \mid U)$ durch $s_{e}(x)=\tau_{\gamma_{x}}(e)$ definiert. Wir zeigen jetzt, dass $d s_{e}(v)=\dot{\tau}_{e}(v)$ gilt, dann folgt sofort: $\dot{\tau}_{e}$ ist linear und $d \pi_{e} \circ \dot{\tau}_{e}=\mathrm{id}$. Es ist $\tau_{\gamma_{[0, t]}}(e)=\tau(\gamma(x, t), e)=s_{e}\left(\gamma_{x}(t)\right)$. Sei $\dot{\gamma}_{x}(0)=v$. Dann ist

$$
d s_{e}(v)=\left.\frac{d}{d t}\right|_{t=0}\left(s_{e} \circ \gamma_{x}\right)=\left.\frac{d}{d t}\right|_{t=0} \tau_{\gamma_{[0, t]}}(e)=\left.\frac{d}{d t}\right|_{t=0}\left(\gamma_{x_{[0, t]}}\right)_{e}=\dot{\tau}_{e}(v)
$$

\subsection*{4.8 Definition}
Der Untervektorraum $H_{e}=\dot{\tau}_{e}\left(T_{\pi(e)} M\right) \subseteq T_{e} E$ heißt der Horizontalraum an der Stelle $e$ des Parallelismus $\tau$. Ein Schnitt $s \in \Gamma E$ heißt horizontal an der Stelle $x$, wenn $d s(v) \in H_{s(x)}$ für alle $v \in T_{x} M$ gilt.

\subsection*{4.9 Notiz}
Die Vereinigung $H=\bigcup_{e \in E} H_{e} \subseteq T E$ ist in kanonischer Weise ein Untervektorraumbündel mit $H \oplus \operatorname{ker} d \pi=T E$. Es ist $\operatorname{ker} d \pi_{e}=T_{e}^{\prime} E:=T_{e} E_{\pi(e)}$.\\
Wir verallgemeinern jetzt den Begriff des Horizontalraums eines Parallelismus, ohne das Vorliegen einer Parallelverschiebung zu fordern.

\subsection*{4.10 Definition}
Wir bezeichnen das Fasertangentialbündel mit $\bigcup T_{e} E_{\pi(e)}=T^{\prime} E$. Ein Zusammenhang ist ein Untervektorraumbündel $H \subseteq T E$, sodass gilt:

$$
H \oplus T^{\prime} E=T E
$$

Ist $H$ ein Zusammenhang auf $E$, so heißt $H_{e}$ auch Horizontalraum an der Stelle $e$, ein Schnitt $s \in \Gamma E$ horizontal and der Stelle $x$, falls $d s_{x}\left(T_{x} M\right) \subseteq H_{s(x)}$ und horizontal, falls $d s_{x}\left(T_{x} M\right) \subseteq H_{s(x)}$ für alle $x \in M$ gilt. Eine Kurve $\gamma: I \rightarrow E$ heißt horizontal and der Stelle $t$, falls $\dot{\gamma}(t) \in H_{\gamma(t)}$ und horizontal, falls $\dot{\gamma}(t) \in$ $H_{\gamma(t)}$ für alle $t \in I$.

\subsection*{4.11 Bemerkung}
In jeder lokal trivialen Faserung gibt es einen Zusammenhang. Man erhält einen Zusammenhang auf $E$, wenn man eine Riemannsche Metrik auf $E$ wählt und $H_{e}:=\left(T_{e}^{\prime} E\right)^{\perp}$ setzt.

\subsection*{4.12 Lemma}
Ist $\tau$ eine Parallelverschiebung in $E$, so ist durch den zugehörigen infinitesimalen Parallelismus ein Zusammenhang gegeben und die Hochhebungen von Kurven sind Horizontalkurven.

\section*{Beweis:}
Zu zeigen: Für $\gamma:[0, L] \rightarrow M$ ist stets $\gamma_{e}(t) \in H_{\gamma_{e}(t)}$ für $t \geq 0$. Dies gilt, denn $\gamma_{e} \mid[t, L]$ ist die Hochhebung von $\gamma \mid[t, L]$ zum Anfangspunkt $\gamma_{e}(t)$, also ist $\dot{\gamma}_{e}(t) \in H_{\gamma_{e}(t)}$ (Axiom 2).

\subsection*{4.13 Satz und Definition}
Sei $E \xrightarrow{\pi} M$ eine lokal triviale Faserung mit Zusammenhang $H \subseteq T E, f$ : $B \rightarrow M$ eine differenzierbare Abbildung. Dann ist auf $f^{*} E$ ein Zusammenhang $f^{*} H$ durch $\left(f^{*} H\right)_{(b, e)}:=\left(d \hat{f}_{(b, e)}\right)^{-1}\left(H_{e}\right)$ definiert. Dabei ist $f^{*} E \xrightarrow{\hat{f}} E$ die kanonische Vektorraumbündelabbildung über $f$.

\section*{Beweis:}
Durch $f^{*} H \subseteq T\left(f^{*} E\right)$ ist ein Untervektorraumbündel definiert, da $f^{*} H$ der Kern der Vektorraumbündelabbildung $T f^{*} E \rightarrow T E \xrightarrow{p r} T^{\prime} E$ ist. Das Faserdifferenzial\\
$d \hat{f}_{(b, e)} \mid T_{(b, e)}^{\prime}\left(f^{*} E\right): T_{(b, e)}^{\prime}\left(f^{*} E\right)=T_{(b, e)}\left(\{b\} \times E_{f(b)}\right) \cong T_{e} E_{f(b)} \rightarrow T_{e} E_{f(b)}=T_{e}^{\prime} E$\\
ist ein Isomorphismus, denn $\hat{f}(b, e)=e$, also $\hat{f} \mid\left(f^{*} E\right)_{b}=\operatorname{pr}_{2}$. Es ist

$$
\left(f^{*} H\right)_{(b, e)}+T_{(b, e)}^{\prime} f^{*} E=T_{(b, e)} f^{*} E
$$

denn ist $v \in T_{(b, e)}\left(f^{*} E\right)$, so ist $d \hat{f}(v) \in T_{f(b)} E=T_{e} E_{f(p)} \oplus H_{e}$, also ist $v \in$ $d \hat{f}^{-1}\left(T_{e} E_{f(p)}\right)+d \hat{f}^{-1}\left(H_{e}\right)$.\\
Die Summe ist direkt: Sei $v \in\left(f^{*} H\right)_{(b, e)} \cap T_{(b, e)}^{\prime}\left(f^{*} E\right)$. Dann ist $d \hat{f}(v)=0$, also ist $v=0$, da $\left.d \hat{f}\right|_{T^{\prime} f^{*} E}$ ein Isomorphismus ist. Also ist $T\left(f^{*} E\right)=\left(f^{*} H\right) \oplus$ $T^{\prime}\left(f^{*} E\right)$.

\subsection*{4.14 Notiz}
Eine differenzierbare Kurve $\alpha: I \rightarrow f^{*} E$ ist genau dann horizontal bezüglich $f^{*} H$, falls $\hat{f} \circ \alpha$ horizontal bezüglich $H$ ist.

\subsection*{4.15 Lemma}
Sei $E \rightarrow M$ eine lokal triviale Faserung mit Zusammenhang $H, \beta:\left(t_{1}, t_{2}\right) \rightarrow M$ eine differenzierbare Kurve, $v$ das auf $\beta^{*} E$ eindeutig definierte horizontale Vektorfeld über $\partial_{t}$. Dann entsprechen die Horizontalkurven über $\beta$ genau den Flusslinien von $v$.\\
Genauer: Ist $\gamma:\left(c_{1}, c_{2}\right) \rightarrow \beta^{*} E$ die eindeutig definierte maximale Integralkurve von $v$ mit $\gamma\left(t_{0}\right)=\left(t_{0}, e\right), e \in E_{\beta\left(t_{0}\right)}$, so ist $\hat{\beta}(\gamma(t))=\operatorname{pr}_{2}(\gamma(t))$ eine Horizontalkurve in $E$ über $\beta$. Umgekehrt ist jede Horizontalkurve von dieser Form.

\section*{Beweis:}
Sei $\gamma$ eine maximale Integralkurve von $v$ mit


\begin{equation*}
\gamma\left(t_{0}\right)=\left(t_{0}, e_{0}\right), e_{0} \in E_{\beta\left(t_{0}\right)} \tag{4.2}
\end{equation*}


Wegen $d \pi(v)=\partial_{t}$ ist $\gamma(t)=(t+c, \tilde{\gamma}(t))$ mit $\tilde{\gamma}(t) \in E_{\beta(t+c)}$. Wegen 4.2) ist also $\gamma(t)=(t, \tilde{\gamma}(t))$. Da $d \hat{\beta}\left(\beta^{*} H\right) \subseteq H$ ist $\gamma$ genau dann horizontal, wenn $\hat{\beta} \circ \gamma$ horizontal ist.\\
Umgekehrt: Jede Kurve $\gamma$ in $E$ über $\beta$ (d.h. mit $\pi \circ \gamma=\beta$ ) definiert eine Kurve in $\beta^{*} E$ durch $\beta^{*} \gamma(t)=(t, \gamma(t))$ die genau dann horizontal bezüglich $f^{*} H$ ist, wenn $\gamma$ horizontal bezüglich $H$ ist.

\subsection*{4.16 Definition}
Ein Zusammenhang $H \subseteq T E$ heißt vollständig, wenn für jede Kurve $\gamma$ : $\left(t_{1}, t_{2}\right) \rightarrow M$ gilt: Jede Horizontalkurve über $\gamma$ ist auf $\left(t_{1}, t_{2}\right)$ definiert.

\subsection*{4.17 Korollar}
\begin{enumerate}
  \item Sei $\beta:\left(t_{1}, t_{2}\right) \rightarrow M, t_{0} \in\left(t_{1}, t_{2}\right)$. Dann gibt es genau eine maximale Horizontalkurve $\beta_{e}:\left(c_{1}, c_{2}\right) \rightarrow E$ zu jedem $e \in E_{\beta\left(t_{0}\right)}$.
  \item Der Zusammenhang einer Parallelverschiebung ist vollständig. Sei nämlich $\gamma:\left(t_{1}, t_{2}\right) \rightarrow M, t_{0} \in\left(t_{1}, t_{2}\right), e \in E_{\gamma\left(t_{0}\right)}$. Setze $\gamma_{e}(t)=\tau_{\gamma\left[t_{0}, t\right]}(e)(*)$ wie in 4.5.
  \item Jeder vollständige Zusammenhang ist ein Zusammenhang einer Parallelverschiebung, definiere nämlich $\tau_{\gamma}$ durch (*). Dann ist 4.3.1 erfüllt nach dem Satz von Picard-Lindelöff, 4.2.2,3 entsprechen den Flussaxiomen, denn für Lösungskurve eines Flusses gilt $\alpha_{x}(t)=\alpha_{\alpha_{x}\left(t_{0}\right)}\left(t-t_{0}\right)$. 4.3.4 (Erstes Ordnungsaxiom) folgt, da $d \pi_{e}: H_{e} \rightarrow T_{\pi(e)} M$ ein Isomorphismus ist.
\end{enumerate}

\subsection*{4.18 Lemma}
Ist $E \rightarrow M$ eine lokal triviale Faserung mit kompakter typischer Faser $F$, so ist jeder Zusammenhang in $E$ vollständig.

\section*{Beweis:}
Sei $E \rightarrow M$ eine lokal triviale Faserung mit Zusammenhang $H, \beta:\left(t_{1}, t_{2}\right) \rightarrow M$ eine Kurve und $v$ das horizontale Vektorfeld über $\partial_{t}$ in $\beta^{*} E$. Sei $\gamma:\left(\tilde{t}_{1}, \tilde{t}_{2}\right) \rightarrow$ $\beta^{*} E$ eine maximale Lösungskurve zu $v$, also o.B.d.A. $\gamma(t)=(t, \tilde{\gamma}(t))$. Wir benutzen, dass Flusslinien endlicher Lebensdauer schließlich jedes Kompaktum verlassen (d.h. ist $\alpha:\left(t_{1}, t_{2}\right) \rightarrow X$ eine Integralkurve eines Vektorfeldes, $t_{2}<$ $\infty$, dann gibt es zu jeder kompakten Teilmenge $K \subseteq X$ ein $T \in\left(t_{1}, t_{2}\right)$ mit $\alpha(t) \notin K$ für $t>T$ ).\\
Angenommen $t_{1}<c_{1} \leq \tilde{t}_{1}<\tilde{t}_{2} \leq c_{2}<t_{2}$. Dann wäre Bild $\gamma \subseteq\left[c_{1}, c_{2}\right] \times F$ (kompakt), also wäre $\gamma$ auf ganz $\mathbb{R}$ definiert. Widerspruch!

\subsection*{4.19 Aufgabe}
Es sei $E \rightarrow M$ eine differenzierbare lokal triviale Faserung mit einem Zusammenhang $H \subset T E$ und $\beta:\left(t_{1}, t_{2}\right) \rightarrow M$ eine differenzierbare Kurve. Man schreibe das Differentialgleichungssystem für die Horizontalkurven $\alpha$ über $\beta$ in lokalen Koordinaten nieder.

\subsection*{4.20 Aufgabe}
Für die triviale Faserung $S^{1} \times \mathbb{R} \rightarrow S^{1}$ gebe man einen nicht vollständigen Zusammenhang an.



\pagebreak

\section{Kovariante Ableitung und Zusammenhangsform}
\subsection*{5.1 Vorbemerkung}
Sei $V$ ein Vektorraum, $V^{\prime} \subseteq V$ ein Untervektorraum. Die Wahl folgender Objekte ist gleichbedeutend:

\begin{enumerate}
  \item Ein Komplement $W \subseteq V$ von $V^{\prime}, V^{\prime} \oplus W=V$
  \item Eine Projektion $V \rightarrow V^{\prime}$
  \item Ein Rechtsinverses von $V \rightarrow V / V^{\prime}$
  \item Eine Spaltung der exakten Sequenz $O \rightarrow V^{\prime} \rightarrow V \rightarrow V / V^{\prime} \rightarrow O$
\end{enumerate}

Entsprechend können Zusammenhänge auf $E \rightarrow M$ auf verschiedene äquivalente Weisen definiert werden.

\subsection*{5.2 Satz und Definition}
Ist $H$ ein Zusammenhang, so heißt die durch $H$ gegebene $T^{\prime} E$-wertige 1-Form $\omega^{\prime} \in \Gamma \operatorname{Hom}\left(T E, T^{\prime} E\right)=\Omega^{1}\left(E, T^{\prime} E\right)$ mit $\omega^{\prime}(v+w)=v$ für $(v, w) \in T_{e}^{\prime} E \oplus H_{e}$ die Zusammenhangsform zu $H$.\\
Umgekehrt: Ist $\omega^{\prime} \in \Omega^{1}\left(E, T^{\prime} E\right)$, sodass $\omega_{e}^{\prime}: E_{e} \rightarrow T_{e}^{\prime} E$ eine Projektion ist, so ist $\omega$ die Zusammenhangsform des Zusammenhangs $H=\operatorname{ker} \omega^{\prime}$.

\subsection*{5.3 Definition}
Ist $H$ ein Zusammenhang auf $E$ mit Zusammenhangsform $\omega^{\prime}$, so ist die kovariante Ableitung (zu H) in Richtung $v \in T_{x} M$ durch

$$
\nabla_{v} s=\omega_{s(x)}^{\prime}(d s(v)) \in T_{e}^{\prime} E
$$

für alle $s \in \Gamma(E \mid U)$ mit $x \in U$ definiert.

\subsection*{5.4 Bemerkung}
Ist $E \rightarrow M$ ein Vektorraumbündel, so ist $T_{e}^{\prime} E=E_{\pi(e)}$, also $\nabla_{v} s \in E_{\pi(v)}$. Insbesondere definiert $\nabla$ dann eine Abbildung

$$
\Gamma T M \times \Gamma E \rightarrow \Gamma E,(v, s) \mapsto\left(\nabla_{v} s\right)
$$

\subsection*{5.5 Definition und Satz}
Sei $R(d \pi) \subseteq \operatorname{Hom}\left(\pi^{*} T M, T E\right)$ das Bündel der Rechtsinversen von $d \pi$, also $R(d \pi)=d \pi^{-1}\left(i d_{\pi^{*} T M}\right)$ also

$$
\alpha_{e} \in R(d \pi)_{e} \Leftrightarrow d \pi_{e} \circ \alpha_{e}=\operatorname{id}_{T_{\pi(e)} M}
$$

$R(d \pi)$ ist ein affines Bündel, also ein Bündel mit typischer Faser $\mathbb{R}^{m n}$ und Strukturgruppe Aff $(m n, \mathbb{R})$, wobei $n$ die Dimension von $E$ und $m$ die Dimension von $M$ ist. Das zugehörige Vektorraumbündel ist $\operatorname{Hom}\left(\pi^{*} T M ; T^{\prime} E\right)$, d.h., jedes $\gamma \operatorname{Hom}\left(\pi^{*} T M, T^{\prime} E\right)$ operiert frei und transitiv auf $\Gamma R(d \pi)$. Die Schnitte in $R(d \pi)$ entsprechen genau den Zusammenhängen $H$ in $E$.

\subsection*{5.6 Definition}
Seien $X$ und $Y$ differenzierbare Mannigfaltigkeiten. Zwei lokal um $x_{0} \in X$ definierte Abbildungen $f$ und $g: U \rightarrow Y$ heißen $k$-äquivalent bei $x_{0}$, wenn gilt:

\begin{enumerate}
  \item $f\left(x_{0}\right)=g\left(x_{0}\right)$
  \item Für eine und dann jede Wahl von Karten um $x_{0}$ und $f\left(x_{0}\right)$ stimmen alle partiellen Ableitungen bis zur Ordnung $k$ von $f$ und $g$ überein.
\end{enumerate}

Man schreibt $[f]_{x_{0}}^{k}$ für die $k$-Äquivalenzklassen bei $x_{0}$. Die Menge dieser Äquivalenzklassen bildet in kanonischer Weise ein differenzierbare Mannigfaltigkeit

$$
J^{k}(X ; Y)=\left\{[f]_{x}^{k} \mid x \in X ; f \text { lokal um } x \text { definiert }\right\}
$$

Insbesondere ist $J^{0}(X ; Y)=X \times Y$. Ist $E \rightarrow M$ eine lokal triviale Faserung, so bezeichnet

$$
J^{k} E:=J^{k}(\pi)=\left\{[f]_{x}^{k} \in J^{k}(M ; E) \mid \pi \circ f=\mathrm{id}\right\}
$$

\subsection*{5.7 Bemerkung}
Für jedes $v \in T_{\pi(e)} M$ sei $E \xrightarrow{E} M$ eine lokal triviale Faserung mit Zusammenhang $H$. Dann ist

$$
\nabla_{v}: J_{e}^{1} E \rightarrow T_{e}^{\prime} E,[s]_{x}^{1} \mapsto \nabla_{v} s
$$

eine wohldefinierte Abbildung.

\subsection*{5.8 Lemma}
\begin{enumerate}
  \item Kanonisch ist $J^{0}(\pi)=E,[s]_{x}^{0} \mapsto s(x)$
  \item $J^{1}(\pi) \rightarrow R(d \pi),[s]_{x}^{1} \mapsto d s_{x}$ ist ein wohldefinierter Diffeomorphismus über $E$ der $\operatorname{Hom}\left(\pi^{*} T M, T^{\prime} E\right)$-äquivariant ist. Damit ist $J^{1}(E) \rightarrow E$ mit $[s]_{x}^{1} \mapsto s(x)$ ein affnes Bündel über $E$ mit Vektorraumbündel Hom $\left(\pi^{*} T M ; T^{\prime} E\right)$. Die Aktion Hom $\left(\pi^{*} T M, T^{\prime} E\right) \times J^{1}(\pi) \rightarrow J^{1}(\pi)$ ist dabei durch $(\alpha,[s]) \mapsto$ $[s+\alpha]$ ) bezüglich Bündelkarten gegeben. Ein Zusammenhang kann damit gelesen werden als Schnitt in $J^{1}(\pi) \rightarrow E$.
\end{enumerate}

\section*{Beweis:}
Seien $[s]$ und $[\tilde{s}] \in J^{1}(\pi)$ mit $s(x)=\tilde{s}(x)$. Dann gilt $[s]_{x}^{1}=[\tilde{s}]_{x}^{1} \Leftrightarrow d s_{x}=d \tilde{s}_{x}$. Folglich ist die Abbildung wohldefiniert und injektiv. Um die Surjektivität zu zeigen, benutzen wir lokale Karten (Details: Übungsaufgabe).

Zusammenfassend ist ein Zusammenhang also eine Spaltung der Sequenz

$$
0 \rightarrow T^{\prime} E \xrightarrow{i} T E \xrightarrow{d \pi} \pi^{*} T M \rightarrow 0
$$

\subsection*{5.9 Lemma}
Sei $\nabla$ die kovariante Ableitung eines Zusammenhangs auf $E \rightarrow M$, dann ist

$$
\nabla: J^{1} E \rightarrow \operatorname{Hom}\left(\pi^{*} T M ; T^{\prime} E\right)
$$

eine Vektorisierung des affinen Bündels $J^{1} E \rightarrow E$, d.h. ein translationäquivarianter Diffeomorphismus über $E$. Der Zusammenhang ist dann durch $\nabla^{-1}(0)$ definiert.

\section*{Beweis:}
Sei $\sigma_{e}: T_{x} M \rightarrow T_{e} E$ das durch den Zusammenhang definierte Rechtsinverse von $d \pi_{e}$ (also $\sigma_{e}\left(T_{x} M\right)=H_{e}$ ). Dann ist

$$
\begin{aligned}
J^{1}(\pi)_{e} & =R(d \pi)_{e} \xrightarrow{\nabla} \operatorname{Hom}\left(T_{x} M, T_{e}^{\prime} E\right) \\
{[s]_{x}^{1} } & =:[\sigma+\varphi] \mapsto \nabla(\sigma+\varphi)=\varphi
\end{aligned}
$$

und für $\psi \in \operatorname{Hom}\left(T_{x} M ; T_{e}^{\prime} E\right)$ ist $[s+\psi]=[\sigma+\varphi+\psi]$ also $\nabla(s+\psi)=\varphi+\psi=$ $\nabla s+\psi$.

\subsection*{5.10 Korollar}
$\nabla[s]=[s]-\sigma(s(x))$, wobei $\sigma$ wie im Beweis von 5.9.

\subsection*{5.11 Definition}
Unter einer kovarianten Ableitung auf einer lokal trivialen Faserung versteht man eine Vektorisierung des affinen Bündels $J^{1} E \rightarrow E$.

\subsection*{5.12 Korollar}
Durch $\nabla \rightarrow \nabla^{-1}(0)$ ist eine Bijektion zwischen kovarianten Ableitungen und dem Raum der Zusammenhänge (gelesen als Schnitt in $J^{1}(\pi)$ ) gegeben.

\subsection*{5.13 Definition}
Ist $\nabla$ eine kovariante Ableitung in $E$ und $\omega^{\prime}$ die zugehörige Zusammenhangsform, so ist für $\alpha: I \rightarrow E$

$$
\frac{\nabla}{d t} \alpha:=\omega^{\prime}(\dot{\alpha}(t))
$$

die kovariante Ableitung von $\alpha$.

\subsection*{5.14 Notiz}
Ist $v \in T_{x} M$ und $\beta$ eine repräsentierende Kurve, also $\beta(0)=v$, dann ist $\frac{\nabla}{d t} s \circ \beta=\nabla_{v} s$, denn $d s_{x}(v)=\left.\frac{d}{d t}\right|_{t=0} s \circ \beta$.

\subsection*{5.15 Notation}
Sei $\tau$ ein Paralleltransport eines Zusammenhangs, und sind $\beta:\left(t_{1}, t_{2}\right) \rightarrow M$ und $\alpha:\left(t_{1}, t_{2}\right) \rightarrow E$ Kurven mit $\pi \circ \alpha=\beta$. Sei $t_{0} \in\left(t_{1}, t_{2}\right)$, dann nennen wir

$$
\alpha_{\left(t_{0}\right)}:\left(t_{1}, t_{2}\right) \rightarrow E_{\beta\left(t_{0}\right)}, \alpha_{\left(t_{0}\right)}(t):=\tau_{\beta_{\left[t, t_{0}\right]}}(\alpha(t))
$$

$\operatorname{mit} \beta_{\left[t, t_{0}\right]}(s)=\beta\left(s t_{0}+(1-s) t\right)$ den $t_{0}$-Monitor von $\alpha$.

\subsection*{5.16 Lemma}
$$
\left.\frac{\nabla}{d t}\right|_{t=t_{0}} \alpha=\dot{\alpha}_{\left(t_{o}\right)}\left(t_{0}\right)
$$

\section*{Beweis:}
Sei $f:\left(t_{1}, t_{2}\right) \times E_{\beta\left(t_{0}\right)} \rightarrow E,(t, e) \mapsto \tau_{\beta_{\left[t_{0}, t\right]}}(e)$. Dann gilt

\begin{enumerate}
  \item $f\left(t, \alpha_{\left(t_{0}\right)}(t)\right)=\alpha(t)$
  \item $f \mid\left\{t_{0}\right\} \times E_{\beta\left(t_{0}\right)}=\mathrm{id}$
  \item $f$ führt horizontale Kurven im Produkt in horizontale Kurven über.
\end{enumerate}

Also ist $d f_{(t, e)}\left(1, \dot{\alpha}_{\left(t_{0}\right)}\left(t_{0}\right)\right)=d f\left((1,0)+\left(0, \dot{\alpha}_{\left(t_{0}\right)}\left(t_{0}\right)\right)=v+\dot{\alpha}_{\left(t_{0}\right)}\left(t_{0}\right)\right.$, wobei $v$ horizontal ist, also ist $\left.\frac{\nabla}{d t}\right|_{t=t_{0}} \alpha(t)=\dot{\alpha}_{\left(t_{0}\right)}\left(t_{0}\right)$.

\subsection*{5.17 Notation}
Wir bezeichnen den affinen Raum der Zusammenhänge über $E \rightarrow M$ mit $\mathcal{C}(E)=\Gamma R(d \pi)=\Gamma J^{1} E$.

\subsection*{5.18 Aufgabe}
Es sei $G$ eine Liegruppe mit einer gegenüber (Rechts- und Links-) Translation invarianten Riemannschen Metrik und $G_{0} \subset G$ eine kompakte Untergruppe. Man zeige: Die Parallelverschiebung des zu den Fasern orthogonalen Zusammenhangs von $G \rightarrow G / G_{0}$ ist isometrisch.

\subsection*{5.19 Aufgabe}
Eine Liegruppe sei durch Links- oder Rechtstranslation parallelisiert. Man zeige, dass bezüglich des dadurch kanonisch gegebenen Zusammenhangs für $T G \rightarrow G$ die einparametrigen Untergruppen geodätisch sind.



\pagebreak

\section{G-Zusammenhänge}

\subsection*{6.1 Definition}
Sei $E \rightarrow M$ ein Faserbündel mit Strukturgruppe G. Dann heißt eine Parallelverschiebung in $E$ eine $G$-Parallelverschiebung, falls sie bezüglich Karten durch eine $G$-Linksaktion gegeben ist.

\subsection*{6.2 Beispiel}
\begin{enumerate}
  \item Ist $P \rightarrow M$ ein Prinzipalbündel, dann ist eine Parallelverschiebung genau dann eine $G$-Parallelverschiebung, wenn für alle $\gamma$ gilt: $\tau_{\gamma}(p g)=\tau_{\gamma}(p) g$
  \item Ist $E \rightarrow M$ ein Vektorraumbündel, so ist eine Parallelverschiebung genau dann eine $G L(n, \mathbb{R})$-Parallelverschiebung, falls $\tau_{\gamma}$ für alle $\gamma$ linear ist.
  \item Ist ( $M, g$ ) eine Riemannsche Mannigfaltigkeit, so ist eine Parallelverschiebung auf $T M$ genau dann eine $O(n)$-Parallelverschiebung, falls $\tau_{\gamma}$ eine Isometrie ist für jedes $\gamma$.
\end{enumerate}

\subsection*{6.3 Satz}
Ist $\tau$ eine $G$-Parallelverschiebung auf $P$ durch

$$
\hat{\tau}_{\gamma}: P_{x} \times_{G} F \rightarrow P_{y} \times_{G} F,[p, v] \mapsto\left[\tau_{\gamma}(p), v\right], \quad \gamma(0)=x ; \gamma(1)=y
$$

eine $G$-Parallelverschiebung auf $P \times_{G} V$ gegeben.\\
Ist $\tau$ eine $G$-Parallelverschiebung auf $E$, so ist durch

$$
\tilde{\tau}_{\gamma}: \operatorname{Iso}_{\gamma}\left(F, E_{x}\right) \rightarrow \operatorname{Iso}_{G}\left(F, E_{y}\right), \varphi \mapsto \tau_{\gamma} \circ \varphi
$$

eine Parallelverschiebung auf $\operatorname{Iso}_{G}(F, E)$ gegeben und es gilt $\tilde{\hat{\tau}}=\tau, \hat{\tilde{\tau}}=\tau$. (Beweis: Übung).

Erinnere: Ist $\tau$ eine Parallelverschiebung, so ist der zugehörige infinitesimale Parallelismus durch

$$
H_{e}=\left\{\dot{\gamma}_{e}(0) \mid \gamma_{e}(t)=\tau_{\gamma_{[0 ; t]}}(e), \gamma:(-\varepsilon, \varepsilon) \rightarrow M, \gamma(0)=\pi(e)\right\}
$$

definiert.

\subsection*{6.4 Satz}
Ein infinitesimaler Parallelismus auf einem $G$-Prinzipalbündel kommt genau dann von einer $G$-Parallelverschiebung, falls gilt $H_{p g}=\left.d R_{g}\right|_{p} H_{p}(*)$.

\section*{Beweis:}
Für die Horizontalkurven gilt

$$
\gamma_{p g}(t)=\gamma_{p}(t) \cdot g
$$

(vgl. 6.2.1), also

$$
\dot{\gamma}_{p g}(0)=\left.d R_{g}\right|_{p} \dot{\gamma}_{p}(0)
$$

Andererseits: Gilt (*), so ist für horizontales $\gamma_{p}$ auch $\gamma_{p g}$ horizontal, also $\tau_{\gamma}(p g)=\tau_{\gamma}(p) g$.

\subsection*{6.5 Definition}
Sei $P \rightarrow M$ ein $G$-Prinzipalbündel, so heißt ein Zusammenhang $H \subseteq T P$ ein $G$-Zusammenhang, falls für alle $p \in P$ und $g \in G$ gilt $d R_{g} H_{p}=H_{p g}$.

\subsection*{6.6 Satz}
Jeder $G$-Zusammenhang auf einem $G$-Prinzipalbündel $P \rightarrow M$ ist vollständig, d.h., er ist ein infinitesimaler Parallelismus einer $G$-Parallelverschiebung.

\section*{Beweis:}
Sei $\beta:\left(t_{1}, t_{2}\right) \rightarrow M$ eine differenzierbare Kurve. Es ist zu zeigen: Jede Horizontalkurve über $\beta$ ist auf ganz ( $t_{1}, t_{2}$ ) definiert.\\
Die Horizontalkurven über $\beta$ entsprechen genau den Lösungskurven des Horizontalvektorfeldes über $\partial_{t} \in \beta^{*} P \cong\left(t_{1}, t_{2}\right) \times G$.\\
Genauer: Ist $\gamma_{\left(t_{0}, p\right)}(t)=:\left(t+t_{0}, \hat{\gamma}(t)\right)$ eine Lösungskurve des horizontalen Vektorfeldes mit $\gamma_{\left(t_{0}, p\right)}(0)=\left(t_{0}, p\right)$ für $p \in P_{\beta\left(t_{0}\right)}$, so ist $\hat{\gamma}\left(t-t_{0}\right)$ eine Horizontalkurve über $\beta$ und umgekehrt. Sei ( $a_{\left(t_{0}, p\right)}, b_{\left(t_{0}, p\right)}$ ) der maximale Definitionsbereich von $\gamma_{\left(t_{0}, p\right)}$. Dann ist $\left(a_{\left(t_{0}, p g\right)}, b_{\left(t_{0}, p g\right)}\right)=\left(a_{\left(t_{0}, p\right)}, b_{\left(t_{0}, p\right)}\right)$ für alle $g \in G$, also hängt $a_{\left(t_{0}, p\right)}$ und $b_{\left(t_{0}, p\right)}$ nicht von $p$ ab. Ist $\left[T_{1}, T_{2}\right] \subseteq\left(t_{1}, t_{2}\right)$, so gibt es also ein $\varepsilon>0$, sodass $b_{(t, p)}>\varepsilon$ für alle $t \in\left[T_{1}, T_{2}\right]$ und $a_{(t, p)}<-\varepsilon$ für alle $t \in\left[T_{1}, T_{2}\right]$ und alle $p \in P$.\\
Angenommen, der maximale Definitionsbereich von $\gamma_{\left(t_{0}, p\right)}$ ist $\left(T_{1}, T_{2}\right)$ mit $T_{2}<$ $t_{2}$. Dann wäre $b_{\left(T_{2}-\frac{\varepsilon}{2}, p\right)}=\frac{\varepsilon}{2}$ für alle $p \in P_{\beta\left(T_{2}-\frac{\varepsilon}{2}\right)}$.

\subsection*{6.7 Definition und Satz}
Ist $H$ ein Zusammenhang auf einem Faserbündel mit Strukturgruppe $G$, so heißt $H$ ein $G$-Zusammenhang, falls er infinitesimaler Parallelismus einer $G$ Parallelverschiebung ist. Ist $E=P \times_{G} F$ mit einem $G$-Zusammenhang $H \subseteq$ $T E$ und $\tilde{H} \subseteq T P$ der zugehörige $G$-Zusammenhang in $P$, so ist $H_{[p, v]}=$ $d f_{v}\left(\tilde{H}_{p}\right)$, wobei $f_{v}: P \rightarrow P \times_{G} F, p \mapsto[p, v]$.\\
Wir schreiben $C^{G}(E)$ für den Raum der $G$-Zusammenhänge auf $E$.

\subsection*{6.8 Bemerkung}
Ist $P$ ein $G$-Prinzipalfaserbündel und $L_{p}: G \rightarrow P, g \mapsto p g$, so ist

$$
T^{\prime} P=P \times \mathfrak{g},\left.(p, X) \mapsto d L_{p}\right|_{1}(X)
$$

ein Vektorraumbündelisomorphismus.\\
Das Vektorraumbündel $\pi: T^{\prime} P \rightarrow P$ ist kanonisch ein Rechts- $G$-Bündel, denn ist $R_{g}: P \rightarrow P, p \mapsto p g$ die Rechtsmultiplikation, so ist

$$
d R_{g}: T_{p}^{\prime} P \rightarrow T_{p g}^{\prime} P
$$

wohldefiniert und $\pi\left(d R_{g}(v)\right)=\pi(v) g$.\\
Damit wird auch $P \times \mathfrak{g} \rightarrow P$ zu einem Rechts-Vektorraumbündel mit der Rechts-G-Aktion

$$
G \times(P \times \mathfrak{g}) \rightarrow(P \times \mathfrak{g}),(g,(p, X)) \mapsto\left(p g, \operatorname{Ad}\left(g^{-1}\right) X\right),
$$

denn $\left.d L_{p g}\right|_{1}\left(\operatorname{Ad}\left(g^{-1}\right) X\right)=\left.\left.d R_{g}\right|_{p} \circ d L_{p}\right|_{1}(X)$, da

$$
L_{p g} \circ \operatorname{konj}\left(g^{-1}\right)\left(\gamma(t)=p g g^{-1} \gamma(t) g=p \gamma(t) g=R_{g} \circ L_{p}(\gamma(t)) .\right.
$$

Insbesondere ist also für $v \in T_{p} P$


\begin{equation*}
d R_{g}(v)=\left.d L_{p g}\right|_{1}\left(\operatorname{Ad}\left(g^{-1}\right)\left(d L_{p}\right)^{-1}(v) .\right. \tag{*}
\end{equation*}


\subsection*{6.9 Definition}
Ist $\omega^{\prime}$ eine Zusammenhangsform eines $G$-Zusammenhangs $H$ auf einem Prinzipalbündel $P \rightarrow M$, so heißt die durch

$$
\left.d L_{p}\right|_{1} ^{-1} \circ \omega_{p}^{\prime}=: \omega_{p}
$$

definierte 1-Form $\omega \in \Omega^{1}(P, \mathfrak{g})$ die ( $G$-)Zusammenhangsform des $G$-Zusammenhangs.

\subsection*{6.10 Satz}
Eine 1-Form $\omega \in \Omega^{1}(P, \mathfrak{g})$ ist genau dann eine $G$-Zusammenhangsform eines $G$-Zusammenhangs, falls gilt:

\begin{enumerate}
  \item $\omega \mid T_{p}^{\prime} P=d L_{p}^{-1}$
  \item $R_{g}^{*} \omega=\operatorname{Ad}\left(g^{-1}\right) \omega$
\end{enumerate}

\section*{Beweis:}
Sei $\omega \in \Omega^{1}(P, \mathfrak{g})$ eine $G$-Zusammenhangsform, dann gilt 1) offenbar und ist $v=v_{H}+v^{\prime}$ mit $v_{H} \in H_{p}, v^{\prime} \in T_{p}^{\prime} P$, so ist $d R_{g}\left(v_{H}\right)$ horizontal und $d R_{g}\left(v^{\prime}\right)$ vertikal, also folgt mit (*):

$$
\begin{gathered}
\left(R_{g}^{*} \omega\right)_{p}(v)=\omega_{p g}\left(d R_{g}\left(v_{H}\right)+d R_{g}\left(v^{\prime}\right)\right)=\omega_{p g}\left(d R_{g}\left(v^{\prime}\right)=d L_{p g}^{-1} \circ d R_{g}\left(v^{\prime}\right)\right. \\
=\operatorname{Ad}\left(g^{-1}\right) \circ\left(d L_{p}\right)^{-1}\left(v^{\prime}\right)=\operatorname{Ad}\left(g^{-1}\right) \omega_{p}(v)
\end{gathered}
$$

Umgekehrt: Erfüllt $\omega$ die Bedingungen 1) und 2), so definiert $H=\operatorname{ker} \omega$ einen $G$-Zusammenhang, denn es ist $H \cap T^{\prime} P=\{0\}$ und $d R_{g}\left(H_{p}\right)=H_{p g}$.

\subsection*{6.11 Definition}
Das Vektorraumbündel $P \times_{\text {Ad }} \mathfrak{g}=: \hat{\mathfrak{g}} \rightarrow M$ heißt das Bündel der infinitesimalen Eichtransformationen.

\subsection*{6.12 Bemerkung}
Ist $\sigma$ ein (lokaler) Schnitt in $P \times_{\text {Ad }} \mathfrak{g}, \sigma(x)=[s(x), v(x)]$, so ist durch

$$
(\exp \sigma)(x)=[s(x), \exp (v(x))]
$$

ein (lokaler) Schnitt im Bündel der Eichtransformationen $P \times_{\text {konj }} G$ wohldefiniert.

\subsection*{6.13 Definition}
Ist $P \rightarrow M$ ein $G$-Prinzipalbündel und $\rho$ eine Darstellung von $G$ auf $V$. Dann heißt $\alpha \in \Omega^{k}(P, V)$

\begin{enumerate}
  \item horizontal, $\alpha \in \Omega_{\mathrm{hor}}^{k}(P, V)$, falls für $v \in T^{\prime} P$ gilt: $i_{v} \alpha=0$,
  \item $\rho$-invariant, $\alpha \in \Omega_{\mathrm{inv}}^{k}(P, V)$, falls gilt $R_{g}^{*} \alpha=\rho\left(g^{-1}\right) \alpha$.
\end{enumerate}

Ist $\omega \in \Omega_{\mathrm{inv}, \mathrm{hor}}^{k}(P, V):=\Omega_{\mathrm{inv}}^{k}(P, V) \cap \Omega_{\mathrm{hor}}^{k}(P, V)$, so heißt $\omega$ auch tensoriell.

\subsection*{6.14 Satz}
Es gilt: $\Omega_{\text {inv.hor }}^{k}(P, V)=\Omega^{k}\left(M, P \times{ }_{\rho} V\right)$.

\section*{Beweis:}
Seien $\alpha \in \Omega_{\text {inv }, \text { hor }}^{k}(P, V)$ und $v_{1}, \ldots, v_{k} \in T_{x} M$. Dann ist $\tilde{\alpha} \in \Omega^{k}\left(M, P \times_{\rho} V\right)$ $\operatorname{durch} \tilde{\alpha}_{x}\left(v_{1}, \ldots, v_{k}\right)=\left[p, \alpha_{p}\left(\tilde{v}_{1}, \ldots, \tilde{v}_{k}\right)\right]$, wobei $\tilde{v}_{j} \in T_{p} P$ mit $d \pi\left(\tilde{v}_{j}\right)=v_{j}$ wohldefiniert.\\
Umgekehrt sind $\alpha \in \Omega^{k}\left(M, P \times_{\rho} V\right)$ und $v_{1}, \ldots, v_{k} \in T_{p} P$, so ist $\hat{\alpha} \in \Omega_{\text {inv }, \text { hor }}^{k}(P, V)$ durch $\alpha_{x}\left(d \pi\left(v_{1}\right), \ldots, d \pi\left(v_{k}\right)\right)=:\left[p, \hat{\alpha}_{p}\left(v_{1}, \ldots, v_{k}\right)\right]$ wohldefiniert.\\
Mithilfe eines Zusammenhangs können wir jetzt eine Horizontalableitung definieren.

\subsection*{6.15 Definition}
Sei $H$ ein $G$-Zusammenhang auf $P$. Dann ist die Horizontalableitung auf $P$ durch

$$
D^{H}: \Omega^{k}(P, V) \rightarrow \Omega_{\mathrm{hor}}^{k+1}(P, V), \omega \mapsto\left(\left(v_{1}, \ldots, v_{k+1}\right) \mapsto(d \omega)\left(\bar{v}_{1}, \ldots, \bar{v}_{k+1}\right)\right)
$$

wobei $\bar{v}_{j}$ der Horizontalanteil von $v_{j}$ ist, definiert.

\subsection*{6.16 Definition und Bemerkung}
Das kovariante Differential $d^{H}: \Omega^{k}\left(M, P \times{ }_{\rho} V\right) \rightarrow \Omega^{k+1}\left(M, P \times{ }_{\rho} V\right)$ ist durch $D^{H}$ und die Identifizierung $\Omega\left(M, P \times{ }_{\rho} V\right)=\Omega_{\text {inv,hor }}(P, V)$ wohldefiniert.

\section*{Beweis:}
Es genügt zu zeigen: Für $\eta \in \Omega_{\mathrm{inv}, \mathrm{hor}}^{k}(P, V)$ ist $D^{H} \eta \in \Omega_{\mathrm{inv}}^{k+1}(P, V)$. Sei $\bar{v}_{j}$ der Horizontalanteil von $v_{j}$. Dann ist

$$
\begin{aligned}
\left(R_{g}^{*} D^{H} \eta_{p}\right)\left(v_{1}, \ldots, v_{k+1}\right) & =\left(D^{H} \eta\right)_{p g}\left(d R_{g}\left(v_{1}\right), \ldots, d R g\left(v_{k+1}\right)\right) \\
& =(d \eta)_{p g}\left(d R_{g}\left(\bar{v}_{1}\right), \ldots, d R_{g}\left(\bar{v}_{k+1}\right)\right) \\
& =d\left(R_{g}^{*} \eta\right)_{p}\left(\bar{v}_{1}, \ldots, \bar{v}_{k+1}\right)=d\left(\rho\left(g^{-1}\right) \eta\right)_{p}\left(\bar{v}_{1}, \ldots, \bar{v}_{k+1}\right) \\
& =\rho\left(g^{-1}\right) D^{H} \eta\left(v_{1}, \ldots, v_{k+1}\right)
\end{aligned}
$$

\subsection*{6.17 Korollar}
Ist $\hat{\eta} \in \Omega_{\mathrm{inv}, \text { hor }}^{k}(P, V)$ die durch $\eta \in \Omega^{k}\left(M, P \times{ }_{\rho} V\right)$ gegebene Differenzialform, so ist

$$
\left(d^{H} \eta\right)_{x}\left(v_{1}, \ldots, v_{k}\right)=\left[p, d \hat{\eta}\left(\bar{v}_{0}, \ldots, \bar{v}_{k}\right)\right]=\left[p, D^{H} \hat{\eta}\left(\tilde{v}_{0}, \ldots, \tilde{v}_{k}\right)\right]
$$

wobei $\bar{v}_{j} \in T_{p} P$ horizontal ist mit $d \pi\left(\bar{v}_{j}\right)=v_{j}$ und $\tilde{v}_{j} \in T_{p} P$ mit $d \pi\left(\tilde{v}_{j}\right)=v_{j}$.

\subsection*{6.18 Satz}
Ist $\omega \in \Omega_{\text {inv }}^{1}(P, \mathfrak{g})$ die Zusammenhangsform eines $G$-Zusammenhangs in $P$, $\eta \in \Omega_{\text {inv }, \text { hor }}^{k}(P, V)$, so ist

$$
D^{H} \eta=d \eta+\rho_{*} \omega \wedge \eta
$$

wobei

$$
\left(\rho_{*} \omega \wedge \eta\right)\left(v_{0}, \ldots, v_{k}\right)=\left.\sum_{i=0}^{k}(-1)^{i} d \rho\right|_{1} \omega\left(v_{i}\right) \cdot \eta\left(v_{0}, \ldots, \hat{v}_{i}, \ldots, v_{k}\right)
$$

\section*{Beweis:}
\begin{enumerate}
  \item Fall: $\left(v_{0}, \ldots, v_{k}\right)$ sind alle horizontal. Dann ist $\left(\rho_{*} \omega \wedge \eta\right)\left(v_{0}, \ldots, v_{k}\right)=0$.
  \item Fall: Mindestens zwei Vektoren sind vertikal. Dann verschwindet die linke Seite und der zweite Summand auf der rechten Seite.\\
Es ist
\end{enumerate}

$$
\begin{aligned}
d \eta\left(v_{0}, \ldots, v_{k}\right)= & \sum_{i=0}^{k}(-1)^{i} v_{i}\left(\eta\left(v_{0}, \ldots, \hat{v}_{i}, \ldots, v_{k}\right)\right)+ \\
& +\sum_{i<j}(-1)^{i+j} \eta\left(\left[V_{i}, V_{j}\right], V_{0}, \ldots, \hat{V}_{i}, \ldots, \hat{V}_{j}, \ldots V_{k}\right)
\end{aligned}
$$

wobei auf der rechten Seit die Vektoren $v_{i}$ zu Vektorfeldern $V_{i}$ ergänzt sind und zwar so, dass die vertikalen Vektoren zu vertikalen Vektorfeldern ergänzt sind. Dann ist $\left[V_{i}, V_{j}\right]$ vertikal, falls $V_{i}, V_{j}$ vertikal sind. Also ist in diesem Fall $d \eta\left(v_{0}, \ldots, v_{k}\right)=0$.\\
3. Fall: Genau ein Vektor ist vertikal, alle anderen horizontal. Sei oBdA $v_{0}$ vertikal und alle anderen Vektoren horizontal. Die linke Seite verschwindet. Der zweite Summand auf der rechten Seite ergibt: $d \rho_{1}\left(\omega\left(v_{0}\right)\right) \eta\left(v_{1}, \ldots, v_{k}\right)$. Sei $v_{0}=\left.d L_{p}\right|_{1}(X), X \in \mathfrak{g}$. Wir setzen $v_{0}$ auf die Faser zu $X_{G}(\tilde{p})=$ $d L_{\tilde{p}}(X)$ und $v_{i}$ als horizontale Vektorfelder $v_{i}$ fort. Dann ist $\left[X_{G}, v_{i}\right]=0$, also $d \eta\left(v_{0}, \ldots, v_{k}\right)=X_{G} \eta\left(v_{1}, \ldots, v_{k}\right)$.\\
Zu zeigen bleibt: $\left.d \rho\right|_{1}(X) \eta\left(v_{1}, \ldots, v_{k}\right)=-X_{G} \eta\left(v_{1}, \ldots, v_{k}\right)$.\\
Es ist

$$
\begin{aligned}
X_{G} \eta\left(v_{1}, \ldots, v_{k}\right) & =\left.\frac{d}{d t}\right|_{t=0} \eta_{p \exp (t X)}\left(v_{1}(p \exp t X), \ldots, v_{k}(p \exp t X)\right) \\
& \left.\stackrel{\text { Inv. }}{=} \frac{d}{d t}\right|_{t=0} \rho(\exp (-t X)) \eta_{p}\left(v_{1}(p), \ldots, v_{k}(p)\right)
\end{aligned}
$$

\subsection*{6.19 Bemerkung}
Für die triviale Darstellung $\rho$ ist $P \times{ }_{\rho} V=M \times V,[p, v] \mapsto(\pi(p), v)$, also $\Omega^{k}\left(M, P \times{ }_{\rho} V\right)=\Omega^{k}(M, V)$ und $d^{H} \eta=d \eta$.

Was das kovariante Differential mit der kovarianten Ableitung zu tun hat, werden wir im nächsten Abschnitt sehen. Wir wenden uns zuerst noch einmal den Zusammenhangsformen in $P$ zu.

\subsection*{6.20 Korollar}
Sind $\omega$ und $\tilde{\omega}$ zwei $G$-Zusammenhangsformen zweier $G$-Zusammenhänge $H$ und $\tilde{H}$ auf $P$, dann ist $\omega-\tilde{\omega} \in \Omega^{1}\left(M, P \times_{\text {Ad }} \mathfrak{g}\right)=\Gamma\left(\operatorname{Hom}\left(T M, P \times_{\text {Ad }} \mathfrak{g}\right)\right)$ und umgekehrt: Ist $A \in \Omega^{1}\left(M, P \times_{\text {Ad }} \mathfrak{g}\right)$ und $\omega$ eine $G$-Zusammenhangsform, so ist auch $\omega+A$ eine Zusammenhangsform auf $P$.\\
Der Raum der $G$-Zusammenhänge auf $P$ ist also ein affner Raum mit Vektorraum $\Omega^{1}\left(M, P \times_{\text {Ad }} \mathfrak{g}\right)$.

\section*{Beweis:}
Nach Satz 6.14 und 6.10 ist $\omega-\tilde{\omega} \in \Omega_{\text {inv,hor }}^{1}(P, \mathfrak{g})=\Omega^{1}(M, \hat{\mathfrak{g}})$\\
Umgekehrt: Ist $\alpha \in \Omega^{1}\left(M, P \times_{\text {Ad }} \mathfrak{g}\right)$, so ist durch $(\omega+\alpha)_{p}(v):=\omega_{p}(v)+$ $d L_{p} \hat{\alpha}_{p}(v)$ für $v \in T_{p} P$ mit $\alpha\left(d \pi_{p}(v)\right)=:\left[p, \hat{\alpha}_{p}(v)\right]$ eine $G$-Zusammenhangsform wohldefiniert und es gilt:

\begin{enumerate}
  \item $(\omega+\alpha)(v)=\omega(v)=d L_{p}^{-1} v$ für $v \in T_{p}^{\prime} P$
  \item $\hat{\alpha}_{p g}\left(d R_{g}(v)\right)=\operatorname{Ad}\left(g^{-1}\right) \hat{\alpha}_{p}(v)$, also $R_{g}^{*}(\omega+\alpha)=\operatorname{Ad}\left(g^{-1}\right)(\omega+\alpha)$
\end{enumerate}

\subsection*{6.21 Beispiel}
\begin{enumerate}
  \item Sei $M=G / K$ ein reduktiver homogener Raum, also $\mathfrak{g}=\mathfrak{k} \oplus \mathfrak{m}$ mit $\operatorname{Ad}(K) \mathfrak{m} \subseteq \mathfrak{m}$ und $T M=G \times_{\operatorname{Ad}(\mathrm{K})} \mathfrak{m}$, dann ist in dem $K$-Prinzipalbündel $G \rightarrow G / K$ ein $K$-Zusammenhnag $H$ durch
\end{enumerate}

$$
H_{g}=d L_{g}(\mathfrak{m})
$$

wohldefiniert.\\
2. Wir betrachten $S U(2)=S^{3} \rightarrow S^{2},\left(z_{1}, z_{2}\right) \mapsto\left[z_{1}: z_{2}\right]$. Dies ist ein $S^{1}$-Prinzipalbündel mit $S^{1}$-Rechtsaktion

$$
S^{3} \times S^{1} \rightarrow S^{3},\left(\left(z_{1}, z_{2}\right), e^{i \theta}\right) \mapsto\left(z_{1} e^{i \theta}, z_{2} e^{i \theta}\right)
$$

Wir identifizieren $\mathbb{R}^{4} \cong \mathbb{C}^{2}$, also $i\left(p_{1}, p_{2}, p_{3}, p_{4}\right)^{T}=\left(-p_{2}, p_{1},-p_{4}, p_{3}\right)^{T}$ und definieren $\omega \in \Omega^{1}\left(S^{3}, i \mathbb{R}\right)$ durch

$$
\omega_{p}(Y)=i<Y, i p>\in i \mathbb{R}
$$

Behauptung: Dies ist eine Zusammenhangsform eines $S^{1}$-Zusammenhangs, denn\\
(a) Für $Y \in T_{p}^{\prime} S^{3}$ ist $\omega_{p}(y)=d L_{p}^{-1}(y)$.\\
(b) $R_{g}^{*} \omega=\operatorname{Ad}\left(g^{-1}\right) \omega$ für alle $g \in S^{1}$.

\section*{Beweis:}
(a) Es ist $T_{p}^{\prime} S^{3}=\mathbb{R} i p$ und $\omega_{p}(i p)=i<i p, i p>=i=d L_{p}^{-1}(i p)$.\\
(b) $R_{e^{i t}}^{*} \omega_{p}(Y)=\omega_{p e^{i t}}\left(d R_{e^{i t}} Y\right)=i<e^{i t} Y, i e^{i t} p>=i<Y$, ip >= $\left.\omega_{p}(Y)=\operatorname{Ad}\left(e^{i t}\right) \omega_{p}(Y)\right)$.\\
Also ist $\omega$ eine Zusammenhangsform eines $S^{1}$-Zusammenhangs. Es ist

$$
H_{e_{1}}=\left\{y \in 0 \times \mathbb{R}^{3} \mid<y, e_{2}>=0\right\}=0 \times R^{2}
$$

\subsection*{6.22 Definition}
Sei $\omega \in \Omega^{1}(P, \mathfrak{g})$ die $G$-Zusammenhangsform eines $G$-Zusammenhangs, $s \in$ $\Gamma(P \mid U)$, so heißt $s^{*} \omega$ Beschreibung des Zusammenhangs mittels $s$ oder lokale Zusammenhangsform bezüglich $s$.

\subsection*{6.23 Lemma}
Sind $s, \tilde{s} \in \Gamma(P \mid U)$ lokale Schnitte, $\tilde{s}(x)=s(x) g(x)$ für $g: U \rightarrow G$, dann ist

$$
\tilde{s}^{*} \omega=\operatorname{Ad}\left(g^{-1}\right) s^{*} \omega+g^{-1} d g
$$

\section*{Beweis:}
Sei $v \in T_{x} M$ und $\gamma$ repräsentierende Kurve, also $\gamma(0)=x, \dot{\gamma}(0)=v$. Sei $g(\gamma(0))=g_{0}$, also $\tilde{s}(x)=s(x) g_{0}$. Dann ist

$$
d \tilde{s}(v)=\left.\frac{d}{d t}\right|_{t=0}(s(\gamma(t)) \circ g(\gamma(t)))=\left.\left.d R_{g}\right|_{s(x)} \circ d s\right|_{x}+d L_{s(x) g_{0}}\left(d g_{0}^{-1} g\right)_{x}
$$

Sei $s^{*} \omega=: A$ und $\tilde{s}^{*} \omega=: \tilde{A}$. Dann ist

$$
\begin{gathered}
\tilde{A}_{x}=\omega \circ d \tilde{s}_{x}=\left(s^{*} R_{g}^{*} \omega\right)_{x}+\left.\omega_{\tilde{s}\left(x_{0}\right)} \circ d L_{\tilde{s}\left(x_{0}\right)} d\left(g_{0}^{-1} g\right)\right|_{x_{0}} \\
=\operatorname{Ad}\left(g^{-1}\right) A_{x}+\left.g^{-1}\left(x_{0}\right) d g\right|_{x_{0}}
\end{gathered}
$$

da

$$
d L_{\tilde{s}(x)} d\left(g_{0}^{-1} g\right)_{x}(v) \in T_{\tilde{s}(x)}^{\prime} P \text { für } v \in T_{x} M
$$

\subsection*{6.24 Bemerkung}
Sei $\left\{\left(U_{\alpha}, \varphi_{\alpha}\right) \mid \alpha \in \Lambda\right\}$ ein Bündelatlas für $P$. Seien $s_{\alpha}$ die durch $\varphi_{\alpha}$ gegebenen Schnitte mit $s_{\beta}=: s_{\alpha} g_{\alpha \beta}$ auf $U_{\alpha} \cap U_{\beta}$ und ist für jedes $\alpha \in \Lambda$ eine 1-Form $A_{\alpha} \in \Omega^{1}\left(U_{\alpha}, \mathfrak{g}\right)$ gegeben, sodass auf $U_{\alpha} \cap U_{\beta}$ gilt:

$$
A_{\beta}=\operatorname{Ad}\left(g_{\alpha \beta}\right)^{-1} A_{\alpha}+g_{\alpha \beta}^{-1} d g_{\alpha \beta}^{-1}
$$

so beschreiben die $A_{\alpha}$ einen $G$-Zusammenhang auf $P$.

\section*{Beweis:}
Sei $X \in T_{s_{\alpha}(x)} P$ Definiere

$$
\omega_{s_{\alpha}(x)}(X)=d L_{s_{\alpha}(x)}^{-1}\left(X-\left.d s_{\alpha}\right|_{x} \circ d \pi_{s_{\alpha}(x)}(X)\right)+A_{\alpha}\left(d \pi_{s_{\alpha}(x)}(X)\right)
$$

Für $X \in T_{s_{\alpha}(x) g} P$ setze $\omega_{s_{\alpha}(x) g}(X)=\operatorname{Ad}\left(g^{-1}\right) \omega_{s_{\alpha}(x)}\left(d R_{g^{-1}}(X)\right)$. Dann ist $s_{\alpha}^{*} \omega=A_{\alpha}$ und $\omega_{\alpha}$ eine $G$-Zusammenhangsform. Wegen Lemma 6.23 ist $\omega \mid \pi^{-1}\left(U_{\alpha}\right)$ unabhängig von der Wahl der Karte ( $U_{\alpha}, \varphi_{\alpha}$ ).

\subsection*{6.25 Aufgabe}
Es sei $P \xrightarrow{\pi} M$ ein $G$-Prinzipalfaserbündel mit einem $G$-Zusammenhang. Für $u \in P$ bezeichne $P(u) \subset P$ die Menge der Punkte, die man mit $u$ durch eine (stückweise differenzierbare) Horizontalkurve verbinden kann, und es sei

$$
\operatorname{Hol}(u):=\{g \in G u g \in P(u)\}
$$

Man zeige, dass $\operatorname{Hol}(u)$ eine Untergruppe von $G$ ist.

\subsection*{6.26 Aufgabe}
Es werde vorausgesetzt, dass $\operatorname{Hol}(u) \subset G$ abgeschlossen ist (in der Tat ist das immer der Fall). Man zeige, dass $P(u) \subset P$ eine Reduktion der Strukturgruppe von $P$ auf $\operatorname{Hol}(u)$ definiert.





\section*{Lineare Zusammenhänge}
\subsection*{Bemerkung}
Ist $E \rightarrow M$ ein Vektorraumbündel, so ist $T_{e} E_{x}=E_{x}$, also $T^{\prime} E=\pi^{*} E$, d.h., die kurze exakte Sequenz, deren Spaltung ein Zusammenhang definiert, ist durch $0 \rightarrow \pi^{*} E \rightarrow T E \rightarrow \pi^{*} T M \rightarrow 0$ gegeben.

\subsection*{Definition}
Ein $G L(n, \mathbb{R})$-Zusammenhang auf einem Vektorraumbündel heißt ein linearer Zusammenhang.

\subsection*{Lemma}
Ein Zusammenhang für die lokal triviale Faserung eines Vektorraumbündels ist genau dann linear, wenn für jeweils zwei Horizontalkurven $\alpha_{1}, \alpha_{2}$ über $\alpha$ auch $\lambda \alpha_{1}+\mu \alpha_{2}$ eine Horizontalkurve ist.

\section*{Beweis:}
Ein Zusammenhang ist nach Definition genau dann linear, wenn er vollständig ist und die zugehörige Parallelverschiebung linear ist. Damit folgt die Behauptung, da die Horizontalkurven $\alpha_{i}$ über $\alpha$ durch

$$
\alpha_{i}(t)=\left[\tau_{\alpha_{[0, t]}}^{P}(p), v_{i}\right]=\tau_{\alpha_{[0, t]}}\left(\alpha_{i}(0)\right) \text { mit } \alpha_{i}(0)=\left[p, v_{i}\right]
$$

gegeben sind. Dabei ist $\tau^{P}$ die Parallelverschiebung im zugrundeliegenden Prinzipalbündel.

\subsection*{Bemerkung}
Ist $E \rightarrow M$ ein Vektorraumbündel, $\operatorname{dim} M=m, \operatorname{dim} E=m+k$, so ist $J^{1} E \rightarrow M$ ein Vektorraumbündel der Faserdimension $k(m+1)$.\\
Ein durch einen Schnitt $\sigma: E \rightarrow J^{1}(\pi)$ gegebener Zusammenhang ist genau dann ein linearer Zusammenhang, wenn er linear über $M$ ist, d.h., wenn er einen Vektorraumbündelhomomorphismus $E \rightarrow J^{1}(\pi)$ beschreibt, denn sind $s_{1}$ und $s_{2}$ horizontale Schnitte an der Stelle $x$, dann ist $H$ nach Lemma 7.3 linear, wenn $\lambda s_{1}+\mu s_{2}$ ein horizontaler Schnitt ist und umgekehrt.

\subsection*{Satz}
Die kovariante Ableitung eines Zusammenhangs ist genau dann linear, wenn der Zusammenhang linear ist.

\section*{Beweis:}
Seien $s_{i} \in \Gamma E, v \in T_{x} M, d s_{i}(v)=: \bar{v}_{i}+v_{i}^{\prime}$ die Zerlegung in horizontalen und vertikalen Anteil, also $\nabla_{v} s_{i}=v_{i}^{\prime}$. Es ist

$$
d\left(\lambda s_{1}(v)+\mu s_{2}(v)\right)=d L_{\lambda}\left(\bar{v}_{1}\right)+d L_{\mu}\left(\bar{v}_{2}\right)+\lambda v_{1}^{\prime}+\mu v_{2}^{\prime}
$$

Nach Lemma 8.3 ist der Zusammenhang genau dann linear, wenn die ersten beiden Summanden stets wieder horizontal sind, also wenn

$$
\nabla\left(\lambda s_{1}+\mu s_{2}\right)=\lambda \nabla s_{1}+\mu \nabla s_{2}
$$

gilt

\subsection*{Bemerkung}
Ist $\nabla$ die kovariante Ableitung eines linearen Zusammenhangs, $f \in C^{\infty}(M)$ und $s \in \Gamma E$, dann gilt

$$
\nabla(f s)=f \nabla s+d f s(*)
$$

denn ist $\beta$ eine $v$ repräsentierende Kurve, so ist nach 5.14 und 5.16

$$
\nabla_{v} s=\left.\frac{d}{d t}\right|_{t=0} \tau_{\beta_{[t, 0]}}(s \circ \beta)
$$

also

$$
\begin{aligned}
\nabla_{v}(f s) & =\left.\frac{d}{d t}\right|_{t=0} \tau_{\beta_{[t, 0]}}((f \circ \beta) \cdot s \circ \beta) \\
& =\left.\frac{d}{d t}\right|_{t=0}\left((f \circ \beta) \cdot \tau_{\beta_{[t, 0]}}(s \circ \beta)\right) \\
& =d f(v) \cdot s(x)+f(x) \nabla_{v} s
\end{aligned}
$$

\subsection*{Lemma}
Ist $\nabla: J^{1} E \rightarrow \operatorname{Hom}(T M, E)$ eine lineare Abbildung, die die Produktformel (*) erfüllt, so ist $\nabla$ die kovariante Ableitung eines linearen Zusammenhangs.

\section*{Beweis:}
Definiere einen Schnitt $\sigma$ in $J^{1} E \rightarrow E$ durch $\sigma(e)=[s]_{x}^{1}-\nabla[s]_{x}^{1}$. Dies ist unabhängig von der Wahl des Schnittes $s$ mit $s(x)=e$ : Aufgrund der Linearität genügt es dazu zu zeigen: für $s$ mit $s(x)=0$ ist $\nabla[s]_{x}^{1}=[s]_{x}^{1}$. Sei $\left(s_{1}, \ldots, s_{n}\right)$ eine lokale Familie lokaler Schnitte, die in einer Umgebung von $x$ eine Basis bilden.\\
Sei $s=\sum \lambda_{i} s_{i}$, also $\lambda_{i}(x)=0$ für $i=1, \ldots n$. Dann ist $d s_{x}=\sum d \lambda_{i}{ }_{x} s_{i}(x)$ und $(\nabla s)_{x}=\sum d \lambda_{i} \times s_{i}(x)$, also $\nabla[s]_{x}^{1}=[s]_{x}^{1}$.

\subsection*{Definition}
Seien $E_{i} \rightarrow M$ Vektorraumbündel. Eine Abbildung $\mathcal{D}: \Gamma E_{1} \rightarrow \Gamma E_{2}$ heißt Differentialoperator $k$-ter Ordung, wenn eine differenzierbare Abbildung über $M$ existiert $D^{k}: J^{k} E_{1} \rightarrow E_{2}$ mit $(\mathcal{D} s)(x)=D^{k}[s]_{x}^{k}$ für alle $x \in M$. Ein Differentialoperator $\mathcal{D}$ heißt linear, wenn es ein $k$ gibt, sodass $D^{k}$ ein Vektorraumbündelhomomorphismus ist.

\subsection*{Satz}
Ein linearer Differenzialoperator $\mathcal{D}: \Gamma E \rightarrow \Omega^{1}(M, E)$ ist genau dann kovariante Ableitung eines Zusammenhangs, wenn die Produktformel

$$
\nabla f s=d f s+f \nabla s
$$

gilt für alle $f \in C^{\infty}(M), s \in \Gamma E$.

\section*{Beweis:}
Wir haben schon gezeigt, dass lineare kovariante Ableitungen in Vektorraumbündeln die Produktformeln erfüllen. Ist die Produktformel erfüllt, so ist $\mathcal{D}$ ein Differenzialoperator 1. Ordnung, denn ist $[s]_{p}^{1}=0$, also lokal $s=\sum \lambda_{i} e_{i}$ mit $\left.d \lambda_{i}\right|_{p}=0$ und $\lambda_{i}(p)=0$, so folgt $\mathcal{D} s(p)=0$, also folgt aus der Linearität, dass

$$
\mathcal{D}: J^{1}(E) \rightarrow T^{*} M \otimes E,[s]_{x}^{1} \mapsto(\mathcal{D} s)(x)
$$

wohldefiniert ist. Der Zusammenhang ist dann durch

$$
\sigma: E \rightarrow J^{1} E, \sigma(e)=[s]_{x}^{1}-\mathcal{D}[s]_{x}^{1}
$$

für $s$ mit $s(x)=e$ wohldefiniert und linear.

\subsection*{Lemma}
Ist $P$ das $E \rightarrow M$ zugrundeliegende Prinzipalbündel, $\omega$ eine $G$-Zusammenhangsform auf $P$, dann gilt für die zugehörige kovariante Ableitung in $E$

$$
\nabla_{X}[s, v]=\left[s, d v(X)+s^{*} \omega(X) v\right]
$$

\section*{Beweis:}
Ist $s$ ein horizontaler Schnitt bei $p$ und $v$ konstant und $X \in T_{p} M$, so sind beide Seiten null.\\
Sei jetzt $\left(v_{1}, \ldots, v_{n}\right)$ eine Basis von $\mathbb{R}^{n}, s$ bei $p$ horizontal und $v=\sum \lambda_{i} v_{i}$. Dann ist die linke Seite für $X \in T_{\pi(p)} M$

$$
\nabla_{X}[s, v]=\sum_{i} d \lambda_{i}(X)\left[s, v_{i}\right]
$$

und die rechte Seite

$$
\sum_{i}\left[s, d\left(\lambda_{i} v_{i}\right)(X)+\lambda_{i}\left(s^{*} \omega\right)(X) v_{i}\right]=\sum_{i} d \lambda_{i}(X)\left[s, v_{i}\right]+0
$$

\subsection*{Bemerkung}
\begin{enumerate}
  \item Sind $v_{j}$ für $j=1, \ldots, n$ lokale Schnitte in $T M \mid U$, sodass $\left(v_{1}(x), \ldots, v_{n}(x)\right)$ für jedes $x \in U$ eine Basis von $T_{x} M$ ist, so ist $s:=\left(v_{1}, \ldots, v_{n}\right) \in \left.\Gamma P_{G L}(M)\right|_{U}$. Dann ist $v_{j}=\left[s, e_{j}\right]$, also
\end{enumerate}

$$
\nabla_{X} v_{j}=\left[s, s^{*} \omega(X) e_{j}\right]=\left[s, \sum_{i}\left(s^{*} \omega\right)(X)_{i j} e_{i}\right]=: \sum_{i}\left(s^{*} \omega(X)\right)_{i j} v_{i} .
$$

\begin{enumerate}
  \setcounter{enumi}{1}
  \item Sind $\omega, \tilde{\omega}$ zwei Zusammenhangsformen, so ist $\omega-\tilde{\omega} \in \Omega^{1}\left(M, P \times_{\mathrm{Ad}} \mathfrak{g}\right)$. Ist $E=P \times_{G} V$, so kann man $P \times_{\text {Ad }} \mathfrak{g}$ auffassen als Teilmenge von $\operatorname{End}(E)$ via
\end{enumerate}

$$
\left(P \times_{\text {Ad }} \mathfrak{g}\right) \times\left(P \times_{G} V\right) \rightarrow\left(P \times_{G} V\right), \quad([p, X],[p, v]) \mapsto[p, X \cdot v]
$$

Dies ist wohldefiniert.\\
Ist $\tilde{\nabla}$ die kovariante Ableitung zu $\nabla+A$ mit $A \in \Omega^{1}(M, \operatorname{End}(E))$, so ist also

$$
\tilde{\nabla}_{X}(v)=\nabla_{X} v+A(X) v
$$

\subsection*{Korollar}
Es gilt $\nabla=d^{H}: \Omega^{0}(M, E) \rightarrow \Omega^{1}(M, E)=\Gamma(\operatorname{Hom}(T M, E))$, denn

$$
\begin{aligned}
\nabla_{X}[s, v] & =\left[s, d v(X)+s^{*} \omega(X)(v)\right] \\
d^{H}[s, \hat{v}] & =\left[s, d \hat{v}(d s(X))+\left(s^{*} \omega\right)(X)(\hat{v})\right]
\end{aligned}
$$

wobei $\hat{v} \in C_{\text {inv }}^{\infty}(P, V), v \circ \pi=\hat{v}$.\\
Erinnere: Ist $\alpha: H \rightarrow G$ ein Liegruppenhomomorphismus, so ist eine $\alpha$ Version von $P \rightarrow M$ ein Paar ( $Q, f$ ) bestehend aus einem $H$-Prinzipalbündel $Q$ und einem Prinzipalbündelhomomorphismus $f: Q \rightarrow P$. Ist $\alpha$ die Inklusion, so nennt man die Äquivalenzklasse von $Q$ eine Reduktion. Wir fassen $Q \subseteq P$ auf. Es ist $Q \times_{H} V=P \times_{G} V$ und $P=Q \times_{H} G$.

\subsection*{Definition}
Ein $G$-Zusammenhang auf $P$ (bzw. $P \times_{\rho} V$ ) heißt reduzibel auf $H$, falls es eine $H$-Version $Q$ von $P$ und einen $H$-Zusammenhang $\omega^{H}$ auf $Q$ gibt, sodass der durch $\omega^{H}$ auf $Q \times_{H} G=P$ induzierte Zusammenhang mit $\omega^{G}$ übereinstimmt. $\omega^{G}$ heißt dann die $G$-Erweiterung von $\omega^{H}$.

\subsection*{Korollar}
Ist $H \subseteq G$ eine abgeschlossene Untergruppe, $Q \subseteq P$ eine Reduktion auf $H$, so ist ein $G$-Zusammenhang genau dann reduzibel auf $H$, falls die Einschränkung der $G$-Zusammenhangsform $\omega^{G}$ auf $T Q$ Werte in $\mathfrak{h}$ annimmt.

\subsection*{Korollar}
Ist $P \rightarrow M$ ein $O(n)$-Prinzipalbündel, $E=P \times_{O(n)} V$ mit der durch $O(n)$ induzierten Bündelmetrik, so gilt für jede kovariante Ableitung auf $E$, die auf $O(n)$ reduzibel ist:

$$
X<v, w>_{E}=<\nabla_{X} v, w>_{E}+<v, \nabla_{X} w>_{E}
$$

\section*{Beweis:}
Sei $v=[s, \hat{v}], w=[s, \hat{w}]$. Dann ist $\langle v, w\rangle_{E}=\langle\hat{v}, \hat{w}\rangle_{\mathbb{R}^{n}}$, also

$$
\begin{aligned}
X<v, w>_{E} & =<d \hat{v}(X), \hat{w}>_{\mathbb{R}^{n}}+<\hat{v}, d \hat{w}(X)>_{R^{n}} \\
<\nabla_{X} v, w>_{E} & =<d \hat{v}(X)+\left(s^{*} \omega\right)(X) \hat{v}, \hat{w}>_{\mathbb{R}^{n}} \\
<v, \nabla_{X} w>_{E} & =<\hat{v}, d \hat{w}(X)+\left(s^{*} \omega\right)(X) \hat{w}>
\end{aligned}
$$

$\operatorname{Da}\left(s^{*} w\right)(X) \in o(n)$, ist also $\left\langle\nabla_{X} v, w\right\rangle+\left\langle v, \nabla_{X} w\right\rangle=\langle d \hat{v}(X), \hat{w}\rangle_{\mathbb{R}^{n}} +<\hat{v}, d \hat{w}(X)>_{\mathbb{R}^{n}}$.

\section*{Die Aktion der Eichtransformationsgruppe auf dem Raum der Zusammenhänge und Rechnen in Bündelkarten}
\subsection*{Notation}
Die Gruppe $\mathcal{G}=\Gamma \operatorname{Aut}(P)=\Gamma\left(P \times_{\text {konj }} G\right)$ heißt Eichtransformationsgruppe. Ein Element $\varepsilon \in \mathcal{G}$ definiert durch $[p, g]=\varepsilon: P \rightarrow P, p \mapsto p g$ einen Bündelautomorphismus auf $P$ und einen Bündelautomorphismus auf $P \times_{G} V$ durch $\varepsilon([p, v])=[\varepsilon(p), v]$.

\subsection*{Definition und Bemerkung}
Ist $H$ ein $G$-Zusammenhang, $\varepsilon \in \Gamma(\operatorname{Aut}(P))$, so ist der eichtransformierte Zusammenhang $\varepsilon_{H}$ durch

$$
{ }^{\varepsilon} H_{\varepsilon(p)}=\left.d \varepsilon\right|_{p}\left(H_{p}\right)
$$

definiert. Durch ${ }^{\varepsilon} H$ ist offenbar wiederum ein $G$-Zusammenhang wohldefiniert.

\subsection*{Notiz}
Für die Parallelverschiebung zu ${ }^{\varepsilon} H$ gilt: ${ }^{\varepsilon} \tau_{\gamma}=\varepsilon \circ \tau_{\gamma} \circ \varepsilon^{-1}$, denn ist $\gamma_{e}$ eine Horizontalkurve zu $H$, so ist $\varepsilon \circ \gamma_{e}$ eine Horizontalkurve $\mathrm{zu}^{\varepsilon} H$, also ${ }^{\varepsilon} \tau_{\gamma}(\varepsilon(p))= \varepsilon \circ \tau_{\gamma}(p)$.

\subsection*{Notiz}
Ist $\omega^{\prime}$ die Vertikalprojektion zu $H,{ }^{\varepsilon} \omega^{\prime}$ die Vertikalprojektion zu ${ }^{\varepsilon} H$, so ist

$$
\varepsilon \omega^{\prime}=d \varepsilon \circ \omega^{\prime} \circ d \varepsilon^{-1}
$$

denn

$$
{ }^{\varepsilon} \omega^{\prime}(d \varepsilon(v))=^{\varepsilon} \omega^{\prime}\left(d \varepsilon(\bar{v})+d \varepsilon\left(v^{\prime}\right)\right)={ }^{\varepsilon} \omega^{\prime}\left(d \varepsilon\left(v^{\prime}\right)\right)=d \varepsilon\left(v^{\prime}\right)=d \varepsilon(\omega(v))
$$

wobei $v=\bar{v}+v^{\prime}, \bar{v} \in H_{e}, v^{\prime} \in T_{e}^{\prime} E$.\\
Damit folgt für eine $G$-Zusammenhangsform $\omega$ auch sofort ${ }^{\varepsilon} \omega_{\varepsilon(p)} \circ d R_{\varepsilon(p)}=\omega_{p}$, also

$$
\varepsilon \omega=\left(\varepsilon^{-1}\right)^{*} \omega
$$

\subsection*{Notiz}
Für die kovariante Ableitung in Vektorraumbündeln gilt

$$
{ }^{\varepsilon} \nabla=\varepsilon \circ \nabla \circ \varepsilon^{-1}
$$

wobei ${ }^{\varepsilon} \nabla$ die kovariante Ableitung zu ${ }^{\varepsilon} H$ ist, denn ${ }^{\varepsilon} \nabla s={ }^{\varepsilon} \omega^{\prime} \circ d s=d \varepsilon \circ \omega^{\prime} d\left(\varepsilon^{-1} s\right)$.\\
Identifiziert man $T_{[p, v]}^{\prime} E=E_{\pi(p)}$, so ist

$$
\begin{aligned}
d \varepsilon_{[p, v]}([p, X]) & =\left.\frac{d}{d t}\right|_{t=0} \varepsilon([p, v+t X]) \\
& \left.=\left.\frac{d}{d t}\right|_{t=0}[\varepsilon(p), v+t X]\right) \\
& =[\varepsilon(p), X] \\
& =\varepsilon([p, X])
\end{aligned}
$$

Für Vektorraumbündel $E=P \times_{G} V$ ist $\Gamma(\operatorname{Aut}(P)) \subseteq \Gamma \operatorname{End}(E)$, ein $G$ Zusammenhang in $P$ induziert sowohl in $E$ also auch in $\operatorname{End}(E)$ einen Zusammenhang.

\subsection*{Lemma}
Ist $\nabla$ ein $G$-Zusammenhang auf $E=P \times_{G} V$ und $\nabla^{\text {End }}$ der vom selben $G$ Zusammenhang induzierte Zusammenhang in $\operatorname{End}(E)=P \times_{\rho} \operatorname{End}(V)$, dann gilt

$$
\nabla^{\text {End }} \phi=\nabla \circ \phi-\phi \circ \nabla
$$

denn $\tau_{\gamma}(\phi) \tau_{\gamma}(e)=\tau_{\gamma}(\phi(e))$ für $\pi(\phi)=\pi(e)$. Damit folgt mit 5.16 für Schnitte $\phi$ und $e$ in End $E$ und $E$ längs einer Kurve $\gamma$

$$
\frac{\nabla}{d t}(\phi(e))=\left(\frac{\nabla}{d t} \phi\right) e+\phi\left(\frac{\nabla}{d t} e\right)
$$

\subsection*{Korollar}
Sei $A \in \Omega^{1}(M, \operatorname{End} E)$ wie in 7.11.2. Dann gilt:

$$
(\nabla+A)^{\text {End }} \phi=\nabla^{\text {End }} \phi+A \circ \phi-\phi \circ A=: \nabla^{\text {End }} \phi+[A, \phi] .
$$

Wir können jetzt also 8.6 auf Eichtransformationen anwenden.

\subsection*{Korollar}
Es gilt

$$
\begin{array}{r}
{ }^{\varepsilon} \nabla-\nabla=-(\nabla \varepsilon) \circ \varepsilon^{-1}=\varepsilon \circ\left(\nabla \varepsilon^{-1}\right), \\
\operatorname{denn}^{\varepsilon} \nabla=\varepsilon \circ \nabla \circ \varepsilon^{-1}=\nabla+\varepsilon \circ\left(\nabla \varepsilon^{-1}\right)=\nabla-(\nabla \varepsilon) \circ \varepsilon^{-1} .
\end{array}
$$

\subsection*{Korollar und Notation}
Es gilt

$$
\varepsilon(\nabla+A)=: \nabla+{ }^{\varepsilon} A \text { mit }^{\varepsilon} A=\varepsilon \circ A \circ \varepsilon^{-1}+\varepsilon \circ \nabla \varepsilon^{-1}
$$

$\operatorname{denn}{ }^{\varepsilon}(\nabla+A)=^{\varepsilon} \nabla+\varepsilon \circ A \circ \varepsilon^{-1}$.\\
Ist $s$ ein Schnitt in $P$, so induziert $s$ eine Karte für jedes induzierte Bündel:

$$
\begin{aligned}
P \times_{G} V & \rightarrow M \times V \\
{[s(x), v] } & \mapsto(x, v)
\end{aligned}
$$

Ist $\sigma \in \Gamma\left(P \times_{G} V\right)$, so wird $\sigma$ bezüglich $s$ durch $\sigma^{(s)} \in C^{\infty}(M, V)$ beschrieben, mit $\sigma(x)=\left[s(x), \sigma^{(s)}(x)\right]$. Ist $\tilde{s}=s \tilde{g}$, so ist $\sigma^{(\tilde{s})}=\tilde{g}^{-1} \sigma^{(x)}$, d.h. der Kartenwechsel bewirkt einen Bündelautomorphismus des trivialen Bündels $M \times V \rightarrow M \times V,(x, v) \mapsto\left(x, \tilde{g}^{-1} v\right)$.\\
Ist $\varepsilon \in \Gamma(P \times$ konj $G)$ eine Eichtransformation, die bezüglich $s$ durch $\varepsilon^{(s)}$ beschrieben wird, so wird sie bezüglich $\tilde{s}$ durch

$$
\tilde{g}^{-1} \varepsilon^{(s)} \tilde{g}=\varepsilon^{(\tilde{s})}
$$

beschrieben.

\subsection*{Satz}
Bezüglich Karten, die durch $s$ gegeben sind, ist $\nabla$ durch $d+s^{*} \omega$ gegeben.\\
Genauer gilt:

$$
\left(\nabla_{X} \sigma\right)^{(s)}=d \sigma^{(s)}(X)+s^{*} \omega(X) \sigma^{(s)}(X)
$$

Bezüglich $\tilde{s}=s g$ ist sie dann durch $d+\tilde{s}^{*} \omega$ gegeben und es gilt

$$
\tilde{s}^{*} \omega=g^{-1} A g+g^{-1} d g
$$

für $A=s^{*} \omega(\operatorname{vgl} .6 .25)$.

\section*{Die Krümmung eines Zusammenhangs}
\subsection*{Notation}
Ist $H$ ein Zusammenhang auf $E, v \in T_{e} E$, so schreiben wir $v=\bar{v}+v^{\prime}$, mit $\bar{v} \in H_{e}$ und $v^{\prime} \in T_{e}^{\prime} E$.\\
Ist $v \in T_{\pi(p)} M$, so bezeichnet $\tilde{v} \in T_{p} E$ einen Vektor mit $d \pi(\tilde{v})=v$ und $\overline{\tilde{v}}=: \bar{v}, \tilde{v}^{\prime}=: v^{\prime}$ den horizontalen und vertikalen Anteil. Vertikale Vektoren werden (auf Prinzipalbündeln durch $d L_{p}(\hat{v}), \hat{v} \in \mathfrak{g}$ ) zu vertikalen Vektorfeldern fortgesetzt, horizontale Vektoren (auf Prinzipalbündeln) werden zu (rechtsinvarianten) horizontalen Vektorfeldern fortgesetzt.

\subsection*{Definition}
Sei $\omega \in \Omega_{\mathrm{inv}}^{1}(P, \mathfrak{g})$ die Zusammenhangsform eines $G$-Zusammenhangs, so ist die Krümmungsform durch $\Omega^{\omega} \in \Omega_{\text {inv }, \text { hor }}^{2}(P, \mathfrak{g})$

$$
\Omega^{\omega}=D^{\omega} \omega
$$

definiert. Die dadurch definierte 2 -Form $F^{\omega} \in \Omega^{2}\left(M, P \times_{\text {Ad }} \mathfrak{g}\right)$ heißt die Krümmung von $\omega$.

\subsection*{Bemerkung}
Ist $E=P \times{ }_{\rho} V$, so ist $F^{\omega} \in \Omega^{2}(M, \operatorname{End}(E))$. Ist der Zusammenhang $\omega$ durch $\nabla$ beschrieben, dann schreibt man auch $F^{\nabla}$.

\subsection*{Lemma}
Sei $\Omega^{\omega}$ Krümmung eines $G$-Zusammenhangs in $P \rightarrow M$. Es gilt

$$
d L_{p} \Omega_{p}^{\omega}(X, Y)=\overline{[X, Y]}-[\bar{X}, \bar{Y}] .
$$

\section*{Beweis:}
$$
\begin{aligned}
d L_{p} \Omega^{\omega}(X, Y) & =d L_{p}(d \omega(\bar{X}, \bar{Y}))=d L_{p}(\bar{X} \omega(\bar{Y})-\bar{Y} \omega(\bar{X})-\omega([\bar{X}, \bar{Y}])) \\
& =-\omega^{\prime}([\bar{X}, \bar{Y}]) \stackrel{*}{=} \overline{[X, Y]}-[\bar{X}, \bar{Y}]
\end{aligned}
$$

(*) gilt, denn

$$
[X, Y]=[\bar{X}, \bar{Y}]+\left[X^{\prime}, \bar{Y}\right]+\left[\bar{X}, Y^{\prime}\right]+\left[X^{\prime}, Y^{\prime}\right]=[\bar{X}, \bar{Y}]+\left[X^{\prime}, Y^{\prime}\right]
$$

also ist $\overline{[X, Y]}=\overline{[\bar{X}, \bar{Y}]}$, denn für $X^{\prime}, Y^{\prime}$ mit

$$
\begin{aligned}
X^{\prime}(p) & =d L_{p}(\hat{X}), \hat{X} \in \mathfrak{g} \\
Y^{\prime}(p) & =d L_{p}(\hat{Y}), \hat{Y} \in \mathfrak{g}
\end{aligned}
$$

ist $\left[X^{\prime}, Y^{\prime}\right]=d L_{p}[\hat{X}, \hat{Y}]$, also vertikal.

\subsection*{Bemerkung}
Auf lokal trivialen Faserungen mit Zusammenhangsform $\omega^{\prime} \in \Omega^{1}\left(E, T^{\prime} E\right)$ ist die Krümmung durch $\Omega^{\omega}(X, Y)=\overline{[X, Y]}-[\bar{X}, \bar{Y}]=-\omega^{\prime}([\bar{X}, \bar{Y}]) \in T^{\prime} E$ definiert.

\subsection*{Korollar}
Es gilt für $\psi \in \Gamma\left(P x_{\rho} V\right)$

$$
F^{\nabla}(X, Y) \psi=\nabla_{X} \nabla_{Y} \psi-\nabla_{Y} \nabla_{X} \psi-\nabla_{[X, Y]} \psi
$$

\section*{Beweis:}
Nach 6.17 gilt:

$$
\begin{aligned}
F^{\nabla}(X, Y)[s, \tilde{\psi}] & =\left[s, \rho_{*}\left(\Omega^{\omega}(\tilde{X}, \tilde{Y})\right) \psi\right] \stackrel{(* *)}{=}-\left[s, \Omega^{\omega}(\tilde{X}, \tilde{Y}) \tilde{\psi}\right] \\
& =[s,-\overline{[X, Y]} \tilde{\psi}+[\bar{X}, \bar{Y}] \tilde{\psi}] \\
& =\nabla_{X} \nabla_{Y} \psi-\nabla_{Y} \nabla_{X} \psi-\nabla_{[X, Y]} \psi
\end{aligned}
$$

wobei (**) gilt, denn

$$
(\xi \tilde{\psi})_{(p)}=\frac{d}{d t}(p \exp (t \xi))=\frac{d}{d t}(\rho(\exp (-t \xi)) \tilde{\psi}(p))=-\rho_{*}(\xi) \tilde{\psi}(p)
$$

Einmal wird $\xi \in \mathfrak{g}$ als Vektor und einmal als Endomorphismus aufgefasst.

\subsection*{Satz (Strukturgleichung für die Krümmung)}
Es gilt

$$
\Omega_{p}^{\omega}(\tilde{X}, \tilde{Y})=d \omega_{p}(\tilde{X}, \tilde{Y})+\left[\omega_{p}(\tilde{X}), \omega_{p}(\tilde{Y})\right]
$$

\section*{Beweis:}
Beide Seiten sind in $\Omega^{2}(P, \mathfrak{g})$. Sei $\tilde{X}=\bar{X}+X^{\prime}$. Aufgrund der Linearität genügt es, die drei folgenden Fälle zu betrachten:

\begin{enumerate}
  \item Fall: $X^{\prime}=Y^{\prime}=0$.
\end{enumerate}

Dann ist $\Omega(\tilde{X}, \tilde{Y})=d \omega(\bar{X}, \bar{Y})$ und stimmt mit der rechten Seite übrein, da $\omega(\bar{X})=\omega(\bar{Y})=0$.\\
2. Fall: $\bar{X}=\bar{Y}=0$.

Dann ist $\Omega_{p}^{\omega}(\tilde{X}, \tilde{Y})=0$ und

$$
\begin{aligned}
d \omega_{p}(\tilde{X}, \tilde{Y})+\left[\omega_{p}(\tilde{X}), \omega_{p}(\tilde{Y})\right] & =d \omega\left(X^{\prime}, Y^{\prime}\right)+\left[d L_{p}^{-1} X^{\prime}, d L_{p}^{-1} Y^{\prime}\right] \\
& =X^{\prime} \omega\left(Y^{\prime}\right)-Y^{\prime} \omega\left(X^{\prime}\right)-\omega\left(\left[X^{\prime}, Y^{\prime}\right]\right)+\omega\left(\left[X^{\prime}, Y^{\prime}\right]\right)=0
\end{aligned}
$$

nach Wahl von $X^{\prime}$ und $Y^{\prime}$ wie in 9.1.\\
3. Fall: $X^{\prime}=0, \bar{Y}=0$.

Dann ist $\Omega^{\omega}(\tilde{X}, \tilde{Y})=0$ und

$$
d \omega(\tilde{X}, \tilde{Y})+[\omega(\tilde{X}), \omega(\tilde{Y})]=\bar{X} \omega\left(Y^{\prime}\right)-Y^{\prime} \omega(\bar{X})-\omega\left(\left[\bar{X}, Y^{\prime}\right]\right)=0
$$

nach Wahl von $\bar{X}$ und $Y^{\prime}$ wie in 9.1.

\subsection*{Bemerkung}
Dies ist kein Widerspruch zu 6.18, denn 6.18 ist nicht anwendbar, denn dort ist $\eta \in \Omega_{\text {inv }, \text { hor }}$.

\subsection*{Notation}
Ist $\alpha \in \Omega^{r}(M, \mathfrak{g}), \beta \in \Omega^{s}(M, \mathfrak{g})$, so schreibt man $[\alpha \wedge \beta] \in \Omega^{r+s}(M, \mathfrak{g})$ für

$$
[\alpha \wedge \beta]\left(v_{1}, \ldots, v_{r+s}\right)=\frac{1}{r!s!} \sum \operatorname{sign}(\tau)\left[\alpha\left(v_{\tau(1)}, \ldots, v_{\tau(r)}\right), \beta\left(v_{\tau(r+1)}, \ldots, v_{\tau(r+s)}\right)\right]
$$

\subsection*{Bemerkung}
Ist $\alpha \in \Omega^{1}(M, \mathfrak{g})$, so ist

$$
[\alpha \wedge \alpha](X, Y)=[\alpha(X), \alpha(Y)]-[\alpha(Y), \alpha(X)]=2[\alpha(X), \alpha(Y)] .
$$

Also kann man die Strukturformel schreiben als

$$
\Omega^{\omega}=d \omega+\frac{1}{2}[\omega \wedge \omega]
$$

\subsection*{Beispiel}
Sei $G / H$ ein reduktiver homogener Raum, also existiert ein $\mathfrak{m} \subseteq \mathfrak{g}$, sodass $\mathfrak{g}=\mathfrak{h} \oplus \mathfrak{m}$ mit

$$
[\mathfrak{h}, \mathfrak{h}] \subseteq \mathfrak{h} \text { und }[\mathfrak{h}, \mathfrak{m}] \subseteq \mathfrak{m} .
$$

Es ist

$$
\begin{aligned}
T G & =G \times \mathfrak{g}=G \times(\mathfrak{h} \oplus \mathfrak{m}) \\
T(G / H) & =G \times_{\mathrm{Ad}} \mathfrak{m}
\end{aligned}
$$

Das $H$-Prinzipalbündel $G$ trägt den kanonischen Zusammenhang: $\omega=\operatorname{pr}_{\mathfrak{h}}(\theta)$, wobei $\theta \in \Omega^{1}(G, \mathfrak{g})$ die Maurer-Cartan-Form ist $\theta=\left(\left.d L_{g}\right|_{1}\right)^{-1}$. Für linksinvariante Vektorfelder $X, Y$ gilt dann $\nabla_{X} Y=0$, also ist der kanonische Zusammenhang nicht notwendig torsionsfrei.\\
Für die Maurer Cartan Form gilt $d \theta=-\frac{1}{2}[\theta \wedge \theta]$, denn

$$
\begin{aligned}
d \theta(X, Y) & =X \theta(Y)-Y \theta(X)-\theta([X, Y])=-\left(\left.d L_{g}\right|_{1}\right)^{-1}[X, Y] \\
& =-\frac{1}{2}[\theta \wedge \theta](X, Y)
\end{aligned}
$$

also

$$
\begin{aligned}
\Omega^{\omega} & \left.=d\left(\operatorname{pr}_{\mathfrak{h}} \theta\right)+\frac{1}{2}\left[\operatorname{pr}_{\mathfrak{h}} \theta \wedge \operatorname{pr}_{\mathfrak{h}} \theta\right]\right) \\
& =-\frac{1}{2} \operatorname{pr}_{\mathfrak{h}}[\theta \wedge \theta]+\frac{1}{2}\left[\operatorname{pr}_{\mathfrak{h}} \theta \wedge \operatorname{pr}_{\mathfrak{h}} \theta\right] \\
& =\frac{1}{2}\left(\left[\operatorname{pr}_{\mathfrak{h}} \theta \wedge \operatorname{pr}_{\mathfrak{h}} \theta\right]-\operatorname{pr}_{\mathfrak{h}}[\theta \wedge \theta]\right)
\end{aligned}
$$

Ist $G / H$ symmetrisch, also $[\mathfrak{m}, \mathfrak{m}] \subseteq \mathfrak{h}$, so ist

$$
\operatorname{pr}_{\mathfrak{h}}[X, Y]=\left[\operatorname{pr}_{\mathfrak{h}} X, \operatorname{pr}_{\mathfrak{h}} Y\right]+\left[\operatorname{pr}_{\mathfrak{m}} X, \operatorname{pr}_{\mathfrak{m}} Y\right]
$$

also

$$
\Omega^{\omega}=-\frac{1}{2}\left[\operatorname{pr}_{\mathfrak{m}} \theta \wedge \operatorname{pr}_{\mathfrak{m}} \theta\right]
$$

Beispiel: $S^{n}=S O(n+1) / S O(n)$ (vgl. 8.11 im Skript Liegruppen). Dann ist\\
$\mathfrak{g}=s o(n+1), \mathfrak{h}=s o(n) \subseteq s o(n+1), \mathfrak{m}=\mathbb{R}^{n} \rightarrow s o(n+1), v \mapsto\left(\begin{array}{cc}0 & -v^{T} \\ v & 0\end{array}\right)$.\\
Es ist

$$
[v, w]=\left(\begin{array}{cc}
0 & 0 \\
0 & \left(-v_{i} w_{j}+v_{j} w_{i}\right)_{i j}
\end{array}\right)
$$

also ist $S^{n}$ symmetrisch, also ist\\
$\Omega(v, w)=\left(v_{i} w_{j}-v_{j} w_{i}\right)_{i, j}$ und $\left\langle\Omega^{\omega}\left(e_{i}, e_{j}\right) e_{k}, e_{l}\right\rangle= \begin{cases}1 & i=k, j=l, i \neq j \\ -1 & i=l, j=k, i \neq j \\ 0 & \text { sonst }\end{cases}$

\subsection*{Bianchi-Identität}
Es gilt:

$$
D^{\omega} \Omega^{\omega}=0, \text { also } d^{\nabla} F^{\nabla}=0 .
$$

\section*{Beweis:}
$$
d \Omega^{\omega}=d d \omega+\frac{1}{2} d[\omega \wedge \omega]=\frac{1}{2}([d \omega \wedge \omega]-[\omega \wedge d \omega])=[d \omega \wedge \omega]
$$

Also ist nach 6.18

$$
D^{\omega} \Omega^{\omega}=d \Omega^{\omega}+\left[\omega \wedge \Omega^{\omega}\right]=[d \omega \wedge \omega]+[\omega \wedge d \omega]+\frac{1}{2}[\omega \wedge[\omega \wedge \omega]]=0
$$

\subsection*{Lemma}
Ist $\rho: G \rightarrow \operatorname{End}(V)$ und $\eta \in \Omega_{\text {inv }, \text { hor }}^{k}(P, V)$, so ist

$$
D^{\omega} D^{\omega} \eta=\rho_{*} \Omega^{\omega} \wedge \eta
$$

\section*{Beweis:}
Es gilt

$$
\begin{aligned}
D^{\omega}\left(d \eta+\rho_{*} \omega \wedge \eta\right) & =d\left(\rho_{*} \omega \wedge \eta\right)+\rho_{*} \omega \wedge d \eta+\rho_{*} \omega \wedge \rho_{*} \omega \wedge \eta \\
& =\rho_{*} d \omega \wedge \eta+\left(\rho_{*} \omega \wedge \rho_{*} \omega\right) \wedge \eta \\
& =\rho_{*} d \omega \wedge \eta+\frac{1}{2} \rho_{*}[\omega \wedge \omega] \wedge \eta \\
& =\rho_{*} \Omega^{\omega} \wedge \eta
\end{aligned}
$$

\subsection*{Aufgabe:}
Zeigen Sie:

$$
F^{\nabla+A}=F^{\nabla}+d^{\nabla} A+[A \wedge A] .
$$

\section*{Das Utiyama-Lemma}
In diesem Kapitel behandeln wir die Frage, welche Information in $\Omega^{\omega}$ über $\omega$ enthalten ist.

\subsection*{Lemma}
Sei $\omega$ ein $G$-Zusammenhang, $\varepsilon \in \mathcal{G}$. Es gilt:

$$
\Omega^{\varepsilon \omega}=\operatorname{Ad}(\varepsilon) \Omega^{\omega} .
$$

Insbesondere gilt: Ist $U \subseteq M$ offen, so gilt: $\Omega^{\omega}=\left.0\right|_{U}$ genau dann, wenn $s^{*} \Omega^{\omega}=0$ bezüglich eines (dann jedes) lokalen Schnittes auf $U$.

\section*{Beweis:}
Wir bezeichnen horizontale Vektoren bezüglich ${ }^{\varepsilon} \omega$ mit $\bar{X}^{\varepsilon}$ und die Krümmung von ${ }^{\varepsilon} \omega$ mit $\Omega^{\varepsilon}$. Nach 9.4 und 8.4 gilt:

$$
\begin{aligned}
d L_{\varepsilon(p)} \Omega_{\varepsilon(p)}^{\varepsilon}\left(d \varepsilon_{p}(X), d \varepsilon_{p}(Y)\right) & =-{ }^{\varepsilon} \omega_{\varepsilon(p)}^{\prime}\left(\left[\overline{d \varepsilon_{p}(X)}{ }^{\varepsilon}, \overline{d \varepsilon_{p}(Y)}{ }^{\varepsilon}\right]\right. \\
& =-{ }^{\varepsilon} \omega_{\varepsilon(p)}^{\prime}\left(\left[d \varepsilon_{p}(\bar{X}), d \varepsilon_{p}(\bar{Y})\right]\right] \\
& =-d \varepsilon_{p} \omega_{p}^{\prime}([\bar{X}, \bar{Y}]) \\
& =d L_{\varepsilon(p)} \Omega_{p}([X, Y])
\end{aligned}
$$

Dabei gilt das letzte Gleichzeitszeichen, da $\varepsilon \circ L_{p}=L_{\varepsilon(p)}$. Also gilt

$$
\varepsilon^{*} \Omega^{\varepsilon}=\Omega
$$

Ist $\tilde{\varepsilon}: P \rightarrow G$ durch $\varepsilon(p)=p \tilde{\varepsilon}(p)$ definiert, so ist

$$
\Omega_{p}=\left(\varepsilon^{*} \Omega^{\varepsilon}\right)_{p}=R(\tilde{\varepsilon}(p))^{*} \Omega_{p}^{\varepsilon}=\operatorname{Ad}\left((\tilde{\varepsilon}(p))^{-1}\right) \Omega_{p}^{\varepsilon}=:\left(\operatorname{Ad}\left(\varepsilon^{-1}\right) \Omega^{\varepsilon}\right)_{p} .
$$

\subsection*{Korollar}
Für die Krümmung $F^{\varepsilon}$ zu ${ }^{\varepsilon} \omega$ gilt damit

$$
F^{\varepsilon}=\varepsilon \circ F^{\omega}
$$

denn für $p \in P_{x}$ ist

$$
F_{x}^{\varepsilon}=\left[\varepsilon(p), \Omega_{\varepsilon(p)}^{\varepsilon}\right]=\left[\varepsilon(p), \Omega_{p}\right]=\varepsilon \circ F_{x}^{\omega}
$$

\subsection*{Definition}
Eine Zusammenhang heißt flach, wenn $\Omega^{\omega}=0$ gilt.

\subsection*{Definition}
Ein Zusammenhang heißt lokal trivial, wenn es einen Bündelatlas gibt, in dessen Karten $\nabla$ durch $d$ beschrieben wird.

\subsection*{Bemerkung}
Ist $A=s^{*} \omega$, so ist $s^{*} \Omega^{\omega}=d A+\frac{1}{2}[A \wedge A]$. Also folgt mit 10.1, dass jeder lokal triviale Zusammenhang flach ist.

\subsection*{Definition}
Sei $M$ eine differenzierbare Mannigfaltigkeit. Dann versteht man unter einer Distribution auf $M$ ein Untervektorraumbündel von $T M$. Eine Distribution $D$ heißt lokal integrierbar, falls es es zu jedem $x \in M$ eine offene Umgebung $U$ und eine Familie von Untermannigfaltigkeiten $\left(N_{\lambda}\right)$ von $U$ gibt, sodass für alle $x^{\prime} \in U$ gilt $x^{\prime} \in N_{\lambda}$ für genau ein $\lambda \in \Lambda$ und $T_{x^{\prime}} N_{\lambda}=D_{x^{\prime}}$.

\subsection*{Bemerkung}
Eindimensionale Distributionen sind stets lokal integrierbar, da Vektorfelder integrierbar sind.

\subsection*{Satz von Frobenius}
Eine Distribution ist genau dann lokal integrierbar, wenn sie involutiv ist, d.h., falls für $X, Y \in D$ stets auch $[X, Y] \in D$ ist (Beweis z.B. in Lee, Differenzierbare Mannigfaltigkeiten).

\subsection*{Satz}
$H$ ist genau dann lokal integrierbar, wenn $\Omega^{H}=0$.

\section*{Beweis:}
Da $\Omega^{\omega}$ horizontal ist, gilt nach 9.4 gilt: $\Omega^{\omega}=0$ genau dann, wenn für horizontale Vektorfelder $X, Y$ auch $[X, Y]$ horizontal ist.

\subsection*{Satz}
Für einen Zusammenhang $\omega$ ein einem $G$-Prinzipalbündel $E \rightarrow M$ sind äquivalent:

\begin{enumerate}
  \item $\omega$ ist flach.
  \item $\omega$ ist lokal trivial.
  \item Für jedes $x \in M$ existiert eine offene Umgebung $U$ und ein $s \in \Gamma U$, sodass $\nabla s=0$ auf $U$ gilt.
  \item Für jedes $x \in M$ gibt es eine offene Umgebung $U$, sodass die Parallelverschiebung $\tau_{\gamma}: E_{x} \rightarrow E_{y}$ unabhängig von der Wahl von $\gamma$ von $x$ nach $y$ ist.
\end{enumerate}

\section*{Beweis:}
Es gilt offenbar $2 . \Rightarrow 3$. und $2 \Rightarrow 4,2 \Rightarrow 1$. wurde bereits in 10.5 gezeigt.\\
$1 \Rightarrow 2$ Ist $\Omega^{\omega}=0$, so gibt es nach dem Satz von Frobenius zu jedem $e \in E$ eine Integralmannigfaltigkeit $\tilde{M}$. Die Einschränkung der Bündelprojektion $\pi \mid \tilde{M} \rightarrow M$ ist ein lokaler Diffeomorphismus bei $x$, da $T_{x^{\prime}} \tilde{M}=H_{x^{\prime}}$ für alle $x^{\prime} \in \tilde{M}$, also ist $s=(\pi \mid U)^{-1}$ ein lokaler horizontaler Schnitt. Ist $s$ ein horizontaler Schnitt, so ist $H$ bezüglich der durch $s$ gegebenen Karte trivial, denn $s^{*} \omega=0$.\\
$4 \Rightarrow 3$ Sei $x \in U, e \in E_{x}, \gamma$ ein Weg von $x$ nach $y$. Dann ist ein Schnitt $s$ durch $s(y)=\tau_{\gamma}(e)$ wohldefiniert. Damit ist $s$ ein Schnitt mit $\nabla s=0$.\\
$3 \Rightarrow 1$ Ist $s$ ein flacher Schnitt, so ist $s(U) \subseteq E$ eine Integralmannigfaltigkeit für $H$, also $s^{*} \Omega^{\omega}=0$ und 1 . folgt mit 10.1.

\subsection*{Bemerkung}
Ist in einem Prinzipalbündel ein flacher lokaler Schnitt $s \in \Gamma P \mid U$, d.h. ein lokaler Schnitt mit $\nabla s=0$ gegeben, so ist natürlich auch für jedes $g \in G$ durch $\tilde{s}=s g$ ein Schnitt mit $\nabla \tilde{s}=0$ gegeben. Für flache Zusammenhänge folgt also in Prinzipalbündeln sofort, dass zu jedem $x \in P$ und $p \in P_{x}$ ein lokaler Schnit $s$ mit $s(x)=p$ und $\nabla s=0$ existiert.\\
Für $G$-Faserbündel folgt aus Satz 11.10 ebenfalls, dass $\omega$ genau dann flach ist, wenn $\omega$ lokal trivial ist, denn die Trivialisierung des Prinzipalbündels liefert eine des assoziierten Faserbündels. Da die Parallelverschiebungen eines $G$-Faserbündels denen im zugrunde liegenden Faserbündel entsprechen, ist die Äquivalenz aus 1. und 4. aus 10.10 für Faserbündel ebenso gültig. Damit folgt natürlich auch die Existenz flacher Schnitte. Umgekehrt folgt nicht aus der Existenz eines flachen Schnitts in einem Faserbündel, dass der Zusammenhang flach ist, denn z.B. der Nullschnitt ist stets flach. Es gilt aber

\subsection*{Korollar}
Ist $E \rightarrow M$ ein $n$-dimensionales Vektorraumbündel mit einem linearen Zusammenhang $\nabla$. Dann ist $\nabla$ genau dann flach, wenn es zu jedem $x \in M$ eine offene Umgebung $U \subseteq M$ gibt und $n$ flache Schnitte $s_{1}, \ldots, s_{n} \in \Gamma E \mid U$, sodass $\left(s_{1}(x), \ldots, s_{n}(x)\right)$ eine Basis von $E_{x}$ ist.

\subsection*{Korollar}
Ist $M$ einfach zusammenhängend, so ist ein $G$-Faserbündel $E \rightarrow M$ mit Zusammenhang $\nabla$ genau dann isomorph zu ( $M \times F, d$ ), wenn seine Krümmung null ist.

\subsection*{Definition}
\begin{enumerate}
  \item Eine Eichtransformation $\varepsilon \in \Gamma \operatorname{Aut}(P)$ heißt $x$-basiert, wenn $\varepsilon(x)=\operatorname{id}_{P_{x}}$.
  \item Sei $Z(P)=R(d \pi)$ das affine Bündel, dessen Schnitte die Zusammenhänge sind. Zwei 1-Jets $[\omega]_{x}^{1}$ und $[\tilde{\omega}]_{x}^{1}$ heißen strikt eichäquivalent, wenn für eine, dann für jede Wahl von Repräsentanten eine $x$-basierte Eichtransformation existiert mit $[\tilde{\omega}]_{x}^{1}=\left[{ }^{\varepsilon} \omega\right]_{x}^{1}$. Wir notieren die Äquivalenzklassen mit $[[\omega]]_{x}^{1}$.
\end{enumerate}

\subsection*{Utiyama-Lemma}
Es gilt:

$$
F_{x}^{\omega}=F_{x}^{\tilde{\omega}} \Leftrightarrow[[\omega]]_{x}^{1}=[[\tilde{\omega}]]_{x}^{1}
$$

\section*{Beweis:}
,, ${ }^{*}$ Ist $[[\omega]]_{x}^{1}=[[\tilde{\omega}]]_{x}^{1}$, so existiert eine Eichtransformation mit $\varepsilon(x)=1$, mit $\tilde{\omega}={ }^{\varepsilon} \omega$, also $\Omega^{\tilde{\omega}}=\operatorname{Ad}(\varepsilon) \Omega^{\omega}$ und folglich $\Omega_{x}^{\tilde{\omega}}=\Omega_{x}^{\omega}$.\\
," $\Rightarrow$ " Sei o.B.d.A. $U=\stackrel{\circ}{D}$ der offene Ball $\operatorname{im} \mathbb{R}^{n}, x=0$. Sei $e \in E_{0}$. Definiere $s$ und $\tilde{s}$ auf $U$ durch $s(y)=\tau_{\gamma_{y}}^{\omega}(e)$, bzw. $\tilde{s}(y)=\tau_{\gamma_{y}}^{\tilde{\omega}}(e)$, wobei $\gamma_{y}(t)=t y$. Die Schnitte $s$ und $\tilde{s}$ sind nicht notwendigerweise horizontal, aber $\nabla_{v} s=$ 0 und $\nabla_{v} \tilde{s}=0$ für radiale Vektoren $v$.\\
Sei $\varepsilon$ durch $\varepsilon(s)=\tilde{s}$ definiert. Dann ist $\varepsilon$ eine $x$-basierte Eichtransformation.\\
Beh: $[\tilde{\omega}]_{x}^{1}=\left[{ }^{\varepsilon} \omega\right]_{x}^{1}$. Da ${ }^{\varepsilon} \omega$ und $\tilde{\omega}$ offenbar denselben radialen Paralleltransport haben, genügt es noch zu zeigen:\\
Lemma: Seien $A_{j} \in \Omega^{1}(\stackrel{\circ}{D}, \mathfrak{g}), j=1,2$. und $F_{j}=d A_{j}+\frac{1}{2}\left[A_{j} \wedge A_{j}\right]$. Ist $F_{1,0}=F_{2,0}$ und haben $A_{1}$ und $A_{2}$ triviale radiale Parallelverschiebung, dann haben $A_{1}$ und $A_{2}$ denselben 1-Jet bei 0, d.h. ist

$$
A_{1}=\sum a_{\mu} d x_{\mu}, A_{2}=\sum \tilde{a}_{\mu} d x_{\mu} \text { mit } \tilde{a}_{\mu}, a_{\mu} \in C^{\infty}(\stackrel{\circ}{D}, \mathfrak{g})
$$

so ist $a_{\mu}(0)=\tilde{a}_{\mu}(0)$ und $\partial_{\nu} a_{\mu}(0)=\partial_{\nu} \tilde{a}_{\mu}(0)$ für alle $\nu, \mu$.

\section*{Beweis:}
Hat $A$ einen trivialen radialen Paralleltransport, so ist $A_{t v}(v)=0$, also

$$
\sum_{\mu} a_{\mu}(t v) d x_{\mu}(v)=\sum_{\mu} a_{\mu}(t v) v_{\mu}=0
$$

für alle $v$ und $t \neq 0$, also $a_{\mu}(0)=0$ aus Stetigkeitsgründen. Ebenso folgt $\tilde{a}_{\mu}(0)=0$. Ableiten von (*) nach $t$ ergibt,

$$
\sum_{\mu, \nu} \partial_{\nu} a_{\mu}(t v) v_{\nu} v_{\mu}=0
$$

also ist ( $\partial a_{\mu}$ ) schiefsymmetrisch und ebenso ( $\partial_{\nu} \tilde{a}_{\mu}$ ). Da $a_{\mu}(0)=0$ für alle $\mu$ ist also

$$
2 \partial_{\mu} a_{\nu}(0)=\frac{\partial}{\partial x_{\mu}} a_{\nu}(0)-\frac{\partial}{\partial x_{\nu}} a_{\mu}(0)=\left(F_{1,0}\right)_{\mu \nu}=\left(F_{2,0}\right)_{\mu, \nu}=2 \partial_{\mu} \tilde{a}_{\nu}(0)
$$

\subsection*{Korollar}
Durch

$$
J_{x}^{1}(Z(p)) / J_{x}^{2}(\operatorname{Aut} P) \rightarrow \Lambda^{2} T_{x}^{*} M \otimes \hat{\mathfrak{g}}, \quad[[\omega]]_{x}^{1} \mapsto F_{x}^{\omega}
$$

ist ein Isomorphismus gegeben.

\section*{Beweis:}
Die Abbildung ist wohldefiniert, da ${ }^{g} A=g A g^{-1}+g d g^{-1}$. Also ist die Wirkung von $\mathcal{G}=\Gamma$ Aut $P$ auf $J^{1} Z(P)$ durch $J^{2}$ Aut $P$ gegeben und die Abbildung ist injektiv nach dem Utiyama-Lemma.\\
Es genügt, die Surjektivität lokal zu zeigen.\\
Sei lokal $F=\sum b_{\mu \nu} d x_{\mu} \wedge d x_{\nu} \in \Omega^{2}(M, \mathfrak{g})$. Sei $A:=\sum x_{\mu} b_{\mu \nu} d x_{\mu} \in \Omega^{1}(M, \mathfrak{g})$. Dann ist $F_{0}=F_{0}^{A}=d A_{0}+\frac{1}{2}\left[A_{0} \wedge A_{0}\right]=d A_{0}$.

\section*{$11 U(1)$-Prinzipalbündel}
\subsection*{Bemerkung}
Sei $P$ ein $U(1)$-Prinzipalbündel. Dann vereinfachen sich viele Objekte, da $U(1)$ abelsch ist.

\begin{enumerate}
  \item $\operatorname{Aut}(P)=P \times_{\text {konj }} U(1)=M \times U(1),[p, g] \mapsto(\pi(p), g)$ und $\Gamma \operatorname{Aut} P= C^{\infty}(M, U(1))$.
  \item $\mathfrak{g}=i \mathbb{R}$. Sind $\omega, \tilde{\omega}$ Zusammenhangsformen, so ist $\omega-\tilde{\omega} \in \Omega^{1}(M, P \times \operatorname{Ad} \mathfrak{g})=\Omega^{1}(M, i \mathbb{R})$.
  \item Ist $\omega$ eine Zusammenhangsform, so ist die Krümmung $F^{\omega} \in \Omega^{2}(M, i \mathbb{R})$. Nach der Bianchi-Identität ist $d F^{\omega}=0$.\\
Ist $\tilde{\omega}=\omega+\eta$, so ist $F^{\tilde{\omega}}=F^{\omega}+d \eta$.
  \item $\Omega_{\text {inv }}^{k}(P, i \mathbb{R}) \ni \omega \Leftrightarrow R_{g}^{*} \omega=\operatorname{Ad}\left(g^{-1}\right) \omega=\omega$.
  \item $\Omega_{\text {inv }, \text { hor }}^{k}(P, i \mathbb{R}) \cong \Omega^{k}\left(M, P \times{ }_{\text {Ad }} i \mathbb{R}\right)=\Omega^{k}(M, i \mathbb{R})$.
\end{enumerate}

Der Isomorphismus $\Omega_{\text {inv }, \text { hor }}^{k}(P, i \mathbb{R}) \rightarrow \Omega^{k}(M, i \mathbb{R})$ ist durch $\eta \mapsto s^{*} \eta$ für einen beliebigen Schnitt $s \in \Gamma P$ gegeben, mit Inversem $\eta \mapsto \pi^{*} \eta$.

\subsection*{Notation}
$$
\begin{aligned}
\bar{\omega} & :=\frac{1}{2 \pi i} \omega \in \Omega_{\mathrm{inv}}^{1}(P) \\
\overline{F^{\omega}} & :=\frac{1}{2 \pi i} F^{\omega} \in \Omega^{2}(M)
\end{aligned}
$$

\subsection*{Definition und Satz}
$\operatorname{Durch}\left[\overline{F^{\omega}}\right] \in H_{d R}^{2}(M)=\operatorname{ker} d^{(2)} / \operatorname{bild} d^{(1)}$ ist eine wohldefinierte Kohomologieklasse unabhängig von der Wahl von $\omega$ gegeben und $c_{1}(P)=\left[\overline{F^{\omega}}\right]$ heißt die 1 . Chernklasse von $P$.

\subsection*{Korollar}
\begin{enumerate}
  \item Sind $P \cong \tilde{P}$ isomorphe $U(1)$-Prinzipalbündel, so ist $c_{1}(P)=c_{1}(\tilde{P})$.
  \item Ist $f: N \rightarrow M$, so ist $c_{1}\left(f^{*} P\right)=f^{*} c_{1}(P)$.
  \item Ist $P$ trivial, so ist $c_{1}(P)=0$.
  \item $P$ besitzt genau dann einen flachen Zusammenhang. wenn $c_{1}(P)=0$ gilt.
\end{enumerate}

\subsection*{Notation}
$$
\mathcal{F}(P):=\left\{\eta \in \Omega^{2}(M) \mid d \eta=0,[\eta]=c_{1}(P)\right\}
$$

\subsection*{Lemma}
Bezeichne $\mathcal{C}(P)$ den Raum der $U(1)$-Zusammenhänge auf $P$.\\
Die Abbildung $\psi: \mathcal{C}(P) / \mathcal{G} \rightarrow \mathcal{F}(P),[\omega] \mapsto \overline{\Omega^{\omega}}=d \bar{\omega} \in \Omega^{2}(M)$ ist wohldefiniert und surjektiv.

\section*{Beweis:}
Sei $\omega \in \mathcal{C}(P), \varepsilon \in \mathcal{G}, \Omega^{\omega}=d \omega$ und $\Omega^{\varepsilon}:=\Omega^{\varepsilon} \omega$.

\begin{enumerate}
  \item Die Abbildung $\psi$ ist wohldefiniert, denn $\Omega^{\varepsilon}=\operatorname{konj}(\varepsilon) \Omega=\Omega$.
  \item Sei $\eta \in \Omega^{2} M$ mit $\left[\overline{\Omega^{\omega}}\right]=[\eta]$ für ein $\omega \in \mathcal{C}(P)$. Dann ist $\overline{\Omega^{\omega}}-\eta=-d \alpha$ für ein $\alpha \in \Omega^{1}(M)$, also $\eta=\overline{\Omega^{\omega+2 \pi i \alpha}}$, also $[\eta]=\psi(\omega+2 \pi i \alpha)$.
\end{enumerate}

\subsection*{Notation}
$$
\begin{array}{ll}
Z^{1}(M, \mathbb{Z}) & :=\left\{\omega \in \Omega^{1}(M, i \mathbb{R}), d \omega=0, \int_{\gamma} \bar{\omega} \in \mathbb{Z} \text { für alle geschlossenen Wege }\right\} . \\
H_{d R}^{1}(M, \mathbb{Z}) & :=Z^{1}(M, \mathbb{Z}) / d \Omega^{0}(M) \\
H_{d R}^{1}(M) & :=\operatorname{ker} d^{(1)} / d \Omega^{0}(M) \\
\operatorname{Pic}(M) & :=H_{d R}^{1}(M) / H_{d R}^{1}(M, \mathbb{Z}) \\
C_{\eta}(P) & :=\left\{\omega \in \mathcal{C}(P) \mid \Omega^{\omega}=\eta\right\}
\end{array}
$$

\subsection*{Satz}
Für jedes $\eta \in \mathcal{F}(P)$ gilt:

$$
C_{\eta}(P) / \mathcal{G} \cong \operatorname{Pic}(M)
$$

\subsection*{Korollar (Satz von Weyl)}
Die Abbildung $\psi: \mathcal{C}(P) / \mathcal{G} \rightarrow \mathcal{F}(P)$ ist surjektiv und $\psi^{-1}(\eta) \cong \operatorname{Pic}(M)$. Ist insbesondere $H_{\mathrm{dR}}^{1}(M)=0$, so ist $\psi$ bijektiv.

\section*{Beweis:}
Sei $\eta \in \mathcal{F}(P)$. Für $\omega, \tilde{\omega} \in \mathcal{C}_{\eta}(P)$ ist also $d \omega=d \tilde{\omega}=2 \pi i \eta$ und $\omega-\tilde{\omega}=2 \pi i \alpha$, für ein $\alpha \in \Omega^{1}(M)$, mit $d \alpha=0$. Sei $\omega \in \mathcal{C}_{\eta}(P)$ fest gewählt.\\
Setze

$$
\varphi_{\omega}: \mathcal{C}_{\eta}(P) \rightarrow H_{d R}^{1}(M), \tilde{\omega} \mapsto\left[\frac{1}{2 \pi i}(\omega-\tilde{\omega})\right]
$$

Die Abbildung $\varphi_{\omega}$ ist surjektiv, denn für $\beta \in \operatorname{ker} d^{(1)}$ ist $\omega+2 \pi i \beta \in \mathcal{C}_{\eta}$ und $\varphi_{\omega}(\omega+2 \pi i \beta)=\beta$.\\
Behauptung Ist $\tilde{\omega}_{2}={ }^{\varepsilon} \tilde{\omega}_{1}$, so ist $\overline{\tilde{\omega}_{1}}-\overline{\tilde{\omega}_{2}} \in H^{1}(M, \mathbb{Z})$.\\
Beweis der Behauptung: Sei $\tilde{\omega}_{2}=g \tilde{\omega}_{1} g^{-1}+g^{-1} d g=\tilde{\omega}_{1}+g^{-1} d g$, für $g \in C^{\infty}(M, U(1))$ und $\gamma:[0,1] \rightarrow M$ ein geschlossener Weg. Sei $f:[0,1] \rightarrow \mathbb{R}$ mit $f(0)=f(1)-k$, mit $k \in \mathbb{Z}$ durch $g(\gamma(t))=e^{2 \pi i f(t)}$ definiert. Dann ist $g d g^{-1}=-2 \pi i f^{\prime}$ und $\int_{\gamma} g d g^{-1}=2 \pi i k$.\\
Also ist

$$
\bar{\varphi}_{\omega}: \mathcal{C}_{\eta}(P) / \mathcal{G} \rightarrow \operatorname{Pic}(M)
$$

wohldefiniert.\\
Behauptung: Die Abbildung $\bar{\varphi}_{\omega}$ ist injektiv: Seien $\left[\omega_{1}\right],\left[\omega_{2}\right] \in \mathcal{C}_{\eta}(P) / \mathcal{G}$ mit $\left[\bar{\omega}_{1}-\bar{\omega}_{2}\right] \in H_{d R}^{1}(M, \mathbb{Z})$. Zu zeigen ist:\\
Es existiert ein $g \in C^{\infty}(M, U(1))$, mit $\omega_{1}=\omega_{2}+g d g^{-1}$.\\
Sei $\alpha=\frac{1}{2 \pi i}\left(\omega_{1}-\omega_{2}\right), \psi_{\alpha}: \pi_{1}(M) \rightarrow \mathbb{Z},[\gamma] \mapsto \int_{\gamma} \alpha$. Dann existiert eine Abbildung

$$
\mu: M \rightarrow S^{1}, \text { mit } \psi_{\alpha}([\gamma])=[\mu \circ \gamma], \text { also } \oint_{\gamma} \alpha=[\mu \circ \gamma]=\frac{1}{2 \pi i} \int_{\gamma} \mu^{-1} d \mu
$$

da $S^{1}$ ein $K(\mathbb{Z}, 1)$-Raum ist, (vgl. z.B. Hatcher, Algebraic Topology, 1B, S. 90, Prop. 1B9).\\
Definiere $f \in C^{\infty}(M)$ durch

$$
f(x)=\int_{x_{0}}^{x}\left(\alpha-\frac{1}{2 \pi i} \mu^{*} \theta\right)
$$

wobei $\theta$ die Maurer-Cartanform auf $U(1)$ ist, also $\mu^{*} \theta=\mu^{-1} d \mu$. Dies ist unabhängig von der Wahl des Weges von $x_{0}$ nach $x$, da für geschlossene Wege gilt: $\int_{\gamma} \alpha=[\mu \circ \gamma]=\frac{1}{2 \pi i} \int_{\gamma} \mu^{*} \theta$. Setze also

$$
g(x)=\mu(x) e^{2 \pi i f(x)}
$$

Damit ist\\
$g^{-1} d g=\mu^{-1} e^{-2 \pi i f}\left(d \mu e^{2 \pi i f}+2 \pi i \mu d f e^{2 \pi i f}\right)=\mu^{-1} d \mu+2 \pi i d f=2 \pi i \alpha=\omega_{1}-\omega_{2}$.

\subsection*{Beispiel}
\begin{enumerate}
  \item Sei $\pi: S^{3} \rightarrow S^{2} ;\left(z_{1}, z_{2}\right) \mapsto\left[z_{1}: z_{2}\right]$ die Hopffaserung mit Zusammenhang
\end{enumerate}

$$
\begin{aligned}
\omega_{p}(Q) & =i<Q, i p>=i\left(-p_{2} q_{1}+q_{2} p_{1}-q_{3} p_{4}+q_{4} p_{3}\right)(\text { vgl. }(6.21)), \text { also } \\
\omega & =i\left(-y_{1} d x_{1}+x_{1} d y_{1}-y_{2} d x_{2}+x_{2} d y_{2}\right) \text { mit } z_{j}=x_{j}+i y_{j} \text { und } \\
\Omega^{\omega} & =d \omega=2 i\left(d x_{1} \wedge d y_{1}+d x_{2} \wedge d y_{2}\right)=-\left(d z_{1} \wedge d \bar{z}_{1}+d z_{2} \wedge d \bar{z}_{2}\right)
\end{aligned}
$$

mit $d z_{1}=d x_{1}+i d y_{1}, d \bar{z}_{1}=d x_{1}-i d y_{1}, d z_{2}=d x_{2}+i d y_{2}, d \bar{z}_{2}=d x_{2}-i d y_{2}$.\\
Sei $U=\left\{\left[z_{1}: z_{2}\right] \in \mathbb{C} P^{1} \mid z_{2} \neq 0\right\} \cong \mathbb{C},\left[z_{1}: z_{2}\right] \mapsto \frac{z_{1}}{z_{2}}$, dann ist auf $U$ ein Schnitt, $s: U \rightarrow S^{3}$ durch $\mathbb{C} \rightarrow S^{3}$

$$
s(z)=\frac{1}{1+|z|^{2}}(z, 1)
$$

gegeben und

$$
\int_{S^{2}} s^{*} \bar{\Omega}_{\omega}=-\frac{1}{2 \pi i} \int_{\mathbb{C}} \frac{1}{\left(1+|z|^{2}\right)^{2}} d z \wedge d \bar{z}=-\frac{1}{\pi} \int_{\mathbb{R}^{2}} \frac{1}{\left(1+x^{2}+y^{2}\right)^{2}} d x d y=-1
$$

(Übungsaufgabe!)\\
2. $P_{S O}\left(S^{2}\right) \rightarrow S^{2}$ ist ebenfalls $U(1)=S O(2)$-Prinzipalbündel. Es stellt sich die Frage, ob $P_{S O}\left(S^{2}\right)=S^{3}$ das Hopfbündel ist.\\
Ist $\omega$ der Levi-Civita-Zusammenhang und $\omega_{S^{2}}$ die kanonische Volumenform auf $S^{2}$, so ist $\Omega^{\omega}=i \omega_{S^{2}}$ und $\int_{S^{2}} \Omega^{\omega}=4 \pi i$, also $\frac{1}{2 \pi i} \int_{S^{2}} \Omega^{\omega}=2$, also ist $P_{S O}\left(S^{2}\right) \rightarrow S^{2}$ nicht isomorph zum Hopfbündel.

\subsection*{Satz}
Ist $M$ eine geschlossene, orientierte, 2-dimensionale Mannigfaltigkeit, so ist

$$
\int_{M} c_{1}(P) \in \mathbb{Z}
$$

Ist $\operatorname{dim} M>2$, so gilt für jede 2 -dimensionale, geschlossene, orientierte Untermannigfaltigkeit $T$ :

$$
\int_{T} c_{1}(P) \in \mathbb{Z}
$$

\section*{Beweis:}
Sei $\operatorname{dim} M=2$. Sei $P \rightarrow M$ ein $U(1)$-Prinzipalbündel, $E=P \times_{U(1)} \mathbb{C}, \sigma: M \rightarrow E$ ein Schnitt. OBdA habe $\sigma$ nur isolierte Nullstellen $N_{1}, \ldots, N_{n}$ (Transversalitätssätze aus der Differentialtopologie, Satz von Sard). Seien $D_{j}$ Koordinatenkugeln um $N_{j}$, sodass $D_{i} \cap D_{j}=\emptyset$ und $P \mid D_{j}$ trivial ist. Dann definiert $\sigma$ einen Schnitt in $P \mid\left(M \backslash\left(\cup \stackrel{\circ}{D}_{j}\right)\right)$. Seien $\sigma_{j}$ Schnitte in $P \mid D_{j}$. Dann folgt mit dem Satz von Stokes aus $\Omega^{\omega}=d^{\omega} \omega$

$$
\begin{aligned}
\int_{M} c_{1}(P) & =\int_{M \backslash D_{j}} \sigma^{*} \Omega^{\omega}+\sum \int_{D_{j}} \sigma_{j}^{*} \Omega^{\omega} \\
& =\sum_{j} \int_{\partial D_{j}}\left(\sigma_{j}^{*} \omega-\sigma^{*} \omega\right)
\end{aligned}
$$

Sei $g_{j}: \partial D_{j}=S^{1} \rightarrow S^{1}$ durch $\sigma_{j}=\sigma g_{j}$ definiert. Dann ist

$$
\begin{gathered}
\sigma_{j}^{*} \omega-\sigma^{*} \omega=g_{j}^{-1} d g_{j}=g_{j}^{*} \theta \\
2 \pi i \int_{M} c_{1}(P)=\sum_{j} \int_{\partial D_{j}} g_{j}^{*} \theta \in 2 \pi \mathbb{Z}(*)
\end{gathered}
$$

\subsection*{Satz}
Die $S^{1}$-Prinzipalbündel über $S^{2}$ sind durch $k \in \mathbb{Z}$ klassifiziert; genauer: Zu jedem $k \in \mathbb{Z}$ gibt es genau eine Isomorphieklasse von $S^{1}$-Prinzipalbündeln mit $c_{1}(P)=k$.

\section*{Beweis:}
Wir betrachten $S^{2}$ als topologischen Raum $S^{2}=D_{1} \cup D_{2} /\left(\partial D_{1} \sim \partial D_{2}\right)$ mit $D_{j} \cong D^{2}$. Sei $g: S^{1} \rightarrow S^{1}$ eine differenzierbare Abbildung und $P_{g}:= \left(D_{1} \times S^{1}\right) \sqcup\left(D_{2} \times S^{1}\right) /\left(x, e^{i \vartheta}\right) \sim\left(x, g(x) e^{i \vartheta}\right)$.\\
Behauptung: Die Isomorphieklasse von $P_{g}$ hängt nur von der Homotopieklasse von $g$ ab.\\
Seien $g \simeq \tilde{g}$, dann existiert eine differenzierbare Abbildung $H:[0,1] \times S^{1} \rightarrow S^{1}$ mit $H(0, x)=g(x), H(1, x)=\tilde{g}(x)$. Sei $Q=\left(D_{1} \times[0,1] \times S^{1}\right) \sqcup\left(D_{2} \times[0,1] \times\right.$\\
$\left.S^{1}\right) /\left(x, t, e^{i \vartheta}\right) \sim\left(x, t, H(t, x) e^{i \vartheta}\right)$. Dann ist $Q \rightarrow S^{2} \times[0,1]$ ein (berandetes) $S^{1}$-Prinzipalbündel.\\
Sei $i_{j}: S^{2} \rightarrow S^{2} \times[0,1], x \mapsto(x, j)$. Dann ist $i_{0}^{*} Q=P_{g}$, und $i_{1}^{*} Q=P_{\tilde{g}}$, also ist $P_{g} \cong P_{\tilde{g}}$, da $i_{0}$ und $i_{1}$ homotop sind (vgl. 1.23).\\
Sei $P(k)$ das $S^{1}$-Prinzipalbündel über $S^{2}$ zu $z \rightarrow z^{k}$, so ist $\int_{S^{2}} c_{1}(P(k))=k$ nach (*).

\subsection*{Beispiel}
$$
\int_{S^{2}} c_{1}\left(P_{S O}\left(S^{1}\right)\right)=2
$$

(vgl. Quantisierung von Dirac-Monopolen (1975) von Wu, Yang. Gesucht: $\nabla$ mit $F^{\nabla}=c \cdot \operatorname{vol}(M), A_{ \pm}=i c(1 \mp \cos \vartheta) d \varphi$, auf $S^{2}$.)

\section*{Charakteristische Klassen für Prinzipalbündel (Chern-Weil-Theorie)}
Wir wollen nun die Ergebnisse aus Kapitel 11 auf beliebige $G$-Prinzipalbündel verallgemeinern. Dazu betrachten wir geeignete Abbildunge $\mathfrak{g} \rightarrow \mathbb{C}$. Wir betrachten in diesem Abschnitt stets $\mathbb{C}$-wertige lineare Abbildungen und Differentialformen.

\subsection*{Definitionen und Notationen}
\begin{enumerate}
  \item Ein homogenes Polynom vom Grad $k$ auf $\mathfrak{g}$ ist eine Abbildung $p: \mathfrak{g} \rightarrow \mathbb{C}$ mit $p(X)=\sum_{i_{1} \leq \ldots \leq i_{k}} p_{i_{1} \ldots i_{k}} x_{i_{1}} \cdots x_{i_{k}}$, wobei $X=\sum x_{j} e_{j}$ bezüglich einer Basis $e_{1}, \ldots, e_{n}$ von $\mathfrak{g}$. Es bezeichne $P^{k}(\mathfrak{g})$ den Vektorraum der homogenen Polynome vom Grad $k$. Ein homogenes Polynom $p$ heißt invariant, $p \in P_{\text {inv }}^{k}(\mathfrak{g})$, falls $p(\operatorname{Ad}(g) X)=p(X)$ gilt.
  \item $\operatorname{Sym}_{\text {inv }}^{k}(\mathfrak{g}):=\left\{f \in \operatorname{Sym}^{k}(\mathfrak{g}) \mid f\left(\operatorname{Ad}(g) X_{1}, \ldots, \operatorname{Ad}(g) X_{k}\right)=f\left(X_{1}, \ldots, X_{k}\right)\right\}$. $\operatorname{Sym}_{\text {inv }}(\mathfrak{g})=\oplus \operatorname{Sym}_{\text {inv }}^{k} \mathfrak{g}$ ist in kanonischer Weise ein Ring:
\end{enumerate}

$$
\left.\left.(g h)\left(X_{1}, \ldots, X_{k+r}\right)=\frac{1}{(k+r)!} \sum_{r \in S} g\left(X_{\tau(1)}, \ldots, X_{\tau(k)}\right)\right) \cdot h\left(X_{\tau(k+1)}, \ldots, X_{\tau(k+r)}\right)\right)
$$

\subsection*{Lemma}
\begin{enumerate}
  \item Durch $\operatorname{Sym}(\mathfrak{g}) \rightarrow P(\mathfrak{g}), f \mapsto \varphi(f): X \mapsto f(X, \ldots, X)$ ist ein graduierter Algebrenisomorphismus gegeben. Die Umkehrung $P^{k} \mathfrak{g} \rightarrow \operatorname{Sym}^{k}(\mathfrak{g})$ ist durch $f \mapsto f_{(1,1,1, \ldots, 1)}$ gegeben, wobei $f_{(1 \ldots 1)}$ durch
\end{enumerate}

$$
\frac{1}{k!} f\left(t_{1} X_{1}+\ldots+t_{k} X_{k}\right)=\sum_{\alpha \in \mathbb{N}_{0}^{k}} t^{\alpha} f_{\alpha}\left(X_{1}, \ldots, X_{k}\right)
$$

definiert ist (Polarisationsformel).\\
2. Offenbar induziert die Abbildung in 1. einen Isomorphismus

$$
\operatorname{Sym}_{\text {inv }}^{k} \mathfrak{g} \rightarrow P_{\text {inv }}^{k} \mathfrak{g} .
$$

\subsection*{Beispiel}
Sei $F(x)=\operatorname{tr}\left(X^{3}\right)$, so ist die zugehörige symmetrische Abbildung durch

$$
f\left(X_{1}, X_{2}, X_{3}\right)=\frac{1}{2} \operatorname{tr}\left(X_{1} X_{2} X_{3}+X_{2} X_{1} X_{3}\right)
$$

gegeben, denn $\operatorname{tr}\left(\left(t_{1} X_{1}+t_{2} X_{2}+t_{3} X_{3}\right)^{3}\right)=t_{1} t_{2} t_{3} \operatorname{tr}\left(X_{1} X_{2} X_{3}+X_{1} X_{3} X_{2}+\right.$ zykl. Perm. $)+\sum_{\alpha \neq(1, \ldots, 1)} t^{\alpha} \ldots$

\subsection*{Bemerkung}
Ist $\mathfrak{g} \subseteq M(n \times n, k)$, so ist $\operatorname{Ad}(g) X=g X g^{-1}$ und $p \in P_{\text {inv }}(\mathfrak{g})$ genau dann, wenn $p(X Y)=p(Y X)$ für alle $X \in G, Y \in \mathfrak{g}$. Für $\mathfrak{g}=M(n \times n, k)$ gilt dies aus Stetigkeitsgründen dann für alle $X, Y \in \mathfrak{g}$.

\subsection*{Definition}
Ist $\eta_{j} \in \Omega_{\text {inv }}^{l_{j}}(P, \mathfrak{g}), f \in \operatorname{Sym}_{\text {inv }}^{k}(\mathfrak{g})$, so ist

$$
f\left(\eta_{1} \wedge \ldots \wedge \eta_{k}\right):=f\left(\eta_{1}, \ldots, \eta_{k}\right) \in \Omega_{\mathrm{inv}}^{l k}(P, \mathbb{C})
$$

mit $l_{1}+\ldots+l_{k}=: l$ durch

$$
\begin{gathered}
f\left(\eta_{1} \wedge \ldots \wedge \eta_{k}\right)\left(X_{1}, \ldots, X_{l k}\right)= \\
=\frac{1}{l_{1}!\ldots l_{k}!} \sum_{\tau} \operatorname{sign}(\tau) f\left(\eta_{1}\left(X_{\tau(1)}, \ldots, X_{\tau(l)}\right), \ldots, \eta_{k}\left(X_{\tau(k(l-1)+1)}, \ldots, X_{\tau(k l)}\right)\right)
\end{gathered}
$$

gegeben. Ist $\eta_{j} \in \Omega_{\text {inv }, \text { hor }}^{l}(P, \mathfrak{g})$, so ist $f\left(\eta_{1} \wedge \ldots \wedge \eta_{k}\right) \in \Omega^{l k}(M, \mathbb{C}), \operatorname{denn} f\left(\sigma^{*} \eta_{1} \wedge\right. \left.\ldots \wedge \sigma^{*} \eta_{k}\right)$ ist unabhängig von der Wahl des Schnittes $\sigma$ wohldefiniert.

\subsection*{Beispiel}
Sei $\operatorname{det}\left(t 11-\frac{1}{2 \pi i} X\right)=: \sum_{k=0}^{n} c_{k}(X) t^{n-k}$. Dann sind $c_{k}$ invariante Polynome

$$
\begin{aligned}
& c_{1}(X)=\operatorname{tr}(X)\left(-\frac{1}{2 \pi i}\right) \\
& c_{n}(X)=\operatorname{det}(X)\left(-\frac{1}{2 \pi i}\right)^{n}
\end{aligned}
$$

Ist $X$ diagonalisierbar, so ist $c_{k}(X)=\left(-\frac{1}{2 \pi i}\right)^{k} \sigma_{k}\left(\lambda_{1}, \ldots, \lambda_{n}\right)$, wobei $\lambda_{j}$ die Eigenwerte und $\sigma_{k}$ die elementarsymmetrischen Polynome vom Grad $k$ sind.

\subsection*{Bemerkung}
Die Polynome $\sigma_{k}$ sind durch $\left(t+X_{1}\right) \ldots\left(t+X_{n}\right)=\sum_{k=0}^{n} t^{k} \sigma_{n-k}\left(x_{1}, \ldots, x_{n}\right)$ definiert.

\subsection*{Bemerkung}
Ist $\Omega=:\left(\Omega_{i j}\right)_{i j} \in \Omega_{\text {inv }, \text { hor }}^{2}(P, \mathfrak{g}), \mathfrak{g} \subseteq g l(n, \mathbb{C})$, so ist\\
$c_{1}(\Omega)=-\frac{1}{2 \pi i} \operatorname{tr}(\Omega)=-\frac{1}{2 \pi i}\left(\Omega_{11}+\ldots+\Omega_{n n}\right) \in \Omega_{\mathrm{inv}, \mathrm{hor}}^{2}(P)=\Omega^{2}(M, \mathbb{C})$\\
$c_{n}(\Omega)=\left(-\frac{1}{2 \pi i}\right)^{n} \operatorname{det} \Omega=\left(-\frac{1}{2 \pi i}\right)^{n} \sum_{\sigma \in(n)} \operatorname{sign}(\sigma) \Omega_{1 \sigma(1)} \wedge \ldots \wedge \Omega_{n \sigma(n)} \in \Omega^{2 n}(M, \mathbb{C})$

\subsection*{Bemerkung}
Ist $f \in P_{\mathrm{inv}}^{k}(\mathfrak{g}), \eta_{j} \in \Omega_{\mathrm{inv}}^{l_{j}}(P)$, so ist

$$
d\left(f\left(\eta_{1} \wedge \ldots \wedge \eta_{k}\right)\right)=\sum_{j=1}^{k}(-1)^{l_{1}+\ldots+l_{j-1}} f\left(\eta_{1} \wedge \ldots \wedge d \eta_{j} \wedge \ldots \wedge \eta_{k}\right)
$$

Ist $\eta_{j} \in \Omega_{\text {inv }, \text { hor }}^{l_{j}}(P)$ und also

$$
f\left(\eta_{1} \wedge \ldots \wedge \eta_{k}\right):=f\left(\sigma^{*} \eta_{1} \wedge \ldots \wedge \sigma^{*} \eta_{k}\right) \in \Omega(M, \mathbb{C})
$$

wohldefiniert, so ist nach 6.19

$$
d\left(f\left(\eta_{1} \wedge \ldots \wedge \eta_{k}\right)\right)=\sum_{j=1}^{k}(-1)^{l_{1}+\ldots+l_{j-1}} f\left(\eta_{1} \wedge \ldots \wedge D^{H} \eta_{j} \wedge \ldots \wedge \eta_{k}\right)
$$

für einen beliebigen Zusammenhang $H$ in $P$.

\subsection*{Lemma}
Ist $\Omega^{\omega}$ die Krümmung eines $G$-Zusammenhangs auf $P, f \in P_{\text {inv }}^{k}(\mathfrak{g})$, so gilt

\begin{enumerate}
  \item $f\left(\Omega^{\omega}\right) \in \Omega^{2 k}(M, \mathbb{C})$.
  \item $d\left(f\left(\Omega^{\omega}\right)\right)=0$
  \item Ist $\omega^{\prime}$ ein weiterer Zusamenhang, so existiert ein $\eta \in \Omega^{1} M$ mit
\end{enumerate}

$$
f\left(\Omega^{\omega}\right)-f\left(\Omega^{\omega^{\prime}}\right)=d \eta
$$

\section*{Beweis:}
\begin{enumerate}
  \item $\sqrt{ }$
  \item Es gilt
\end{enumerate}

$$
\begin{aligned}
d f\left(\Omega^{\omega}\right) & =d\left(f\left(\Omega^{\omega} \wedge \ldots \wedge \Omega^{\omega}\right)\right. \\
& =f\left(d^{\omega} \Omega^{\omega} \wedge \ldots \wedge \Omega^{\omega}\right)+\ldots+f\left(\Omega^{\omega} \wedge \ldots \wedge d^{\omega} \Omega^{\omega}\right)=0
\end{aligned}
$$

\begin{enumerate}
  \setcounter{enumi}{2}
  \item Seien $\omega_{0}$ und $\omega_{1}$ Zusammenhangsformen, $\alpha=\omega_{1}-\omega_{0}$ und $\omega_{t}=\omega_{0}+t \alpha$. Sei $\Omega^{t}:=\Omega^{\omega_{t}}$, also $\Omega^{t}=\Omega^{0}+t d \alpha+t\left[\omega_{0} \wedge \alpha\right]+\frac{1}{2} t^{2}[\alpha \wedge \alpha]$. Dann ist
\end{enumerate}

$$
\frac{d}{d t} \Omega^{t}=d \alpha+\left[\omega_{0} \wedge \alpha\right]+t[\alpha \wedge \alpha]=d^{\omega_{t}} \alpha
$$

Sei $\hat{f}:[0,1] \rightarrow \Omega^{2 k}(M), \hat{f}(t)=f\left(\Omega^{t}\right)$. Dann ist

$$
\begin{aligned}
f\left(\Omega^{1}\right)-f\left(\Omega^{0}\right)=\int_{0}^{1} \frac{d}{d t} \hat{f}(t) & =k \int_{0}^{1} f\left(\frac{d}{d t} \Omega^{t} \wedge \ldots \wedge \Omega^{t}\right) d t \\
& =k \int_{0}^{1} f\left(d^{\omega_{t}} \alpha \wedge \Omega^{t} \wedge \ldots \wedge \Omega^{t}\right) d t \\
& =k \int_{0}^{1} d f\left(\alpha \wedge \ldots \wedge \Omega^{t}\right) d t
\end{aligned}
$$

da $d^{\omega_{t}} \Omega^{t}=0$ gilt nach der Bianchi-Indentität und $\alpha \in \Omega_{\text {inv }, \text { hor }}^{1}(P, \mathfrak{g})$ und $\Omega^{t} \in \Omega_{\text {inv }, \text { hor }}^{2}(P, \mathfrak{g})$. Sei

$$
\eta=k \int_{0}^{1} f\left(\alpha \wedge \Omega^{t} \wedge \ldots \wedge \Omega^{t}\right) d t
$$

Dann ist $d \eta=f\left(\Omega^{1}\right)-f\left(\Omega^{0}\right)$.

\subsection*{Korollar und Definition}
Sei $P \rightarrow M$ ein $G$-Prinzipalbündel

$$
\begin{aligned}
\mathcal{W}^{P}: \operatorname{Sym}_{\text {inv }}(\mathfrak{g}) & \rightarrow H_{d R}(M, \mathbb{C}) \\
f & \mapsto\left[f\left(\Omega^{\omega}\right)\right]
\end{aligned}
$$

ist ein wohldefinierter Algebrenhomomorphismus. $\mathcal{W}^{P}$ heißt Weil-Homomorphismus.

\section*{Beweis:}
Es ist nur

$$
(f \cdot g)\left(\Omega^{\omega}\right)=f\left(\Omega^{\omega}\right) \wedge g\left(\Omega^{\omega}\right)
$$

nachzurechnen.

\subsection*{Beispiel}
\begin{enumerate}
  \item $\mathcal{W}^{M \times G}(f)=0$ für jedes $f \in \operatorname{Sym}_{\text {inv }}(\mathfrak{g})$.
  \item Ist $G=S^{1}$ und $f(X)=X$, so ist $\mathcal{W}^{P}(f)=c_{1}(P)$.
\end{enumerate}

\subsection*{Lemma}
\begin{enumerate}
  \item Ist $P \cong \tilde{P}$, so ist $\mathcal{W}^{P}=\mathcal{W}^{\tilde{P}}$.
  \item Ist $\varphi: N \rightarrow M$ ein Diffeomorphismus, so ist $\mathcal{W}^{\varphi^{*}} P(f)=\varphi^{*} \mathcal{W}^{P}(f)$ für $f \in \operatorname{Sym}_{\text {inv }}(\mathfrak{g})$.
\end{enumerate}

\section*{Beweis:}
\begin{enumerate}
  \item Ist $\varphi: \tilde{P} \rightarrow P$ ein $G$-Prinzipalbündelisomorphismus, $\omega$ ein Zusammenhang auf $P$, so ist $\varphi^{*} \omega$ ein Zusammenhang auf $\tilde{P}$ mit Krümmung $\varphi^{*} \Omega^{\omega}$. Ist $\tilde{\sigma}$ ein Schnitt in $\tilde{P}$, so ist $\varphi \circ \tilde{\sigma}=: \sigma$ einer in $P$ und $\sigma^{*} \Omega^{\omega}=\tilde{\sigma}^{*} \varphi^{*} \Omega^{\omega}$.
  \item Ist $P \rightarrow M$ ein $G$-Prinzipalbündel mit Zusammenhang $\omega$ und $\hat{\varphi}: \varphi^{*} P \rightarrow P$ die kanonische Abbildung über $\varphi$, so ist $\hat{\varphi}^{*} \omega$ ein Zusammenhang in $\varphi^{*} P$ mit Krümmung $\hat{\varphi}^{*} \Omega^{\omega}$. Ein lokaler Schnitt $s: U \rightarrow P$ definiert einen lokalen Schnitt $\tilde{s}: \varphi^{-1}(U) \rightarrow \varphi^{*} P$ für den $\hat{\varphi} \circ \tilde{s}=s \circ \varphi$ gilt, also $\tilde{s}^{*} \hat{\varphi}^{*} \Omega^{\omega}=\varphi^{*} s^{*} \Omega^{\omega}$.
\end{enumerate}

\subsection*{Lemma}
Ist $\lambda: H \rightarrow G$ ein Liegruppenhomomorphismus, $P \rightarrow M$ ein $G$-Prinzipalbündel und $\tilde{\lambda}: Q \rightarrow P$ eine $H$-Struktur von $P$, so ist $\mathcal{W}^{Q}\left(\lambda^{*} f\right)=\mathcal{W}^{P}(f)$, mit $\lambda^{*} f=\left.f \circ d \lambda\right|_{1}$. Ist $Q \rightarrow P$ eine Reduktion, so ist $\omega^{Q}(f \mid \mathfrak{h})=\omega^{p}(f)$.

\section*{Beweis:}
Ein $H$-Zusammenhang $\omega$ auf $Q$ definiert einen $G$-Zusammenhang $\tilde{\omega}$ auf $P= Q \times_{\lambda} G$, für den $\tilde{\lambda}^{*} \tilde{\omega}=\left.d \lambda\right|_{1} \circ \omega$ gilt, also

$$
d \lambda_{1} \circ \Omega^{\omega}=\tilde{\lambda}^{*} \Omega^{\tilde{\omega}}
$$

Ein Schnitt $s$ in $Q$ induziert einen Schnitt $\tilde{s}=\tilde{\lambda} \circ s$ in $P$ und es gilt

$$
\tilde{s}^{*} \Omega^{\tilde{\omega}}=s^{*} d \lambda_{1} \circ \Omega^{\omega}=d \lambda_{1} \circ s^{*} \Omega^{\omega}
$$

Für $f \in \operatorname{Sym}_{\text {inv }}(\mathfrak{g})$ ist also\\
$f\left(\tilde{s}^{*} \Omega^{\tilde{\omega}} \wedge \ldots \wedge \tilde{s}^{*} \Omega^{\tilde{\omega}}\right)=f\left(d \lambda_{1} \circ s^{*} \Omega^{\omega} \wedge \ldots \wedge d \lambda_{1} \circ s^{*} \Omega^{\omega}\right)=\lambda^{*} f\left(s^{*} \Omega^{\omega} \wedge \ldots \wedge s^{*} \Omega^{\omega}\right)$.

\subsection*{Definition}
Sei $c_{j}$ wie in 12.6. Dann heißt $\mathcal{W}^{P}\left(c_{j}\right)$ die $j$-te Chern-Klasse von $P$. Man schreibt $\mathcal{W}^{P}\left(c_{j}\right)=: c_{j}(P)$.

\section*{Chern-Klassen für komplexe Vektorraumbündel}
Jedes komplexe Vektorraumbündel besitzt eine komplexe Bündelmetrik und damit eine Reduktion auf $U(n)$. Nach 12.13 gilt für jede Reduktion $Q$ von $P$

$$
\mathcal{W}^{Q}\left(\left.f\right|_{\mathfrak{h}}\right)=\mathcal{W}^{P}(f)
$$

Damit können wir die Chern-Klasse von komplexen Vektorraumbündeln definieren.

\subsection*{Definition und Satz}
Sei $E$ ein komplexes Vektorraumbündel, $P$ das zugrundeliegende $\operatorname{GL}(n, \mathbb{C})$ Prinzipalbündel oder eine $U(n)$-Reduktion, so ist die $k$-te Chern-Klasse von $E$ durch

$$
c_{k}(E)=c_{k}(P) \in H^{2 k}(M, \mathbb{R})
$$

wohldefiniert. $c(E)=c_{0}(E)+c_{1}(E)+\ldots+c_{n}(E)=\operatorname{det}\left(\mathbb{l}-\frac{1}{2 \pi i} \Omega\right)$ heißt totale Chern-Klasse.\\
Ist $M$ eine $n$-dimensionale, geschlossene, orientierte Mannigfaltigkeit, $2\left(i_{1}+\right. \left.{ }_{\ldots}+i_{k}\right)=n$ und $\left[\alpha_{j}\right] \in c_{i_{j}}(E)$, so heißt

$$
\int_{M} \alpha_{1} \wedge \ldots \wedge \alpha_{k}
$$

die $\left(i_{1}, \ldots, i_{k}\right)$-te Chernzahl (oder topologische Quantenzahl) von $E$.\\
Ist $M$ eine komplexe Mannigfaltigkeit, so bezeichnet man als $\left(i_{1}, \ldots, i_{k}\right)$-te Chernzahl von $M$ die $\left(i_{1}, . ., i_{k}\right)$-te Chernzahl von $T M$.

\section*{Beweis:}
Ist $A \in u(n)$, also $\bar{A}=-A^{t}$, also

$$
\begin{aligned}
\sum_{k=0}^{n} \overline{c_{k}(A)} t^{n-k} & =\operatorname{det} \overline{\left(t \cdot 11-\frac{1}{2 \pi i} A\right)} \\
& =\operatorname{det}\left(t \cdot 11+\frac{1}{2 \pi i} \bar{A}\right) \\
& =\operatorname{det}\left(t \cdot 11-\frac{1}{2 \pi i} A^{t}\right) \\
& =\operatorname{det}\left(t \cdot 11-\frac{1}{2 \pi i} A\right) \\
& =\sum_{k=0}^{n} c_{k}(A) t^{n-1}
\end{aligned}
$$

also ist $c_{k}(A) \in \mathbb{R}$.

\subsection*{Beispiel}
$$
\begin{aligned}
c_{k}(E) & =\frac{(-1)^{k}}{(2 \pi i)^{k}} \sum_{i_{1}<\ldots<i_{k}} \sum_{\sigma \in S(k)} \operatorname{sign}(\sigma) \Omega_{i_{1} i_{\sigma(1)}} \wedge \ldots \wedge \Omega_{i_{k} i_{\sigma(k)}} \\
c_{0}(E) & =1 \\
c_{1}(E) & =-\frac{1}{2 \pi i} \operatorname{tr} \Omega \\
c_{2}(E) & =-\frac{1}{4 \pi^{2}}\left(\sum_{i<j} \Omega_{i i} \wedge \Omega_{j j}-\Omega_{i j} \wedge \Omega_{j i}\right)=\frac{1}{8 \pi^{2}}(\operatorname{tr}(\Omega \wedge \Omega)-\operatorname{tr} \Omega \wedge \operatorname{tr} \Omega) \\
c_{n}(E) & =\left(-\frac{1}{2 \pi i}\right)^{n} \operatorname{det} \Omega
\end{aligned}
$$

\subsection*{Korollar}
Besitzt $E$ eine Reduktion auf $S U(n)$, so ist $c_{1}(E)=0$ und $c_{2}(E)=\frac{1}{8 \pi^{2}}\left[\operatorname{tr}\left(F^{\nabla} \wedge\right.\right. \left.\left.F^{\nabla}\right)\right]$, denn für $A \in s u(2)$ ist $\operatorname{Spur} A=0$.

\subsection*{Lemma}
\begin{enumerate}
  \item Ist $E \rightarrow M$ ein $k$-dimensionales Vektorraumbündel $c_{k}(E) \in H_{d R}^{2 k}(M)$, so ist
\end{enumerate}

$$
\begin{aligned}
c_{0}(E) & =1 \\
c_{l}(E) & =0 \text { für } l>k
\end{aligned}
$$

\begin{enumerate}
  \setcounter{enumi}{1}
  \item Ist $f: N \rightarrow M$ differenzierbar, so ist $f^{*} c(E)=c\left(f^{*} E\right)$.
  \item $c\left(E_{1} \oplus E_{2}\right)=c\left(E_{1}\right) \cdot c\left(E_{2}\right)$, genauer
\end{enumerate}

$$
\sum_{j=0}^{k} c_{j}\left(E_{1}\right) c_{k-j}\left(E_{2}\right)=c_{k}\left(E_{1} \oplus E_{2}\right)
$$

\begin{enumerate}
  \setcounter{enumi}{3}
  \item Sei $l_{1}=\left\{(L, \xi) \in \mathbb{C} P^{1} \times \mathbb{C}^{2} \mid \xi \in L\right\} \rightarrow \mathbb{C} P^{1}$ das kanonische Linienbündel, $S^{3} \rightarrow S^{2}$ das Hopfbündel, so ist
\end{enumerate}

$$
S^{3} \times_{S^{1}} \mathbb{C} \cong l_{1},\left[\left(\omega_{1}, \omega_{2}\right), z\right] \mapsto\left(\left[\omega_{1}, \omega_{2}\right],\left(\omega_{1} z, \omega_{2} z\right)\right)
$$

also ist $c_{1}\left(l_{1}\right)=-1$.

\section*{Beweis:}
(von 3.) Sei $\omega_{j}$ ein Zusammenhang auf $E_{j}$ mit Krümmung $\Omega^{j}$, dann ist kanonisch $\omega_{1} \oplus \omega_{2}$ ein Zusammenhang auf $E_{1} \oplus E_{2}$ mit Krümmung $\Omega^{1} \oplus \Omega^{2}$, also\\
ist

$$
\begin{aligned}
\sum_{k=1}^{n} c_{k}(E) t^{n-k} & =\left(\sum_{k_{1}=1}^{n_{1}} c_{k_{1}}\left(E_{1}\right) t^{n_{1}-k_{1}}\right)\left(\sum_{k_{2}=1}^{n_{2}} c_{k_{2}}\left(E_{2}\right) t^{n_{2}-k_{2}}\right) \\
& =\sum_{k=1}^{n}\left(\sum_{j=1}^{k} c_{j}\left(E_{1}\right) c_{k-j}\left(E_{2}\right)\right) t^{n-k}
\end{aligned}
$$

wobei $E_{j}$ ein $n_{j}$-dimensionales Vektorraumbündel und $n_{1}+n_{2}=n$ ist.

\subsection*{Bemerkung}
Es gilt sogar, dass die Chernklassen im Bild von

$$
H_{\mathrm{sing}}^{*}(M, \mathbb{Z}) \rightarrow H_{d R}(M, \mathbb{R})
$$

liegen, d.h., es gilt

$$
\int_{\sigma} c_{k}(E) \in \mathbb{Z}
$$

für alle glatten $k$-dimensionalen geschlossenen Untermannigfaltigkeiten von $M$.\\
Durch 1.-4. im Lemma wird axiomatisch eine Klasse $c^{*}(E) \in H_{\text {sing }}(M, \mathbb{Z})$ definiert (ersetze $(4)$ durch $\left(4^{\prime}\right): c_{1}\left(l_{1}\right)$ erzeugt $H_{\text {sing }}^{2}\left(\mathbb{C} P^{1}, \mathbb{Z}\right)$.\\
Weitere Eigenschaften der Chern-Klassen:

\subsection*{Satz}
\begin{enumerate}
  \item Ist $E_{1} \cong E_{2}$, so gilt $c\left(E_{1}\right)=c\left(E_{2}\right)$.
  \item Ist $E$ trivial, so ist $c(E)=1$, also $c_{k}(E)=0$ für $k>0$.
  \item Ist $E$ ein komplexes Vektorraumbündel der Dimension $n$, das $k$ überall linear unabhängige Schnitte besitzt, so ist $c_{j}(E)=0$ für $j>n-k$.
  \item Es gilt
\end{enumerate}

$$
c_{k}\left(E^{*}\right)=(-1)^{k} c_{k}(E)
$$

\section*{Beweis:}
\begin{enumerate}
  \item $\sqrt{ }$ 2. $\sqrt{ }$
  \item Existieren $k$ überall linear unabhängige Schnitte, so ist $E \cong\left(M \times \mathbb{R}^{k}\right) \oplus F$, wobei $F \rightarrow M$ ein $(n-k)$-dimensionales Vektorraumbündel ist, also ist $c(E)= 1 \cdot c(F)$.
  \item Für $\eta \in \Gamma E^{*}$ gilt $\left(\nabla_{X}^{*} \eta\right)(e)=X \eta(e)-\eta\left(\nabla_{X} e\right)(Y)$ und folglich
\end{enumerate}

$$
F^{\nabla^{*}}(X, Y) \eta=-\eta\left(F^{\nabla}(X, Y) \cdot\right)
$$

also $\Omega^{*}=-\Omega$.

\subsection*{Beispiele}
\begin{enumerate}
  \item Sei $\nu=S^{n} \times \mathbb{R}$ das Normalenbündel, dann ist $T S^{n} \oplus \nu=S^{n} \times \mathbb{R}^{n+1}$. Sei für ein reelles Bünde $E$ das komplexifizierte durch $E^{\mathbb{C}}=E \otimes \mathbb{C}$ gegeben. Dann ist $T^{\mathbb{C}} S^{n} \oplus \nu^{\mathbb{C}}=S^{n} \times \mathbb{C}^{n+1}$, also ist $c\left(T^{\mathbb{C}} S^{n}\right)=1$.
  \item Sei $E \rightarrow M$ ein komplexes Vektorraumbündel der Dimension $n$, dann ist $c_{1}\left(\Lambda^{n} E\right)=c_{1}(E)$, denn für die kovariante Ableitung $\hat{\nabla}$ auf $\Lambda^{1} E$ gilt $\hat{\nabla}\left(e_{1} \wedge \ldots \wedge e_{n}\right):=\sum e_{1} \wedge \ldots \wedge \nabla e_{j} \wedge \ldots \wedge e_{j}$, also
\end{enumerate}

$$
\begin{aligned}
F^{\hat{\nabla}}\left(e_{1} \wedge \ldots \wedge e_{n}\right) & =\sum_{j} e_{1} \wedge \ldots F^{\nabla} e_{j} \wedge e_{n} \\
& =\sum_{j, k} e_{1} \wedge \ldots<F^{\nabla} e_{j}, e_{k}>e_{k} \wedge \ldots \wedge e_{n} \\
& =\sum_{j}<F^{\nabla} e_{j}, e_{j}>e_{1} \wedge \ldots \wedge e_{n} \\
& =\operatorname{Spur}\left(F^{\nabla}\right) e_{1} \wedge \ldots \wedge e_{n}
\end{aligned}
$$

also ist $c_{1}\left(\Lambda^{n} E\right)=-\frac{1}{2 \pi i} \operatorname{tr}\left(F^{\hat{\nabla}}\right)=c_{1}(E)$.\\
3. Als letztes Beispiel berechnen wir die Chern-Klassen für $\mathbb{C} P^{n}$. Der projektive Raum $\mathbb{C} P^{n}$ ist eine komplexe Mannigfaltigkeit, also ist das Tangentialbündel ein komplexes Vektorraumbündel. Dazu betrachten wir zuerst $\pi: S^{2 n+1} \rightarrow \mathbb{C} P^{n}$ die verallgemeinerte Hopffaserung. Sei ( $\cdot, \cdot$ ) das kanonische hermitsche Skalarprodukt auf $\mathbb{C}^{n+1}$. Dann ist

$$
\begin{aligned}
T_{z} S^{2 n+1} & =\left\{x \in \mathbb{C}^{n+1} \mid \operatorname{Re}(x, z)=0\right\} \\
T_{z}^{\prime} S^{2 n+1} & =\left\{X \in \mathbb{C}^{n+1} \mid X \in \mathbb{R} i z\right\}
\end{aligned}
$$

Also ist durch $\pi^{*} l_{n}^{\perp}=: \tilde{l}_{n}^{\perp} \subseteq T_{z} S^{2 n+1}$ mit

$$
l_{n,[z]}^{\perp}=\left\{x \in \mathbb{C}^{n+1} \mid(x, z)=0\right\}
$$

ein $S^{1}$-Zusammenhang auf $S^{2 n+1} \rightarrow \mathbb{C} P^{n}$ gegeben.\\
Sei $l_{n}=\left\{(w, L) \in \mathbb{C} P^{n} \times \mathbb{C}^{n+1} \mid L \in w\right\}$ das kanonische Geradenbündel über $\mathbb{C} P^{n}$ und $\tilde{l}_{n}=\pi^{*} l_{n}$. Dann ist für $z \in S^{2 n+1}$ ein Isomorphismus

$$
\operatorname{Hom}\left(\tilde{l}_{n}, \tilde{l}_{n}^{\perp}\right)_{z} \rightarrow T_{[z]} \mathbb{C} P^{n}, \psi \mapsto d \pi_{z}((\psi(\mathbb{C} z))
$$

wohldefiniert.\\
Damit ist

$$
\begin{aligned}
T \mathbb{C} P^{n} \oplus\left(\mathbb{C} P^{n} \times \mathbb{C}\right) & \cong \operatorname{Hom}\left(l_{n}, l_{n}^{\perp}\right) \oplus \operatorname{Hom}\left(l_{n}, l_{n}\right) \\
& \cong \operatorname{Hom}\left(l_{n}, \mathbb{C} P^{n} \times \mathbb{C}^{n+1}\right) \\
& \cong l_{n}^{*} \oplus \ldots \oplus l_{n}^{*}
\end{aligned}
$$

also

$$
c\left(T \mathbb{C} P^{n}\right)=\left(1-c_{1}\left(l_{n}\right)\right)^{n+1}
$$

Wir kommen jetzt noch zu den Chern-Klassen reeller Mannigfaltigkeiten, deren Tangentialbündel fast komplex ist.

\subsection*{Definition}
\begin{enumerate}
  \item Sei $M$ eine reelle Mannigfaltigkeit. Ein Vektorraumbündelhomomorphismus heißt fast komplexe Struktur auf $M$, wenn $J^{2}=-J d_{T M}$ gilt. $(M, J)$ heißt eine fast komplexe Mannigfaltigkeit.
  \item Ist ( $M, g$ ) eine Riemannsche Mannigfaltigkeit und $J$ eine orthogonale, fast komplexe Struktur auf $M$, so heißt ( $M, g, J$ ) fast hermitsche Mannigfaltigkeit.
\end{enumerate}

\subsection*{Bemerkung}
\begin{enumerate}
  \item Eine fast hermitsche Mannigfaltigkeit hat immer gerade Dimension $2 m$, und das Tangentialbündel hat eine Reduktion auf $U(m)$, wobei $U(m) \subseteq S O(2 m)$ durch $f: A+i B \rightarrow\left(\begin{array}{cc}A & -B \\ B & A\end{array}\right)$ aufgefasst wird. Umgekehrt besitzt eine geradedimensionale Mannigfaltigkeit genau dann eine orthogonale fast komplexe Struktur, wenn das Tangentialbündel auf $U(m)$ reduziert werden kann.
  \item Eine fast komplexe Mannigfaltigkeit ist eine komplexe Mannigfaltigkeit, falls
\end{enumerate}

$$
[X, Y]=[J X, J Y]-J([J X, Y]-J[X, J Y]) .
$$

\begin{enumerate}
  \setcounter{enumi}{2}
  \item Ist ( $M^{2 k}, J$ ) eine fast komplexe Mannigfaltigkeit, so ist durch $J_{x}^{*}$ eine lineare Abbildung auf $\Lambda^{k} T^{*} M$ gegeben, mit $\left(J^{*}\right)^{2}=-1$. Die Eigenräume dieser Abbildung auf $T^{*} M \otimes_{\mathbb{R}} \mathbb{C}$ mit $\left(T^{*} M\right)^{(1,0)}$ und $\left(T^{*} M\right)^{(0,1)}$ und mit
\end{enumerate}

$$
\begin{gathered}
\Lambda^{(p, q)}\left(T^{-} M\right)^{\mathbb{C}}:=\operatorname{Span}\left\{u \wedge v \mid u \in \Lambda^{p}\left(T^{*} M\right)^{(1,0)}, v \in \Lambda^{q}\left(T^{*} M\right)^{(0,1)}\right\} \text { und } \\
\Omega^{(p, q)}(M):=\Gamma \Lambda^{(p, q)}\left(T^{*} M\right)^{\mathbb{C}}
\end{gathered}
$$

\subsection*{Definition}
Eine fast hermitsche Mannigfaltigkeit ( $M, g, J$ ) heißt kählersch, falls $\nabla^{g} J=0$ gilt.

\subsection*{Bemerkung}
\begin{enumerate}
  \item Jede Kähler-Mannigfaltigkeit ist eine komplexe Mannigfaltigkeit (vgl. [Ko-No]).
  \item Auf einer Kähler-Mannigfaltigkeit ist die Kähler-Form durch $\Omega(X, Y):= g(J X, Y)$ definiert. Diese ist geschlossen, und $(M, \Omega)$ wird damit zu einer symplektischen Mannigfaltigkeit.
\end{enumerate}

\subsection*{Bemerkung}
Wegen $\left(\nabla_{X} J\right)(Y)=X(J(Y))-J\left(\nabla_{X} Y\right)$ ist ( $M, g, J$ ) genau dann kählersch, wenn für die Parallelverschiebung des Levi-Civita-Zusammenhangs gilt

$$
\tau_{\gamma}(J v)=J\left(\tau_{\gamma}(v)\right)
$$

also wenn der Levi-Civita-Zusammenhang eine $U(m)$-Reduktion besitzt.\\
Für $X=A+i B \in u(m)$ ist $\operatorname{Spur} X=i \operatorname{Spur}(B)=-\frac{i}{2} \operatorname{Spur}\left(J_{0} f(X)\right)$, mit $J_{0}=\left(\begin{array}{cc}0 & -1 l \\ 1 l & 0\end{array}\right)$.\\
Damit gilt also für die Krümmung des Levi-Civita-Zusammenhangs $R^{g}$ aufgefasst als Krümmung $F$ im $U(m)$-Bündel $\Omega$

$$
-\frac{1}{2 \pi i} \operatorname{Spur}(F)=\frac{1}{4 \pi} \operatorname{Spur}\left(J R^{g}\right) .
$$

Um dies noch besser zu verstehen, berechnen wir die Ricci-Krümmung einer Kähler-Mannigfaltigkeit.

\subsection*{Lemma}
Ist ( $M, g, J$ ) eine Kähler-Mannigfaltigkeit, so gilt:

$$
\operatorname{ric}(X, Y)=\operatorname{ric}(J X, J Y)=\frac{1}{2} \operatorname{Spur}\left(J \circ R^{g}(X, J Y)\right)
$$

\section*{Beweis:}
Da $J$ parallel und orthogonal ist, ist

$$
\begin{aligned}
g\left(R^{g}(X, Y) Z, J V\right) & =-g\left(J R^{g}(X, Y) Z, V\right) \\
& =-g\left(R^{g}(X, Y) J Z, V\right)
\end{aligned}
$$

also wegen der Blocksymmetrie des Krümmungstensors auch

$$
R^{g}(J X, Y)=-R^{g}(X, J Y)
$$

und folglich wegen $J^{2}=-$ id auch $R^{g}(X, Y)=R^{g}(J X, J Y)$ und

$$
g\left(R^{g}(X, Y) Z, V\right)=g\left(R^{g}(X, Y) J Z, J V\right)
$$

Betrachtet man die Spur mittels der einer orthonormalen Basis

$$
\left(e_{1}, \ldots, e_{m}, J e_{1}, \ldots, J e_{m}\right)=:\left(e_{1}, \ldots, e_{2 m}\right)
$$

so rechnet man sofort nach:

$$
\operatorname{ric}(X, Y)=\operatorname{ric}(J X, J Y)
$$

und

$$
\begin{aligned}
\operatorname{ric}(X, Y) & =\sum_{i=1}^{2 m} g\left(R^{g}\left(X, e_{i}\right) J e_{i} J Y\right) \\
& =-\sum_{i=1}^{2 m}\left(g\left(R^{g}\left(e_{i}, J e_{i}\right) X, J Y\right)-g\left(R^{g}\left(J e_{i}, X\right) e_{i}, J Y\right)\right. \\
& =\sum_{i=1}^{2 m}\left(g\left(R^{g}(X, J Y) J e_{i}, e_{i}\right)-\operatorname{ric}(J X, J Y)\right)
\end{aligned}
$$

also

$$
\begin{aligned}
2 \operatorname{ric}(X, Y) & =\sum_{i=1}^{2 m} g\left(J R^{g}(X, J Y) e_{i}, e_{i}\right) \\
& =\operatorname{Spur}\left(J \circ R^{g}(X, J Y)\right)
\end{aligned}
$$

\subsection*{Korollar und Definition}
Sei $(M, g, J)$ eine Kähler-Mannigfaltigkeit, $\rho=\operatorname{ric}(J \cdot, \cdot) \in \Omega^{2}(M)$, so ist $d \rho=$ 0 und

$$
c_{1}(M)=\left[\frac{1}{2 \pi} \rho\right],
$$

insbesondere gilt für die 1. Chern-Klasse eine Ricci-flachen Kähler-Mannigfaltigkeit $c_{1}(M)=0$. Die 2-Form $\rho$ hießt die Ricci-Form.

\section*{Pontrjagin-Klassen}
Betrachtet man reelle Vektorraumbündel, so kann man diese komplexifizieren und ihre Chern-Klassen betrachten (vgl. 13.7,1). Pontrjagin-Klassen entstehen auf diese Weise.

\subsection*{Notationen und Bemerkungen}
Seien $V$ ein reeller und $W$ ein komplexer Vektorraum. Dann heißt

\begin{enumerate}
  \item $V^{\mathbb{C}}=V \otimes \mathbb{C}$ die Komplexifizierung von $V$.
  \item Der $W$ zugrunde liegende reelle Vektorraum $W^{\mathbb{R}}$ die Reellifizierung von $W$, d.h. $W^{\mathbb{R}}$ entsteht durch Einschränkung der skalaren Multiplikation in $W$ auf $\mathbb{R}$.
\end{enumerate}

Es gilt:

\begin{enumerate}
  \item $V^{*} \cong V, W^{*} \cong \bar{W}$.
  \item $\left(V^{\mathbb{C}}\right)^{\mathbb{R}}=V \oplus V$.\\
$\left(W^{\mathbb{R}}\right)^{\mathbb{C}}=W \oplus \bar{W}$.
  \item $\left(V^{*}\right)^{\mathbb{C}}=\left(V^{\mathbb{C}}\right)^{*}$.
\end{enumerate}

Analoge Schreibweisen und Sachverhalte gelten auch für Vektorraumbündel.

\subsection*{Bemerkung}
Ist $V$ ein reeller Vektorraum, so ist

$$
V^{\mathbb{C}}=V \otimes \mathbb{C} \cong V \oplus V,
$$

wobei die Multiplikation mit $i$ auf $V \oplus V$ durch $i(u, v)=(-v, u)$ gegeben ist. Ist also $V=W^{\mathbb{R}}$, so ist

$$
V^{\mathbb{C}}=\left(W^{\mathbb{R}}\right)^{\mathbb{C}} \cong W \oplus W .
$$

Dabei ist die Abbildung

$$
\left(W^{\mathbb{R}}\right)^{\mathbb{C}} \rightarrow W \oplus W
$$

durch $f(u, v)=(u, i v, u-i v)$ gegeben und $\mathbb{C}$ linear mit der Multiplikation $i(u, v)=(i u,-i v)$, denn

$$
\begin{aligned}
f(i(v, v)) & =f(-v, u)=(i u-v,-v-i u) \\
& =i(u+i v, u-i v) \\
& =i f(u, v)
\end{aligned}
$$

Man schreibt daher $\left(W^{\mathbb{R}}\right)^{\mathbb{C}}=W \oplus \bar{W}$.

\subsection*{Korollar}
Ist $E \rightarrow M$ ein reelles Vektorraumbündel, so ist $c_{k}\left(E^{\mathbb{C}}\right)=0$ für $k$ ungerade, da

$$
c_{k}\left(E^{\mathbb{C}}\right)=c_{k}\left(\left(E^{*}\right)^{\mathbb{C}}\right)=c_{k}\left(\left(E^{\mathbb{C}}\right)^{*}\right)=(-1)^{k} c_{k}\left(E^{\mathbb{C}}\right)
$$

\subsection*{Definition}
Sei $E$ ein reelles $n$-dimensionales Vektorraumbündel, dann heißt

$$
p_{k}(E)=(-1)^{k} c_{2 k}\left(E^{\mathbb{C}}\right) \in H_{d R}^{4 k}(M, \mathbb{R})
$$

die $k$-te reelle Pontrjagin-Klasse von $E$,

$$
p(E):=p_{0}(E)+p_{1}(E)+\ldots+p_{\left[\frac{n}{2}\right]}(E)
$$

die totale Pontrjagin-Klasse von $E$.

\subsection*{Korollar}
\begin{enumerate}
  \item Ist $E_{1} \cong E_{2}$, so gilt
\end{enumerate}

$$
p\left(E_{1}\right)=p\left(E_{2}\right)
$$

\begin{enumerate}
  \setcounter{enumi}{1}
  \item Ist $\varphi: N \rightarrow M$ differenzierbar, so ist
\end{enumerate}

$$
p\left(\varphi^{*} E\right)=\varphi^{*} p(E)
$$

\begin{enumerate}
  \setcounter{enumi}{2}
  \item Es gilt
\end{enumerate}

$$
p\left(E_{1} \oplus E_{2}\right)=p\left(E_{1}\right) p\left(E_{2}\right)
$$

\section*{Beweis:}
\begin{enumerate}
  \item und 2 . folgen aus 13.4 und 13.6 .\\
zu 3. Es ist
\end{enumerate}

$$
\begin{aligned}
(-1)^{k} \sum_{j=0}^{2 k} c_{j}\left(E_{1}^{\mathbb{C}}\right) c_{2 k-j}\left(E_{2}^{\mathbb{C}}\right) & \left.=\sum_{r=0}^{k}(-1]\right)^{r} c_{2 r}\left(E_{1}^{\mathbb{C}}\right) \cdot(-1)^{k-r} c_{2 k-2 r}\left(E_{2}^{\mathbb{C}}\right) \\
& =\sum_{r=0}^{k} p_{1}\left(E_{1}\right) p_{k-r}\left(E_{2}\right)
\end{aligned}
$$

\subsection*{Bemerkung}
Ist $E \oplus F=M \times \mathbb{R}^{N}$, so ist $p(E)$ invertierbar, also $p(E) p(F)=1$, für

$$
p(E) \in \bigoplus_{k=0}^{\left[\frac{n}{2}\right]} H^{4 k}(M), p(F) \in \bigoplus_{k=0}^{\left[\frac{N-n}{2}\right]} H^{4 k}(M)
$$

\subsection*{Lemma}
Seien $g_{k}: g l(n, \mathbb{R}) \rightarrow \mathbb{R}$ definiert durch

$$
\operatorname{det}\left(t 11-\frac{1}{2 \pi} X\right)=: \sum_{k=0}^{n} g_{k}(X) t^{n-k}
$$

Dann ist

$$
p_{k}(E)=\mathcal{W}^{p}\left(g_{2 k}\right)
$$

\section*{Beweis:}
Ist $E=P \times_{G L(n, \mathbb{R})} \mathbb{R}^{n}$, so ist $E^{\mathbb{C}}=P \times_{G L(n, \mathbb{R})} \mathbb{C}^{n}$, also ist $P$ eine $G L(n, \mathbb{R})$ Reduktion des $E^{\mathbb{C}}$ zugrunde liegenden Prinzipalbündels. Also ist (vgl. 12.6 und 12.14)

$$
c_{2 k}\left(E^{\mathbb{C}}\right)=\mathcal{W}^{P}\left(\left.c_{2 k}\right|_{g l(n, \mathbb{R})}\right)
$$

Wegen

$$
\operatorname{det}\left(t 1 l-\frac{1}{2 \pi i} X\right)=(-i)^{n} \operatorname{det}\left(i t 1 l-\frac{1}{2 \pi} X\right)
$$

ist $c_{k}(i X)=i^{k} c_{k}(X)$ und folglich gilt für $X \in g l(n, \mathbb{R})$

$$
(-1)^{k} c_{2 k}(X)=c_{2 k}(i X)=g_{k}(X)
$$

und damit

$$
p_{k}(E)=(-1)^{k} c_{2 k}\left(E^{\mathbb{C}}\right)=\left[g_{2 k}\left(F^{\nabla}\right)\right]
$$

\subsection*{Beispiel}
Wir haben gesehen, dass $c\left(T \mathbb{C} P^{n}\right)=(1-a)^{n+1}$, für $a=c_{1}\left(l_{n}\right)$ gilt. Also ist

$$
\begin{aligned}
c\left(T \mathbb{C} P^{n} \otimes \mathbb{C}\right)= & =c\left(T \mathbb{C} P^{n} \oplus \overline{T \mathbb{C} P^{n}}\right) \\
& =(1-a)^{n+1}(1+a)^{n+1} \\
& =\left(1-a^{2}\right)^{n+1}
\end{aligned}
$$

Also ist

$$
\begin{aligned}
p\left(T \mathbb{C} P^{1}\right) & =\left(1+a^{2}\right)^{2}=1 \\
p\left(T \mathbb{C} P^{2}\right) & =\left(1+a^{2}\right)^{3}=1+3 a^{2} \\
p\left(T \mathbb{C} P^{3}\right) & =\left(1+a^{2}\right)^{4}=1+4 a^{2}
\end{aligned}
$$

\subsection*{Anwendung}
Die Pontrjagin-Klasse kann oft helfen, zu entscheiden, ob eine Mannigfaltigkeit $M$ in $\mathbb{R}^{N}$ einzubetten ist. Ist nämlich $M \subseteq \mathbb{R}^{N}$, so ist $T M \oplus T^{\perp} M=M \times \mathbb{R}^{n}$, also $p_{1}(T M) p\left(T^{\perp} M\right)=1$.\\
Jede $n$-dimensionale Manigkfaltigkeit lässt sich in den $\mathbb{R}^{2 n}$ einbetten, jede nicht kompakte sogar in den $\mathbb{R}^{2 n-1}$, jedoch stellt sich oft die Frage, ob es auch in niedriger Kodimension möglich ist, so existiert z.B. eine Einbettung des $\mathbb{R} P^{1}$ in $\mathbb{R}^{2}$, aber keine von $\mathbb{R} P^{2}$ in den $\mathbb{R}^{3}$ (nur eine Immersion).\\
Wir zeigen hier: $\mathbb{C} P^{4}$ kann nicht in den $\mathbb{R}^{9}$ eingebettet werden. Angenommen doch, dann wäre $T^{\perp} \mathbb{C} P^{4}$ ein 1-dimensionales Vektorraumbündel mit Pontrjagin-Klasse

$$
\left(1+a^{2}\right)^{-5}=\left(1+5 a^{2}+15 a^{4}\right)^{-1}=1-5 a^{2}+10 a^{4}
$$

Dies ist aber nicht die Pontrjagin-Klasse eines 1-dimensionalen Vektorraumbündels. Genauer sieht man, dass, falls $\mathbb{C} P^{4} \subseteq \mathbb{R}^{n}$ ist, mindestens $n \geq 8+4=12$ gelten muss.

Für die $k$-te Pontrjagin-Klasse $p(T M)$ einer Riemannschen Mannigfaltigkeit $M$ gilt damit für eine Orthonormalbasis $\left(e_{1}, \ldots, e_{n}\right)$ von $T M$

$$
p_{k}\left(\mathcal{R}^{g}\right)=\left(\frac{1}{2 \pi}\right)^{2 k} \sum_{i_{1}<\ldots<i_{2 k}} \operatorname{sign}(\tau) \mathcal{R}^{g}\left(e_{i_{1}} \wedge e_{i_{\tau(1)}}\right) \wedge \ldots \wedge \mathcal{R}^{g}\left(e_{i_{2 k}} \wedge e_{i_{\tau(2 k)}}\right)
$$

wobei $\mathcal{R}: \Lambda^{2} T^{*} M \rightarrow \Lambda^{2} T^{*} M$ der durch den Riemannschen Krümmungstensor gegebene Endomorphismus ist, also

$$
\mathcal{R}\left(e_{i} \wedge e_{j}\right)=\sum_{k<l} g\left(R\left(e_{i}, e_{j}\right) e_{k}, e_{l}\right) e_{k} \wedge e_{l}
$$

Dabei haben wir $T^{*} M=T M$ auf Riemannschen Mannigfaltigkeiten benutzt. Da Spur $\left(\mathcal{R}\left(e_{i} \wedge e_{j}\right)\right)=0$ gilt, ist also

$$
p_{1}\left(R^{g}\right)=\frac{1}{2 \pi^{2}} \sum_{i<j} \mathcal{R}\left(e_{i} \wedge e_{j}\right) \wedge \mathcal{R}\left(e_{j} \wedge e_{i}\right)
$$

Wir wollen jetzt in einem Spezialfall eine nähere Beziehung zur Differentialgeometrie zeigen. Dazu betrachten wir eine Zerlegung des Krümmungstensors. Sei $s$ die Skalarkrümmung und ric die Ricci-Krümmung einer Riemannschen Mannigfaltigkeit, dann betrachte man die Zerlegung des Krümmungstensors

$$
R^{g}=\frac{s}{2 n(n-1)} g \oplus g+\frac{1}{(n-2)}\left(\text { ric }-\frac{s}{n} g\right) \otimes g+W^{g} .
$$

Dabei ist ( $\triangle$ das Kulkarni-Normizu-Produkt\\
$(a \circlearrowleft b)(x, y, u, v)=a(x, u) b(y, v)-a(x, v) b(y, u)+a(y, v) b(x, u)-a(y, u) b(x, v)$.\\
Sind $a, b \in S^{2} V^{*}$, so hat $a \otimes b$ dieselben Symmetrieeigenschaften wie der Krümmungstensor. Fassen wir also die Summanden als Schnitte in End $\left(\Lambda^{2} T M\right)$ auf, so ist $g\left(\triangle g=2 \cdot \operatorname{Id}_{\Lambda^{2} T M}\right.$ und da Spur(ric $\left.-\frac{s}{n} g\right)=0$, ist auch die Spur des Schoutentensors Spur(ric $-\frac{s}{n} g$ ) $\triangle g=0$. $W^{g}$ heißt der Weyl-Tensor. Seine Bedeutung liegt darin, dass er eine konforme Invariante ist. Er hat dieselben Symmetrieeigenschaften wie der Krümmungstensor und ist spurfrei, also $\sum_{i=1}^{n} e_{i} W\left(u, e_{i}, e_{i}, v\right)=0$ (vgl. Besse, Einstein manifolds, S. 59 ff).\\
In Dimension 4 liefert nun der Sternoperator eine Aufspaltung von $\Lambda^{2} T M= \Lambda^{2} T_{+} M \oplus \Lambda^{2} T_{-} M$. Bezüglich dieser Aufspaltung ist $\mathcal{R}$ durch

$$
\mathcal{R}=\left(\begin{array}{cc}
W_{+} & 0 \\
0 & W_{-}
\end{array}\right)+\frac{s}{12} 11+\left(\begin{array}{cc}
0 & B \\
B^{*} & 0
\end{array}\right)
$$

wobei $W=\left(\begin{array}{cc}W_{+} & 0 \\ 0 & W_{-}\end{array}\right)$der Weyl-Tensor ist, beschrieben wird. Es gilt $\operatorname{Spur} W_{+}=\operatorname{Spur} W_{-}=0$. Um dies zu erläutern, erinnern wir an die Definition des *-Operators.

\subsection*{Exkurs: *-Operator}
Die Metrik $g$ auf $T M$ induziert ein Skalarprodukt $<,>$ auf $\Lambda^{2} T_{x} M$ folgendermaßen. ist ( $e_{1}, \ldots, e_{n}$ ) eine Sylvesterbasis in $T_{x} M$ mit $\varepsilon_{i}=g\left(e_{i}, e_{i}\right)$, so definieren wir

$$
<e_{i_{1}} \wedge \ldots \wedge e_{i_{k}}, e_{j_{1}} \wedge \ldots \wedge e_{j_{k}}>=\operatorname{sign}\left(\begin{array}{ccc}
i_{1} & \ldots & i_{k} \\
j_{1} & \ldots & j_{k}
\end{array}\right) \cdot \varepsilon_{i_{1}} \ldots \varepsilon_{i_{k}}
$$

für

$$
\left\{i_{1}, \ldots, i_{k}\right\}=\left\{j_{1}, \ldots, j_{k}\right\} \text { und }<e_{i_{1}} \wedge \ldots \wedge e_{i_{k}} e_{j_{1}} \wedge \ldots \wedge e_{j_{k}}>=0
$$

falls $\left\{i_{1}, \ldots, i_{k}\right\} \neq\left\{j_{1}, \ldots, j_{k}\right\}$.\\
Aufgrund der Blocksymmetrie des Riemannschen Krümmungstensors $R^{g}$ ist $\mathcal{R}^{g}$ selbstadjungiert bezüglich dieses Skalarprodukts.\\
Der Hodge-Operator

$$
*: \Omega^{k}(M) \rightarrow \Omega^{n-k}(M)
$$

auf einer orientierbaren Riemannschen Mannigfaltigkeit ist durch

$$
<* \omega, \eta>\operatorname{det}=\omega \wedge \eta
$$

gegeben. In Dimension 4 erhält man dadurch eine Abbildung

$$
*: \Lambda^{2} T^{*} M \rightarrow \Lambda^{2} T^{*} M
$$

und auf Riemannschen Mannigfaltigkeiten gilt für 2-Formen stets ** = id. Wir bezeichnen mit $\Lambda_{ \pm}^{2} T^{*} M$ die Eigenräume zu 1 und -1 .\\
Wir untersuchen nun die Zerlegung des Krümmungstensors $\mathcal{R} \in \operatorname{End}\left(\Lambda^{2} T M\right)$ bezüglich der Zerlegung $\Lambda^{2} T M=\Lambda_{+}^{2} T M \oplus \Lambda_{-}^{2} T M, \mathcal{R}=\left(\begin{array}{cc}A_{+} & S_{-} \\ S_{+} & A_{-}\end{array}\right)$. Aufgrund der Blocksymmetrie folgt sofort, dass $S_{+}=S_{-}^{T}$ gilt. Außerdem rechnet man leicht nach, dass der Schouten-Tensor nur zu $S_{+}$und $S_{-}$beiträgt. Für die Rechnungen wählt man dazu eine Orthonormalbasis aus Eigenvektoren des Ricci-Tensors $\left(e_{1}, e_{2}, e_{3}, e_{4}\right)$ und damit eine Orthonormalbasis von $\Lambda_{ \pm}^{2} T M$ durch

$$
\begin{aligned}
\alpha_{1}^{ \pm} & =\frac{1}{\sqrt{2}}\left(e_{1} \wedge e_{2} \pm e_{3} \wedge e_{4}\right) \\
\alpha_{2}^{ \pm} & =\frac{1}{\sqrt{2}}\left(e_{2} \wedge e_{3} \pm e_{4} \wedge e_{1}\right) \\
\alpha_{3}^{ \pm} & =\frac{1}{\sqrt{2}}\left(e_{3} \wedge e_{1} \pm e_{2} \wedge e_{4}\right)
\end{aligned}
$$

Ist $\left(e_{1}, e_{2}, e_{3}, e_{4}\right)$ eine Orthonormalbasis von Eigenvektoren von $\alpha$ zum Eigenwert $\lambda_{j}$, so ist

$$
(\alpha \bigoplus g)\left(e_{i} \wedge e_{j}, e_{i} \wedge e_{j}\right)=\lambda_{i}+\lambda_{j}
$$

$$
(\alpha \oplus g)\left(e_{i} \wedge e_{j}, e_{r} \wedge e_{s}\right)=0 \text { falls }\{i, j\} \neq\{r, s\} .
$$

Nutzt man nun aus, dass die Spur von $Z=2\left(\right.$ ric $\left.-\frac{s}{4} s g\right)$ verschwindet, so ist

$$
g(Z \oslash g)\left(\alpha_{i}^{ \pm}, \alpha_{j}^{ \pm}\right)=0
$$

für alle $i, j$.\\
Mithilfe der Spurfreiheit des Weyl-Tensors führt eine ähnliche Rechnung zu dem Ergebnis, dass $W^{g}$ die Zerlegung von $\Lambda^{2} T M$ respektiert also $A_{ \pm}=\left(W^{g}+\right. \left.\frac{s}{12} \mathbb{1}\right)\left.\right|_{\Lambda_{ \pm}^{2} T M}=: W_{ \pm}+\frac{s}{12}$ und $Z \oplus g=\left(\begin{array}{cc}0 & S_{-} \\ S_{+} & 0\end{array}\right)$ gilt.\\
Nutzt man die Biancchi-Identität aus, so folgt, dass

$$
\operatorname{Spur} A_{+}=\sum_{i=1}^{3} g\left(\mathcal{R}\left(\alpha_{i}^{+}\right), \alpha_{i}^{+}\right)=-\frac{1}{4} s=\operatorname{Spur} A_{-}
$$

gilt, also ist

$$
\operatorname{Spur}\left(A_{ \pm}-\frac{s}{12} \mathbb{1}\right)=\operatorname{Spur}\left(A_{ \pm}+\frac{1}{3} \operatorname{Spur} A \mathbb{1}\right)=0
$$

und damit Spur $W_{+}=\operatorname{Spur} W_{-}=0$.

\subsection*{Satz}
Auf einer 4-dimensionalen orientierten Mannigfaltigkeit ist

$$
p_{1}(T M)=\frac{1}{4 \pi^{2}}\left(\left\|W_{+}\right\|^{2}-\left\|W_{-}\right\|^{2}\right) \omega_{M}
$$

bezüglich einer beliebigen Riemannschen Metrik auf $M$.

\subsection*{Korollar}
Ist $p_{1}(T M) \neq 0$, so ist $M$ nicht konform flach.

\section*{Beweis:}
Es gilt

$$
4 \pi^{2} p_{1}\left(R^{g}\right)=\sum_{i=1}^{3}\left(\mathcal{R}\left(\alpha_{i}^{+}\right) \wedge \mathcal{R}\left(\alpha_{i}^{+}\right)+\mathcal{R}\left(\alpha_{i}^{-}\right) \wedge \mathcal{R}\left(\alpha_{i}^{-}\right)\right)
$$

Es ist

$$
\sum_{i=1}^{3}\left(\mathcal{R}\left(\alpha_{i}^{ \pm}\right) \wedge \mathcal{R}\left(\alpha_{i}^{ \pm}\right)\right)=
$$

$=\sum_{i=1}^{3} W_{ \pm}\left(\alpha_{i}^{ \pm}\right) \wedge W_{ \pm}\left(\alpha_{i}^{ \pm}\right)-\frac{s}{6} W_{ \pm}\left(\alpha_{i}^{ \pm}\right) \wedge \alpha_{i}^{ \pm}+S_{ \pm}\left(\alpha_{i}^{ \pm}\right) \wedge S_{ \pm}\left(\alpha_{i}^{ \pm}\right)+\frac{s^{2}}{144} \alpha_{i}^{ \pm} \wedge \alpha_{i}^{ \pm}$,\\
denn

$$
\left(W_{ \pm}\left(\alpha_{i}^{ \pm}\right)+\frac{s}{12} \alpha_{i}^{ \pm}\right) \wedge S_{ \pm}\left(a_{i}^{ \pm}\right)=0
$$

da $\alpha_{i}^{ \pm} \wedge \alpha_{j}^{\mp}=0$.\\
Wegen $\alpha_{i}^{ \pm} \wedge \alpha_{i}^{ \pm}= \pm \omega_{M}$ und $\alpha_{i}^{ \pm} \wedge \alpha_{j}^{ \pm}=0$ für $i \neq j$, ist

$$
\begin{aligned}
\sum_{i=1}^{3} W_{ \pm}\left(\alpha_{i}^{ \pm}\right) \wedge \alpha_{i}^{ \pm} & =\sum_{i, j=1}^{3} g\left(W_{ \pm}\left(\alpha_{i}^{ \pm}\right), \alpha_{j}^{ \pm}\right) \alpha_{j}^{ \pm} \wedge \alpha_{i}^{ \pm} \\
& = \pm \sum_{i=1}^{3} g\left(W_{ \pm}\left(\alpha_{i}^{ \pm}\right), \alpha_{i}^{ \pm}\right) \omega_{M} \\
& = \pm \operatorname{Spur} W_{ \pm} \omega_{M}=0
\end{aligned}
$$

und die Terme $\frac{s^{2}}{144} \alpha_{i}^{ \pm} \wedge \alpha_{i}^{ \pm}$heben sich gegenseitig auf.\\
Weiter ist

$$
\begin{aligned}
\sum_{i=1}^{3} W_{ \pm}\left(\alpha_{i}^{ \pm}\right) \wedge W_{ \pm}\left(\alpha_{i}^{ \pm}\right) & = \pm \sum_{i=1}^{3}<W_{+}\left(\alpha_{i}^{ \pm}\right), W_{+}\left(\alpha_{i}^{ \pm}\right)>\omega_{m} \\
& = \pm\left\|W_{ \pm}\right\|^{2} \omega_{m}
\end{aligned}
$$

Schließlich nutzt man noch

$$
S_{ \pm}\left(\alpha_{i}^{ \pm}\right)=g\left(\mathcal{R}\left(\alpha_{i}^{ \pm}\right), \alpha_{i}^{\mp}\right) \alpha_{i}^{\mp}
$$

und folglich

$$
S_{+}\left(\alpha_{i}^{+}\right) \wedge S_{+}\left(\alpha_{i}^{+}\right)=-S_{-}\left(\alpha_{i}^{-}\right) \wedge S_{-}\left(\alpha_{i}^{-}\right)
$$

um das Ergebnis zu erhalten.

\section*{Die Euler-Klasse}
\subsection*{Notation}
Ist $A \in s o(n)$, so fassen wir im Folgenden $A \in \Lambda^{2} \mathbb{R}^{n}$ auf, also $A=\left(a_{i j}\right)_{i, j}$ lesen wir auch als

$$
A=\sum_{i<j} a_{i j} e_{i} \wedge e_{j}
$$

\subsection*{Definition}
Für $A \in s o(2 m)$ ist die Pfaffsche Zahl $\operatorname{Pf}_{2 m}(A)$ durch

$$
A \wedge \ldots \wedge A=m!\mathrm{Pf}_{2 m}(A) \operatorname{det}
$$

wohldefiniert.

\subsection*{Bemerkung}
Es gilt $\operatorname{Pf}\left(X^{T} A X\right)=\operatorname{Pf}(A) \operatorname{det} X$, also ist $\operatorname{Pf} S O(2 m)$-invariant (aber nicht $G L(n, \mathbb{R})$-invariant!).

\subsection*{Lemma}
\begin{enumerate}
  \item Die Abbildung
\end{enumerate}

$$
\mathrm{Pf}_{2 m}: \operatorname{so}(2 m) \rightarrow \mathbb{R}
$$

ist ein $S O(2 m)$-invariantes, homogenes Polynom vom Grad $m$. Es gilt:

$$
2^{m} m!\operatorname{Pf}_{2 m}(A)=\sum_{\tau} \operatorname{sign}(\tau) a_{\tau(1) \tau(2)} \ldots a_{\tau(2 m-1) \tau(m)}
$$

Insbesondere gilt

$$
\begin{aligned}
\operatorname{Pf}_{2}(A) & =a_{12} \\
\operatorname{Pf}_{4}(A) & =a_{12} a_{34}-a_{13} a_{24}+a_{14} a_{23}
\end{aligned}
$$

\begin{enumerate}
  \setcounter{enumi}{1}
  \item $\left(\mathrm{Pf}_{2 m}(A)\right)^{2}=\operatorname{det} A$
  \item Bezeichnet $j: u(m) \rightarrow s o(2 m)$ die übliche Inklusion, so gilt
\end{enumerate}

$$
\operatorname{Pf}_{2 m}(j(X))=(-1)^{m(m-1) / 2} i^{m} \operatorname{det}(X)
$$

\section*{Beweis:}
\begin{enumerate}
  \item Folgt direkt aus der Definition des Dachprodukts.
  \item Schiefsymmetrische Matrizen können in gerader Dimension durch Konjugation mit einer orthogonalen Matrix in Standardform $\left(a_{i j}\right)$ mit $a_{2 i-1,2 i}= a_{2 i, 2 i-1}=\lambda_{i}$ für $i=1, \ldots, m$ und $a_{i j}=0$ sonst überführt werden.\\
In dieser Basis ( $v_{1}, \ldots, v_{2 m}$ ) gilt dann
\end{enumerate}

$$
A=\sum_{i=1}^{m} \lambda_{i} v_{2 i-1} \wedge v_{2 i}
$$

und folglich

$$
A \wedge \ldots \wedge A=m!\lambda_{1} \ldots \lambda_{m} \operatorname{det}
$$

also

$$
\operatorname{Pf}_{m}(A)=\lambda_{1} \cdot \ldots \cdot \lambda_{m} \text { und } \operatorname{det}(A)=\lambda_{1}^{2} \ldots \lambda_{m}^{2} .
$$

\begin{enumerate}
  \setcounter{enumi}{2}
  \item Ist $X=A+i B$, so ist
\end{enumerate}

$$
\begin{aligned}
\left(\mathrm{Pf}_{2 m}(j(X))^{2}\right. & =\operatorname{det}\left(\begin{array}{cc}
A & -B \\
B & A
\end{array}\right)=\operatorname{det}\left(\begin{array}{cc}
A+i B & -B \\
0 & A-i B
\end{array}\right) \\
& =\operatorname{det}(A+i B) \operatorname{det}(A-i B)=|\operatorname{det} X|^{2}
\end{aligned}
$$

Also ist $\operatorname{Pf}_{2 m}(j(X))=f(X) \operatorname{det}(X)$ für $f: u(n) \rightarrow S^{1}$ eine Abbildung $f: u(n) \rightarrow S^{1}$.\\
Wegen $\bar{X}=-X^{T}$ ist $\overline{\operatorname{det} X}=(-1)^{m} \operatorname{det} X$, also $f(X)= \pm i^{m}$, also ist $f$ konstant und für $X=i 11$ ergibt sich der gesuchte Vorfaktor.

\subsection*{Definition}
Sei $E$ ein orientiertes $2 m$-dimensionales Riemannsches Vektorraumbündel, dann heißt

$$
e(E)=\frac{1}{(2 \pi)^{m}} \mathcal{W}^{P}\left(\operatorname{Pf}_{2 m}\right) \in H_{d R}^{2 m}(M)
$$

die Euler-Klasse von $(E, g)$.\\
Dabei ist $P$ die Reduktion von $E$ auf $S O(2 m)$.\\
Ist $E=T M$, so heißt

$$
e(M)=e(T M)
$$

auch die Euler-Klasse von M

\subsection*{Bemerkung}
Beachte. dass zur Berechnung der Euler-Klasse eine Riemannsche Metrik und Orientierung nötig ist!\\
Ist $\Omega$ die Krümmung eines isometrischen Zusammenhangs, so gilt

$$
e(E)=\left(\frac{1}{4 \pi}\right)^{m} \frac{1}{(2 m)!} \sum \operatorname{sign}(\tau) \Omega_{\tau(1) \tau(2)} \ldots \Omega_{\tau(2 m-1) \tau(2 m)}
$$

\subsection*{Satz}
Es gilt für ein $2 m$-dimensionales orientiertes Riemannsches Vektorraumbündel

\begin{enumerate}
  \item $e(E) e(E)=p_{m}(E)$,
  \item $e(E \oplus F)=e(E) e(F)$
\end{enumerate}

\section*{Beweis:}
\begin{enumerate}
  \item $\left.\mathrm{Pf}_{2 m}(X)\right)^{2}=\operatorname{det}(X)$.
  \item Betrachtet man Pf für Bockdiagonalmatrizen, so folgt die Behauptung.
\end{enumerate}

\subsection*{Bemerkung}
Ist $E$ ein komplexes Vektorraumbündel, so ist

$$
e\left(E^{\mathbb{R}}\right)=(-1)^{m(m-1) / 2} c_{m}(E) .
$$

Die $n$-te Chernklasse eines komplexen Vektorraumbündels hängt also nicht von der komplexen Struktur ab. Um dies zu zeigen, benutzt man zur Berechnung der Euler-Klasse den Realteil einer hermitschen Bündelmetrik

Chern-, Pointrjagin- und Euler-Klassen erzeugen invariante Polynomringe. Es gilt

\subsection*{Lemma}
Sei $P$ ein Polynom mit $P\left(T A T^{-1}\right)=P(A)$ für alle $A \in s o(k), T \in S O(k)$.

\begin{enumerate}
  \item Ist $k$ ungerade, so ist $P O(k)$-invariant und kann in Pontrjagin-Klassen geschrieben werden.
  \item Ist $k$ gerade, so ist $P=P_{0}+e \cdot P_{1}$, wobei $e$ die Euler-Klasse und $P_{i} O(k)$-invariant sind.
\end{enumerate}

\section*{Beweis:}
Gilkey, Invariance theory

\subsection*{Satz}
\begin{enumerate}
  \item Sei $\operatorname{dim} M=2, G$ die Gaussche Krümmung, so ist
\end{enumerate}

$$
\operatorname{Pf}_{2}\left(R^{g}\right)=G \omega_{M}
$$

\begin{enumerate}
  \setcounter{enumi}{1}
  \item Ist $M$ eine 4-dimensionale orientierte Riemannsche Mannigfaltigkeit, so ist
\end{enumerate}

$$
\mathrm{Pf}_{2}\left(R^{g}\right)=\frac{1}{8}\left(\|R\|^{2}-4\|\mathrm{ric}\|^{2}+\mathrm{scal}^{2}\right) \omega_{M} .
$$

\section*{Beweis:}
Es ist $\operatorname{Pf}_{2 m}\left(R^{g}\right)=\frac{(-1)^{m}}{2^{m} m!} \sum_{\tau} \operatorname{sign}(\tau) \mathcal{R}\left(e_{\tau(1)}\right) \wedge e_{\tau(2)} \wedge \ldots \mathcal{R}\left(e_{\tau(2 m-1)} \wedge e_{\tau(2 m)}\right)$.\\
1.

$$
\mathrm{Pf}_{2}\left(R^{g}\right)\left(e_{1}, e_{2}\right)=g\left(R\left(e_{1}, e_{2}\right) e_{1}, e_{2}\right)=G .
$$

\begin{enumerate}
  \setcounter{enumi}{1}
  \item Folgt aus einer längeren Rechnung.
\end{enumerate}

\subsection*{Satz (Chern-Gauß-Bonnet)}
Für eine $2 m$-dimensionale kompakte orientierte Mannigfaltigkeit gilt

$$
\int_{M^{2 m}} e\left(M^{2 m}\right)=\chi(M)=\sum_{j=1}^{2 m}(-1)^{j} \operatorname{dim} H_{d R}^{j}\left(M^{2 m}\right)
$$

\section*{Beweis:}
Siehe Alfred Gray: Tubes, oder Peter Gilkey: Invariance Theory

\subsection*{Beispiel}
$M=S^{2}$. Es bezeichnen $\theta, \varphi$ die Kugelkoordinaten. Dann ist

$$
R^{g}\left(\partial_{\theta}, \partial_{\varphi}\right)=\sin \theta d \theta \wedge d \varphi
$$

Also ist $e(M)=\frac{1}{2 \pi} \sin \theta d \theta \wedge d \varphi$.

\section*{Multiplikative und additive Klassen}
Ist $E \rightarrow M$ ein komplexes Vektorraumbündel, $E=L_{1} \oplus \ldots \oplus L_{n}$ mit komplexen Linienbündeln $L_{i}$, so ist

$$
c(E)=\prod_{i=1}^{n}\left(1+c_{1}\left(L_{i}\right)\right)
$$

also ist

$$
c_{k}(E)=\sigma_{k}\left(x_{1}, \ldots x_{n}\right)
$$

wobei $x_{j}=c_{1}\left(L_{j}\right) \in H^{2}(M)$. Analog ist

$$
p_{k}\left(E^{\mathbb{R}}\right)=\sigma_{k}\left(x_{1}^{2}, \ldots, x_{n}^{2}\right)
$$

da

$$
\left.() L_{i}^{\mathbb{R}}\right)^{\mathbb{C}} \cong L_{i} \oplus \bar{L}_{i},
$$

also $c_{2}\left(\left(L_{i}^{\mathbb{R}}\right)^{\mathbb{C}}\right)=-x_{i}^{2}$. Natürlich können wir nicht jedes Vektorraumbündel in Linearbündel zerlegen. Wir können aber ein anschauliches Verfahren finden, um neue charakteristische Klassen zu definieren.

\subsection*{Satz}
Es bezeichne $\sigma_{k}$ die elementarsymmetrischen Polynome vom Grad $k$, also

$$
\sigma_{k}\left(x_{1}, \ldots, x_{n}\right):=\sum_{i_{1}<\ldots i_{k}} x_{i_{1}} \ldots x_{i_{k}} \text { und } \sigma_{k}=0 \text { für } k>n \text {, }
$$

$Q$ eine formale Potenzreihe mir rationalen Koeffizienten und konstantem Summanden 1, dann existieren Potenzreihen $S_{Q}$ und $P_{Q}$ in $n$ Variablen, sodass

$$
\begin{aligned}
\sum_{i=1}^{n} Q\left(x_{i}\right) & =S_{Q}\left(\sigma_{1}\left(x_{1}, \ldots, x_{n}\right), \ldots, \sigma_{n}\left(x_{1}, \ldots, x_{n}\right)\right) \\
& =: n+S_{1}\left(\sigma_{1}\left(x_{1}, \ldots, x_{n}\right)\right)+S_{2}\left(\sigma_{1}\left(x_{1}, \ldots, x_{n}\right),\left(\sigma_{2}\left(x_{1}, \ldots, x_{n}\right)\right)+\ldots\right. \\
\prod_{i=1}^{n} Q\left(x_{i}\right) & =P_{Q}\left(\sigma_{1}\left(x_{1}, \ldots, x_{n}\right), \ldots, \sigma_{n}\left(x_{1}, \ldots, x_{n}\right)\right) \\
& =1+P_{1}\left(\sigma_{1}\left(x_{1}, \ldots, x_{n}\right)\right)+\ldots+
\end{aligned}
$$

für Polynome $S_{j}, P_{j}$ vom Grad $j$.\\
Es gilt für $Q=1+a_{1} x+\ldots$

$$
\begin{aligned}
S_{1}\left(\sigma_{1}\right) & =a_{1} \sigma_{1} \\
S_{2}\left(\sigma_{1}, \sigma_{2}\right) & =a_{2}\left(\sigma_{1}^{2}-2 \sigma_{2}\right) \\
S_{3}\left(\sigma_{1}, \sigma_{2}, \sigma_{2}\right) & =a_{3} \sigma_{1}^{3}-3 a_{3} \sigma_{1} \sigma_{2}+3 a_{3} \sigma_{3} \\
P_{1}\left(\sigma_{1}\right) & =a_{1} \sigma_{1} \\
P_{2}\left(\sigma_{2}\right) & =a_{2}\left(\sigma_{1}^{2}-2 \sigma_{2}\right)+a_{1}^{2} \sigma_{2} \\
P_{3}\left(\sigma_{1}, \sigma_{2}\right) & =a_{3} \sigma_{1}^{3}+\left(a_{1} a_{2}-3 a_{3}\right) \sigma_{1} \sigma_{2}+\left(3 a_{3}-3 a_{1} a_{2}+a_{1}^{3}\right) \sigma_{3}
\end{aligned}
$$

Auf diese Weise können neue charakteristische Klassen aus Chern- und PontrjaginKlassen gebildet werden. Wir führen dies nur für die Chern-Klassen vor:

$$
\begin{aligned}
Q_{+}(E) & =S_{Q}\left(c_{1}(E), \ldots, c_{n}(E)\right) \\
Q_{\bullet}(E) & =P_{Q}\left(c_{1}(E), \ldots, c_{n}(E)\right)
\end{aligned}
$$

und analog für die Pontrjagin-Klassen. Charakteristische Klassen der Form $Q_{+}(E)$ heißen additive charakteristische Klassen, der Form $Q_{\bullet}$ multiplikative charakteristische Klassen. Aus $c\left(E_{1} \oplus E_{2}\right)=c\left(E_{1}\right) c\left(E_{2}\right)$ folgt für additive charakteristische Klassen:

$$
Q_{+}\left(E_{1} \oplus E_{2}\right)=Q_{+}\left(E_{1}\right)+Q_{+}\left(E_{2}\right)
$$

Für multiplikative charakteristische Klassen gilt

$$
Q_{\bullet}\left(E_{1} \oplus K_{2}\right)=Q_{\bullet}\left(E_{1}\right) Q_{\bullet}\left(E_{2}\right)
$$

\subsection*{Beispiele}
\begin{enumerate}
  \item Der Chern-Charakter
\end{enumerate}

$$
Q(x)=e^{x}=1+x+\frac{x^{2}}{2}+\ldots
$$

Die zugehörige additive charakteristische Klasse ist der Chern-Charakter

$$
\begin{gathered}
\operatorname{ch}(E)=n+c_{1}(E)+\frac{1}{2}\left(c_{1}^{2}(E)-2 c_{2}(E)\right)+\frac{1}{6}\left(c_{1}^{3}(E)-3 c_{1}(E) c_{2}(E)+3 c_{3}(E)\right)+\ldots \\
=\operatorname{Spure}^{c(E)}
\end{gathered}
$$

Ist $P \rightarrow S^{4}$ ein $S U(2)$-Bündel, so ist

$$
\operatorname{ch}(\Omega)=2+\operatorname{Spur}\left(\frac{i \Omega}{2 \pi}\right)+\frac{1}{2} \operatorname{Spur}\left(\frac{-\Omega}{2 \pi}\right)^{2} .
$$

$\int_{S^{4}} \mathrm{ch}_{2}(\Omega)$ heißt Instanton-Zahl.\\
Der Vorteil des Chern-Charakters gegenüber den Chern-Klassen ist, dass er sich bei Addition eines trivialen Bündels ändert. Es gilt\\
(a) $\operatorname{ch}\left(f^{*} E\right)=f^{*} \operatorname{ch}(E)$\\
(b) $\operatorname{ch}(E \oplus F)=\operatorname{ch}(E)+\operatorname{ch}(F)$\\
(c) $\operatorname{ch}(E \otimes F)=\operatorname{ch}(E) \operatorname{ch}(F)$.\\
(direkter Beweis: Nakahara, Differentialgeometrie, Topologie und Phy$s i k)$.\\
2. Die Todd-Klasse ist die multiplikative Klasse, die zu $\frac{x}{1-e^{-x}}=1+\frac{x}{2}+ \frac{x^{2}}{12}-\frac{x^{4}}{720} \pm \ldots$ gebildet wird.\\
3. Das $L$ Polynom ist die multiplikative charakteristische Klasse, die aus den Pontrjagin-Klassen für reelle Vektorraumbündel zu

$$
L(x)=\frac{x}{\tanh x}
$$

gebildet wird, also

$$
L(x)=\prod\left(1+\sum_{n \geq 1}(-1)^{n-1} \frac{2^{2 n}}{(2 n)!} B_{n} x_{j}^{2 n}\right),
$$

wobei $B_{n}$ die Bernoulli-Zahlen sind, also

$$
L(E)=1+\frac{1}{3} p_{1}(E)+\frac{1}{45}\left(-p_{1}^{2}+7 p_{2}\right)+\ldots
$$

\begin{enumerate}
  \setcounter{enumi}{3}
  \item Ist $E$ ein reelles Vektorraumbündel, so ist das $\hat{A}$-Geschlecht die multiplikative charakteristische Klasse, die zu
\end{enumerate}

$$
\hat{A}(x)=\frac{x / 2}{\sinh (x / 2)}=1-\frac{x}{24}+\frac{7 x^{2}}{45 \cdot 2^{7}}-\ldots
$$

aus den Pontrjagin-Klassen gebildet wird. Für kompakte 4-dimensionale reelle Mannigfaltigkeiten heißt

$$
\int_{M} \hat{A}_{m}(T M)
$$

das $\hat{A}$-Geschlecht von $M$. Also

$$
\hat{A}(E)=1-\frac{1}{24} p_{1}+\frac{1}{5760}\left(7 p_{1}^{3}-4 p_{2}\right)+\ldots
$$

Additive und multiplikative charakteristische Klassen sind in der Indextheorie wichtig. Sie setzen den analytischen Index eines Differentialoperators in Beziehung zu topologischen Invarianten, z.B. der Satz von Chern-Gauß-Bonnet für den Index von $d+d^{*}$

$$
\chi(M)=\sum_{j=1}^{2 m}(-1)^{j} \operatorname{dim} H_{d R}^{j}(M)=\int_{M} e(M)
$$

Atiyah-Singer für den Index des Dirac-Operators $D$

$$
\operatorname{Ind}\left(D^{+}\right)=\int_{M} \operatorname{ch}(E) \cdot \hat{A}(T M)
$$

Hirzebruchscher Signatursatz

$$
\operatorname{sign}(M)=\int_{M} L(M)
$$

Riemann-Roch-Hirzebruch für die Dimension der Dolbeault-Kohomologiegruppen $h^{q}=h^{0, q}$

$$
\sum_{q}(-1)^{q} h_{q}=\int_{M} \operatorname{ch}(E) \operatorname{Tod}(T M) \text { mit } h_{q}=\operatorname{dim}_{\mathbb{C}} H_{i}(M, E)
$$

\subsection*{Bemerkung}
Die Signatur einer $4 n$-dimensionalen kompakten orientierten Mannigfaltigkeit ist die Signatur der Bilinearform

$$
\Lambda: H^{2 n} M \times H^{2 n} M \rightarrow \mathbb{R},([\omega],[\eta]) \mapsto \int_{M} \omega \wedge \eta
$$

Die Hirzebruchsche Signaturformel sagt, dass diese durch

$$
\int_{M} L\left(p_{1}(M), \ldots, p_{n}(M)\right)
$$

mit $L(x)=\prod_{i=1}^{n} \frac{x_{i}}{\tanh x_{i}}$ gegeben ist.





\end{document}